%! Public Domain Software 2022-2026
%! 2026-03-04
%! 1.0
%! Elayson Abreu
%! abntexto.classe@gmail.com

% !TeX TS-program = lualatex
\documentclass{abntexto}

\usepackage[german,english,brazil,shorthands=off]{babel}
\usepackage[cmyk]{xcolor}
\usepackage{lipsum}
\usepackage[
    style       = abnt,
    maxbibnames = 100,
    extrayear   = true,
]{biblatex} \addbibresource{abntexto.bib}
\usepackage{unicode-math}
\usepackage[colorlinks,linktoc=page]{hyperref}
\usepackage{microtype}
\usepackage[brazilian,capitalize,noabbrev]{cleveref}\AtBeginDocument{\def\ref#1{\cref{#1}}}
\usepackage{fontspec}
\usepackage{listings}
\usepackage{lua-widow-control} \clubpenalty=10000\widowpenalty=10000

\makeatletter

% VERBATIM
% ================================================

\def\adef#1{\catcode`#1=13 \begingroup \lccode`\~=`#1\lowercase{\endgroup\def~}}
\def\setverb{\def\do##1{\catcode`##1=12}\dospecials\space}
\def\verbchar#1{%
    \ifx\savedttchar\undefined\else \catcode\savedttchar=\savedttcharc \fi
    \chardef\savedttchar=`#1\relax
    \chardef\savedttcharc=\catcode`#1\relax
    \adef{#1}{\leavevmode\hbox\bgroup\setverbchar\readverb}%
    \bgroup\lccode`\~=`#1\lowercase{\egroup\def\readverb ##1~}{##1\egroup}%
    \catcode`#1=13
}
\def\setverbchar{\setverb
    \adef{ }{\ }%
    \adef{\{}{\bgroup\normalcolor\char`\{}%
    \adef{\}}{\char`\}\egroup}%
    \adef{[}{\bgroup\normalcolor\char`[}%
    \adef{]}{\char`]\egroup}%
    \ttfamily
}

\edef\bslash{\csstring\\}
\def\mytarget#1{\vbox to0pt{\kern-12pt \hypertarget{#1}{}\vss}}
\def\1{`}\def\2{``}
{\catcode`\`=13 \AtBeginDocument{%
    \verbchar{`}%
    \def\^`{\bgroup \docpoint}%
    \def\docpoint#1`{\egroup\leavevmode\edef\tmp{\csstring#1}%
        \expandafter\parsecn \string#1\relax
        \ifcsname \cn s/^\tmp\endcsname \else
            \mytarget{\cn s/^\tmp}\csgdef{\cn s/^\tmp}{}%
        \fi
        \hyperlink{\cn s/\tmp}{\color{blue}\ttfamily\if\cn c\bslash\fi\tmp}%
    }%
    \def\parsecn #1#2\relax {\edef\cn{\if\bslash#1c\else n\fi}}%
    \def\`{\bgroup \mainpoint}%
    \def\mainpoint #1`{\egroup\leavevmode\edef\tmp{\csstring#1}%
        \expandafter\parsecn \string#1\relax
        \ifcsname \cn s/\tmp\endcsname \moremainpoints \else
            \mytarget{\cn s/\tmp}\csgdef{\cn s/\tmp}{}%
        \fi
        \ifcsname \cn s/^\tmp\endcsname
            \hyperlink{\cn s/^\tmp}{\ttfamily\if\cn c\bslash\fi\tmp}\else
            {\color{red}\ttfamily\string#1}%
        \fi
    }%
    \def\moremainpoints{\errmessage{Second main documentation point \if\cn c\bslash\fi\tmp}}
}}

\let\lchevron=<
{\catcode`<=13 \AtBeginDocument{\catcode`<=13
    \def<#1>{\ifhmode\else\expandafter\hbox\fi\bgroup\normalcolor\ttfamily
             $\color{gray}\langle$\/{\rmfamily\itshape #1\/}$\color{gray}\rangle$\egroup}}%
}

\def\type#1{\NoCaseChange{\protect\typeA{#1}}}
\def\typeA#1{\bgroup\normalshape\ttfamily{\escapechar=-1\relax\expandafter}\detokenize{#1}\egroup}
% Precisamos de \m@ne aqui, \m@ne=-1.
\pdfstringdefDisableCommands{%
    \def\type#1{\bgroup\escapechar\m@ne\expandafter\egroup\detokenize{#1}}%
    \let\escapechar=\relax
    \let\m@ne=\relax
}

\definecolor{P}{RGB}{242, 121, 0}   % Primitivas.
\definecolor{C}{RGB}{23, 230, 23}   % Comentários.
\definecolor{B}{RGB}{245, 250, 247} % Fundo.

\lstset{
    language        = [latex]tex,
    basicstyle      = \abntsmall\singlesp\ttfamily,
    texcsstyle      = *\color{P}, % O listings não permite ":"
    commentstyle    = \color{C},  % no nome dessas cores. Bug.
    backgroundcolor = \color{B},
    frame           = leftline,
    rulecolor       = \color{lightgray},
    columns         = fullflexible,
    breaklines      = false,
    keepspaces      = true,
    aboveskip       = 0pt plus1.2ex minus.67ex,
    belowskip       = 0pt plus1.2ex minus.67ex,
    abovecaptionskip = 0pt,
    belowcaptionskip = 0pt,
%    tabsize         = 1, % Sem efeito, porque o documento usa Espaço na tabulação.
}

% #1 = linewidth
% #2 = border
% #3 = padding
% #4 = padding-top
% #5 = padding-right
% #6 = padding-bottom
% #7 = padding-left
% #8 = margin-left
% #9 = margin-right
\def\mylstset#1#2#3#4#5#6#7#8#9{\lstset{
    linewidth          = \dimexpr#1 -#2 -#3 -#5\relax,
    framerule          = \dimexpr#2\relax,
    framesep           = \dimexpr#3\relax,
    xleftmargin        = \dimexpr#2 +#3 +#7 +#8\relax,
    xrightmargin       = \dimexpr#9\relax,
    framextopmargin    = \dimexpr#4\relax,
    framexrightmargin  = \dimexpr#5\relax,
    framexbottommargin = \dimexpr#6\relax,
    framexleftmargin   = \dimexpr#7\relax,
}}
\mylstset{\linewidth}{1pt}{0pt}{0pt}{0pt}{0pt}{0pt}{0pt}{0pt}

\lstnewenvironment{lst}[1][\placepos]{%
    \let\medskipamount=\smallskipamount
    \place[#1]\vbox\bgroup
}{%
    \egroup\endplace
}

\def\hookXP{}
\lstnewenvironment{xp}[1][]{%
    \lstset{
        frame           = leftline,
        basicstyle      = \normalsize\ttfamily\singlesp\vskip-\baselineskip,
        texcsstyle      = \normalcolor,
        commentstyle    = \normalcolor,
        backgroundcolor = \color{white},
        gobble          = 4,
    }%
    \mylstset{\linewidth}{0pt}{0pt}{0pt}{0pt}{0pt}{\parindent}{0pt}{0pt}%
    \expandafter\lstset \expandafter{\hookXP,#1}%
    \endgraf\null
}{}

% ETC.
% ================================================

\setmainfont{XITS}[
    UprightFont    = *-Regular,
    BoldFont       = *-Bold,
    ItalicFont     = *-Italic,
    BoldItalicFont = *-BoldItalic,
    Extension      = .otf
]
\setmathfont{XITSMath-Regular.otf}
\setmonofont{InconsolataN}[
    UprightFont = *-Regular,
    BoldFont    = *-Bold,
    Extension   = .otf
]

\definelegendplace{code}{Código}{loc}
\definelegendplace{output}{Saída}{loo}
\definelegendplace{scr}{\emph{Screenshot}}{loscr}

\Crefname{code}{Código}{Códigos}
\Crefname{output}{Saída}{Saídas}
\Crefname{appendix}{Apêndice}{Apêndices}
\Crefname{codesection}{Seção}{Seções}
\Crefname{codesubsection}{Seção}{Seções}
\Crefname{scr}{\emph{Screenshot}}{\emph{Screenshots}}

\let\over=\@@over
\AtBeginDocument{\def\refname{Referências}} % Removido \^ da definição.
\def\hooksection{}
\def\tocsubsectionfont{\itshape}
\def\subsectionfont   {\itshape}
\AtBeginDocument{\pretocmd\appendix{\clearpage}{}{}}
\pretocmd\annex{\clearpage}{}{}
\def\etex{\leavevmode\hbox{$\varepsilon$-\TeX}}
\def\pdflatex{\leavevmode\hbox{PDF\LaTeX}}
\def\xelatex{\leavevmode\hbox{Xe\LaTeX}}
\def\lualatex{\leavevmode\hbox{Lua\LaTeX}}
\def\me{Elaboração própria.}
\def\bibfont{\raggedright\interlinepenalty=10000\singlesp\bibitemsep=\baselineskip}
\appto\textual{
    \def\@evenhead{\abntsmall \firstmarks0\hfil\thepage}
    \def\@oddhead {\abntsmall \botmarks1\hfil\thepage}
    \def\sectionmark   #1{\marks0{Seção \thesection: \unexpanded{#1}}}
    \def\subsectionmark#1{\marks1{Subseção \thesubsection: \unexpanded{#1}}}
    \preto\@oddhead {\lower4pt\rlap{\vrule width\hsize height.4pt}}
    \preto\@evenhead{\lower4pt\rlap{\vrule width\hsize height.4pt}}
}
\patchcmd{\eletroniclayout}{false}{true}{}{}
\let\onesidelayout=\eletroniclayout
\let\twosidelayout=\eletroniclayout
\hypersetup{
    pdfauthor   = Elayson Abreu,
    pdftitle    = Classe ABNTexto: normas ABNT em LaTeX,
    pdfsubject  = Manual da classe ABNTexto,
    pdfcreator  = LaTeX with abntexto,
    pdfkeywords = ABNTexto. LaTeX. Classe. ABNT,
}
\makeatother



\begin{document}
\leavevmode\vskip1.9in
\begingroup \centering \LARGE\sffamily\bfseries
    Classe ABNTexto\vskip-3pt \Large\itshape\mdseries
    Normas ABNT em {\upshape\LaTeX}\vskip2em \normalfont\large
    Elayson Abreu\vskip-3pt
    \href{mailto:abntexto.classe@gmail.com}%
                {abntexto.classe@gmail.com}\vskip2ex
    \today
\vskip2in\endgroup

\noindent\hfil {\Large\ttfamily 1.0}
\newpage

\nonum\notoc\section{Sumário}
\maketoc
\newpage

\section{Introdução}

O `abntexto` é uma classe do \LaTeX\ criada para facilitar a utilização das normas da Associação Brasileira de Normas Técnicas (ABNT) em trabalhos acadêmicos (monografias, dissertações, teses). Ela carrega os recursos necessários para a elaboração de trabalhos nessas normas.

Este \emph{não} é um manual sobre LaTeX (ou TeX, a ferramenta que o constitui) nem tampouco uma descrição detalhada sobre as normas ABNT. Para aprender sobre LaTeX leia \textcite{massago-2022-latex2e-via-exemplos}, em português, ou \textcite{martinsen-2022-latex2e-an-unnoficial}, disponível em inglês, espanhol e francês. Além disso, também se pode ler \textcite{olsak-2012-tex-in-nutshell} e \textcite{eijkhout-1992-tex-by-topic}, disponíveis em inglês, para aprender sobre o sistema \TeX.

Já existe uma classe para aplicação das normas ABNT cujo nome é `abntex2` \cite{araujo-2012-a-classe-abntex2}. Mas, diferente dessa, o `abntexto` não depende do `memoir`, uma classe de código-fonte e manual extensos e que implementa funções as quais, em sua maioria, não são pertinentes a um documento que use as diretrizes da ABNT.

As normas ABNT as quais o `abntexto` se propõe a seguir seguem listadas:

\begin{topics}
    \item NBR 14724:2024. Informação e documentação --- Trabalhos acadêmicos --- Apresentação \cite{abnt-2024-14724-trabalhos-academicos};
%    \item NBR 12225:2004. Informação e documentação --- Lombada --- Apresentação \cite{abnt-2004-12225-lombada};
%    \item NBR 6028:2021. Informação e documentação --- Resumo --- Apresentação \cite{abnt-2021-6028-resumo-resenha-recensao};
    \item NBR 6027:2012. Informação e documentação --- Sumário --- Apresentação \cite{abnt-2012-6027-sumario};
    \item NBR 10520:2023. Informação e documentação --- Citações em documentos --- Apresentação \cite{abnt-2023-10520-citacoes};
%    \item NBR 6023:2018. Informação e documentação --- Referências --- Elaboração \cite{abnt-2018-6023-referencias};
%    \item NBR 5892:2019. Informação e documentação --- Representação e formatos de tempo --- Datas e horas --- Apresentação \cite{abnt-2019-5892-representacao-e-formatos-de-tempo};
    \item NBR 6024:2012. Informação e documentação --- Numeração progressiva das seções de um documento --- Apresentação \cite{abnt-2012-6024-numeracao-progressiva-das};
    \item IBGE. Normas de apresentação tabular \cite{ibge-1993-normas-de-apresentacao}.
\end{topics}

\subsection{Nota de compatibilidade}

No salto da versão anterior, 3.2.1-beta, para esta, 4.0.6-beta, muito do código-fonte foi reescrito; funcionalidades foram adicionadas; comandos obsolesceram. Por consequência, documentos antigos irão certamente quebrar em \emph{layout} e aparência.

Para que os usuários tenham tempo de atualizar para esta nova versão, o arquivo `abntexto-3-2-1-beta.cls` foi deixado no repositório da classe no CTAN\footnote{\url{https://ctan.org/pkg/abntexto}.}. Todavia, para documentos antigos que já foram concluídos, não faz sentido atualizá-los para a versão atual. Nesse caso, seria pertinente copiar o arquivo `abntexto-3-2-1-beta.cls` para o mesmo diretório do documento cuja fase de escrita já foi concluída, chamá-lo por meio de `\documentclass{abntexto-3-2-1-beta}`, e, dessa forma, garantir que o resultado da compilação seja sempre o mesmo.

Para usuários que desejam atualizar seus documentos, o `abntexto` preservou a funcionalidade de comandos obsoletados até onde foi possível e muitas mensagens de alerta/erro foram criadas.

\subsection{Mudanças e novidades}

A sintaxe de `\place{<conteúdo>}` foi substituída por `\begin{place}[<opcional>]`. Além disso, após requisições, a classe ganhou suporte para objetos flutuantes. O posicionamento de `place` pode ser alterado em `<opcional>` com as palavras-chave de objetos flutuantes: `!, t, b, p, h`. Além destas, a classe criou uma nova: `here`\footnote{Semelhante a especificação `H` do pacote `float` \cite{float-package}.}. Com ela, o usuário impõe que `<conteúdo>` seja exibido no PDF na mesma posição em que foi inserido no documento.

Comandos da forma `\<ext>labelwidth` tais como `\toclabelwidth`, `\appendixlabelwidth`, entre outros, não precisam mais ser configurados, porque o `abntexto` alinha os rótulos de suas respectivas entradas automaticamente com a ajuda dos novos comandos `\eqbox` e `\eqboxsize`.

O pacote `enumitem` é carregado pela classe e o ambiente `topics` foi reescrito em função do referido pacote.

A partir de agora, use `\abntsmall` ao invés de `\small`. Por enquanto, o significado deste será igual ao daquele.

\subsection{\emph{Feedback}}

Sugestões, dúvidas, relatos de \emph{bugs} e doações são bem-vindas:\hfil\break {\ttfamily abntexto.classe@gmail.com}.

Chave PIX: {\ttfamily abntexto.classe@gmail.com}

\section{Chamando a classe}

O autor da classe recomenda o uso do \lualatex\ na compilação, embora o \pdflatex\ e o \xelatex\ também sejam suportados. Digite as linhas do \ref{cd:inicio} em conjunto com o \lualatex\ e use `abntexto`.

\legend{code}{Início}
\src\me
\label{cd:inicio}
\begin{lst}
    \documentclass{abntexto}

    \usepackage[english,brazil]{babel}
    \usepackage{fontspec}
    \usepackage{unicode-math}

    \setmainfont{XITS}[
        UprightFont    = *-Regular,
        BoldFont       = *-Bold,
        ItalicFont     = *-Italic,
        BoldItalicFont = *-BoldItalic,
        Extension      = .otf
    ]
    \setmathfont{XITSMath-Regular.otf}

    \begin{document}
        Texto
    \end{document}
\end{lst}

O arquivo `abntexto-exemplo.tex`, disponível no diretório da classe no CTAN, mostra como elaborar os elementos pré-textuais.

\section{Fonte tipográfica}

O tamanho da fonte é 12\,pt de imediato e pode ser configurado redefinindo `\normalsize` por meio de \^`\definesize`, cuja sintaxe é
\begin{xp}
    \definesize\<seletor>{<tamanho de fonte>}{<entrel. simples>}%
                         {<configs. extras>}
\end{xp}
O primeiro argumento é o nome do seletor (`\normalsize`, por exemplo). O segundo argumento é o tamanho da fonte. O terceiro é a entrelinha simples que deve ser selecionada levando em conta o valor do segundo argumento. Por exemplo, uma fonte com tamanho 10\,pt deve ter por volta de 12\,pt de entrelinha simples. O quarto argumento é para ajuste de valores dependentes do tamanho da fonte, por exemplo `\abovedisplayskip` ou `\bigskipamount`.

A classe `article`, na qual se baseia o `abntexto`, fornece os seletores habituais `\tiny`, `\footnotesize`, `\Large` etc. No entanto, esta classe se limitou a redefinir apenas dois: `\normalsize` para o tamanho padrão do texto e \^`\abntsmall` para satisfazer a exigência da \textcite{abnt-2024-14724-trabalhos-academicos} acerca do tamanho menor e uniforme. Até a versão anterior da classe, 3.2.1-beta, o comando `\small` era usado, mas foi substituído por `\abntsmall` para não haver conflito com o significado original de `\small` da classe `article`.

O espaçamento é de 1,5 por padrão e pode ser alterado com os atalhos \^`\singlesp` e \^`\onehalfsp` que alternam para espaçamento simples e de 1,5, respectivamente\footnote{O \type{\\doublesp} está obsoleto.}. A fórmula desses comandos foi atualizada, portanto documentos com versões antigas do `abntexto` terão o \emph{layout} alterado.

Nesta versão, `\spacing` está obsoleto porque esse comando ainda usa a fórmula antiga. Além disso, ele tem uma limitação: só pode ser usado após um seletor de tamanho, porque tais seletores redefinem `\baselineskip` diretamente. Por exemplo, `\abntsmall\spacing{<fator>}` é correto, porém `\spacing{<fator>}\abntsmall` é ineficaz. Essa limitação não existe ao usar os reimplementados `\singlesp` e `\onehalf` de forma que tanto `\singlesp\onehalfsp` quanto `\onehalfsp\singlesp` produzem o mesmo efeito.

\section{\emph{Layout}}

A folha é A4 por padrão e as margens esquerda e direita são de, respectivamente, 3\,cm, 2\,cm para páginas ímpares e de, respectivamente, 2\,cm, 3\,cm para páginas pares. Além disso, as margens superior e inferior têm, nessa ordem, 3\,cm, 2\,cm para todo o documento. Essas margens foram ajustadas com o pacote `geometry`, carregado na classe.

Existem dois comandos que configuram o \emph{layout} da página: \^`\pretextual` e \^`\textual`. O primeiro é iniciado automaticamente durante `\begin{document}`, invoca \^`\onesidelayout` e remove a numeração de página. O segundo, `\textual`, também é executado automaticamente pelo primeiro `\section` numerado, invoca \^`\twosidelayout` e ativa a numeração de página.

O comando `\onesidelayout` configura as páginas para impressão de somente um lado do papel (anverso) ao passo que `\twosidelayout` espelha as margens do documento, habilitando a impressão para os dois lados do papel (anverso e verso).

Os trabalhos acadêmicos podem vir a ser disponiblizados num repositório \emph{online}. Nesse caso, o usuário pode querer reconfigurar as margens de forma a se ter uma leitura mais confortável em dispositivos eletrônicos. Para esse propósito foi criado o comando \^`\eletroniclayout`, apesar dessa configuração não estar prevista na ABNT.

Há duas formas de ativá-lo. A primeira é, no preâmbulo, redefinir `\pretextual` e `\textual` substituindo as instruções `\onesidelayout` e `\twosidelayout` por `\eletroniclayout`. A segunda possibilidade é, também no preâmbulo, redefinir `\onesidelayout` e `\twosidelayout` diretamente conforme \ref{cd:eletroniclayout}. O comando `\eletroniclayout` centraliza o corpo do texto configurando 2,5\,cm nas margens esquerda e direita.

\legend{code}{\emph{Layout} para dispositivos eletrônicos}
\src\me
\label{cd:eletroniclayout}
\begin{lst}
    \documentclass{abntexto}
    \let\onesidelayout=\eletroniclayout
    \let\twosidelayout=\eletroniclayout
    \begin{document}
        texto
    \end{document}
\end{lst}

Outros comandos de estilo de página tais como `\pagestyle`\{{\ttfamily headings, empty, myheadings, plain}\} herdados da classe `article` não fazem sentido no `abntexto`. Ao invés de usá-los, é aconselhável redefinir `\textual`, `\pretextual` ou ambos.

\section{Sumário}

A classe `article` criou o comando `\tableofcontents` para impressão do Sumário, cuja definição será mantida por compatibilidade. No `abntexto`, porém, é recomendado usar \^`\maketoc`, cuja definição é ligeiramente diferente.

Os comandos que configuram a fonte tipográfica das entradas do Sumário são \^`\tocsectionfont`, \^`\tocsubsectionfont` etc. (ver \ref{cd:contents}). O último token nessas definições pode conter parâmetro, por exemplo:
\begin{xp}
    \def\tocsectionfont{\bfseries\MakeUppercase}
    \def\tocparagraphfont{\fbox}
\end{xp}

\legend{code}{Aparência do Sumário}
\src\me
\label{cd:contents}

\begin{lst}
    \documentclass{abntexto}
    \usepackage{xcolor}
    \def \tocsectionfont    {\color{red}\itshape\MakeUppercase}
    \def \tocsubsectionfont {\color{blue}\bfseries}

    \begin{document}
        \nonum\notoc\section{Sumário}
        \maketoc
        \newpage

        \section{Título}
        \subsection{Título}
    \end{document}
\end{lst}

\section{Secionamento}\label{sec:secionamento}

A personalização das macros de seção (\^`\section`, \^`\subsection`, \^`\subsubsection`, \^`\paragraph` e \^`\subparagraph`) é análoga a personalização do Sumário e deve ser feita redefinindo \^`\sectionfont`, \^`\subsectionfont`, \^`\subsubsectionfont`, \^`\paragraphfont` e \^`\subparagraphfont` como exemplificado em \ref{cd:sec} abaixo.

\src\me
\legend{code}{Aparência das seções}
\label{cd:sec}
\begin{lst}
    \documentclass{abntexto}
    \usepackage{xcolor}
    \def \sectionfont       {\color{blue}\scshape}
    \def \subsectionfont    {\color{red}\itshape}
    \def \subsubsectionfont {\color{yellow}\sffamily}

    \begin{document}
        \nonum\notoc\section{Sumário}
        \maketoc

        \section{Texto}
        \subsection{Texto}
        \subsubsection{Texto}
    \end{document}
\end{lst}

Também há como alterar o espaço vertical acima e abaixo das seções redefinindo \^`\abovesection`, \^`\belowsection`. Por exemplo, `\def\abovesubsection{2cm}` insere um espaço de 2\,cm acima de `\subsection`. Se o leitor deseja zerar essas instruções, deve escrever `\def\abovesubsection{0pt}` ao invés de `\def\abovesubsection{}`.

\subsection{O uso de \type{\\chapter} e \type{\\part}}

De acordo com a \textcite{abnt-2024-14724-trabalhos-academicos}: \enquote{[\dots] O trabalho acadêmico não pode ser dividido em capítulos; deve ser organizado em seções [\dots]}.

O usuário pode querer utilizar capítulos no trabalho, apesar de a unidade secional `\section` ser suficiente para um trabalho acadêmico. Para esse propósito, o `abntexto` disponibiliza a instrução \^`\chapter` para a utilização de capítulos no documento, porém não em conformidade com a \textcite{abnt-2024-14724-trabalhos-academicos}, mas no estilo da classe `book`. Para ativá-los é preciso chamar a instrução \^`\usechapters` no preâmbulo.

Também é possível usar \^`\part` em conjunto com o comando \^`\useparts`. Esse útlimo contém `\usechapters` em sua definição, portanto, ao chamar `\useparts`, não se deve chamar `\usechapters` para que este não sobrescreva as definições daquele.

Os níveis secionais `\section` e `\chapter` podem ser prefixados com \^`\nonum` que remove a numeração do título: `\nonum\section`. Os títulos prefixados com `\nonum` serão centralizados \cite{abnt-2024-14724-trabalhos-academicos}. Veja um exemplo de uso em \ref{cd:nonum}.

\legend{code}{Exemplo de título sem numeração e sem ingresso ao Sumário}
\src\me
\label{cd:nonum}

\begin{lst}
    \nonum\notoc\section{Resumo}
    Texto

    \nonum\notoc\section{Sumário}
    \maketoc
\end{lst}

As unidades secionais supracitadas também podem ser prefixadas com \^`\notoc` que instrui o título a não ingressar no Sumário: `\notoc\section`. Além disso, é possível usar ambos os prefixos ao mesmo tempo, conforme mostrado em \ref{cd:nonum}.

O primeiro parágrafo após uma seção é indentado por padrão, dispensando, assim, o uso do pacote `indentfirst`.

Infelizmente, comandos como `\printbibliography`, `\printglossary` e `\printindex` inserem o título internamente. Para que esse título seja exibido corretamente no `abntexto`, é necessário usar o ambiente \^`corrprint`, como demonstrado em \ref{cd:corrprint}.

\legend{code}{Ambiente \type{corrprint}}
\src\me
\label{cd:corrprint}

\begin{lst}
    \begin{corrprint}
        \printbibliography % ou \printglossary, \printindex e semelhantes
    \end{corrprint}
\end{lst}

Existe a possibilidade de o usuário desejar criar um Índice usando o ambiente `theindex` --- embora alternativas automatizadas sejam preferíveis \cite{martinsen-2022-latex2e-an-unnoficial}. Nesse caso, o ambiente `corrprint` funcionará como esperado conforme explicitado em \ref{cd:corrprint-theindex}. O ambiente `corrprint` também pode ser utilizado com o ambiente `thebibliography`, assim como no exemplo com `theindex`.

\legend{code}{Ambientes \type{corrprint} e \type{theindex}}
\src\me
\label{cd:corrprint-theindex}

\begin{lst}
    \begin{corrprint}
        \begin{theindex}
            \item carro, 1
                \subitem árvore, 2
            \indexspace
            \item mar, 3
                \subitem continente, 4
        \end{theindex}
    \end{corrprint}
\end{lst}

\section{Alíneas}

Use o ambiente \^`topics` para criar alíneas. Ele dispõe de dois níveis sendo que o segundo usa um travessão como marcador conforme o exemplo em \ref{cd:topics}. Naturalmente, é permitido usar `\label` após um item.

\legend{code}{Alíneas}
\src\me
\label{cd:topics}

\begin{lst}
    \begin{topics}
        \item \label{al:x} 1
            \begin{topics}
                \item \label{subal:x} x
                \item y
                \item z
            \end{topics}
        \item \label{al:y} 2
        \item 3
    \end{topics}
\end{lst}

\section{Áreas de legenda\label{sec:area-de-legenda}}

Use o ambiente \^`place` para inserir áreas de legenda. Ele substitui integralmente os ambientes `figure` e `table`, bem como o comando `\caption`, não suportados pela classe. Além disso, tem-se \^`\legend`\nobreak`{<tipo>}{<legenda>}` para titular; \^`\src`\nobreak`{<fonte>}` (\emph{source}) para inserir a fonte como mostrado em \ref{cd:place}. Porém, existem condições: `place` deve constar em último lugar e `\label` deve suceder `\legend`.

\legend{code}{Áreas de legenda}
\src\me
\label{cd:place}

\begin{lst}
    \begin{document}
        \legend{figure}{Um título}
        \src{\texttt{www.site.com.br}. Acesso em 01 jan. 2020.}
        \label{fig:teste}
        \begin{place}
            \includegraphics[width=0.4\linewidth]{example-image}
        \end{place}

        Ver \ref{fig:teste}.
    \end{document}
\end{lst}

Além disso, existe a possibilidade de alterar o espaço vertical acima e abaixo de uma área de legenda com `\def`\^`\aboveplace``{<valor>}` e `\def`\^`\belowplace``{<valor>}`. Todavia, tais comandos não têm efeito em objetos flutuantes. Para modificar o espaçamento em `floats` o leitor pode fornecer valores diferentes para `\floatsep`, `\intextsep` e `\textfloatsep`, definidos no LaTeX.

A partir desta versão, o `abntexto` passou a suportar objetos flutuantes. Por exemplo, em
\begin{xp}
    \begin{place}[<especificação>]
\end{xp}
A `<especificação>` é composta pelas palavras-chave habituais: `!`, `t`, `b`, `p`, e `h`. Alternativamente, é possível usar a palavra-chave `here`, definida nesta classe. Com ela, o conteúdo de `place` é mostrado na página exatamente onde foi escrito no documento.

A especificação padrão é `here` caso nenhuma seja explicitamente assumida, mas o usuário pode mudar o padrão digitando, por exemplo, `\def`\nobreak\^`\placepos``{tbp}` no preâmbulo do documento.

\subsection{Criando áreas de legenda}

Até a versão anterior, 3.2.1-beta, o `abntexto` fornecia \^`\definelegendplace` para criação de novos espaços de legenda sob a sintaxe
\begin{xp}
    \definelegendplace{<tipo>}{<nome>}{<ref rótulo>}{<extensão>}
\end{xp}
Nesta versão, no entanto, o terceiro argumento foi excluído em favor do suporte mais extensivo do pacote `cleveref.sty` (não carregado pela classe). Veja o \ref{app:cleveref} para mais detalhes. A partir de agora, a nova sintaxe é
\begin{xp}
    \definelegendplace{<tipo>}{<nome>}{<extensão>}
\end{xp}
Esta classe continuará fornecendo suporte sintático, porém sem funcionalidade, para a versão anterior do referido comando, mas vai exibir uma mensagem de alerta pedindo para que o terceiro argumento seja removido.

Acerca do novo `\definelegendplace`: a instrução
\begin{xp}
    \definelegendplace{diagram}{Diagrama}{lod},
\end{xp}
por exemplo, criará o comando `\makelod` para escrever no documento as entradas do arquivo auxiliar `.lod`. Além disso, o tipo de legenda `diagram` estará disponível para ser usado em `\legend`: `\legend{diagram}{<título>}`.

\subsection{Figuras lado a lado}

O `abntexto` tem suporte nativo para inserção de figuras lado a lado com o ambiente \^`multiplace`. Elas se comportam como caracteres então podemos usar `\hfil` para centralizá-las. Veja o exemplo no \ref{cd:multiplace} e o resultado em \ref{out:result}.

\legend{code}{Figuras lado a lado}
\src\me
\label{cd:multiplace}
\begin{lst}
    \legend{figure}{Duas figuras}
    \src{Elaboração do autor.}

    \begin{multiplace}
        \sublegend{Um título} \label{fig:teste1}
        \subsrc{\texttt{www.site.com.br}. Acesso em 01 jan. 2020.}
        \hfil
        \begin{subplace}[.25\linewidth]
            \includegraphics[width=\linewidth]{example-image-a}
        \end{subplace}
        %
        \sublegend{Outro título} \label{fig:teste2}
        \subsrc{\texttt{www.site.com.br}. Acesso em 01 jan. 2020.}
        \hfil
        \begin{subplace}[.25\linewidth]
            \includegraphics[width=\linewidth]{example-image-b}
        \end{subplace}
    \end{multiplace}

    \parindent=1.5cm Ver \ref{fig:teste1} e \ref{fig:teste2}.
\end{lst}

\legend{output}{Resultado do \ref{cd:multiplace}}
\label{out:result}
\src\me
\begin{place}\vbox\bgroup
    \legend{figure}{Duas figuras}
    \src{Elaboração do autor.}
    %
    \begin{multiplace}
        \sublegend{Um título} \label{fig:teste1}
        \subsrc{\texttt{www.site.com.br}. Acesso em 01 jan. 2020.}
        \hfil
        \begin{subplace}[.25\linewidth]
            \includegraphics[width=\linewidth]{example-image-a}
        \end{subplace}
        %
        \sublegend{Outro título} \label{fig:teste2}
        \subsrc{\texttt{www.site.com.br}. Acesso em 01 jan. 2020.}
        \hfil
        \begin{subplace}[.25\linewidth]
            \includegraphics[width=\linewidth]{example-image-b}
        \end{subplace}
    \end{multiplace}

    \parindent=1.5cm Ver \ref{fig:teste1} e \ref{fig:teste2}.
\egroup\end{place}

Há uma legenda principal impressa com `\legend` e duas sublegendas para cada objeto inseridas com \^`\sublegend`\nobreak`{<título>}`, bem como um par de ambientes \^`subplace`\footnotemark. Além disso, a ordem das instruções é similar àquela aplicada no ambiente `place`: `subplace` deve ser posto após ambos `\sublegend` e \^`\subsrc``{<fonte>}`, e `\label` deve suceder `\legend`.

\footnotetext{A sintaxe \type{\\subplace{<conteúdo>}} está obsoleta. Use `subplace` na forma de ambiente: \type{\\begin{subplace}}.}

Linhas em branco produzem quebra de parágrafo em `multiplace` e, provavelmente, não são desejadas.

Normalmente, espaços em branco depois de `\end{<env>}` são significativos, mas, no caso do ambiente `subplace`, isso não é verdade, porque existe um `\ignorespacesafterend` em sua definição, um comando definido no LaTeX.

Existe, ainda, a possibilidade de um arranjo diferente no `multiplace` usando o parâmetro opcional de `\sublegend`, como se pode ver no \ref{cd:multialt} em conjunto com o resultado em \ref{out:alt}.

\legend{code}{Uso alternativo de \type{multiplace}}
\src{Elaborado pelo autor.}
\label{cd:multialt}
\begin{lst}
    \begin{multiplace}
        \sublegend[table]{Tabela de exemplo} \label{tab1}
        \subsrc{\me}
        \hfil
        \begin{subplace}[.25\linewidth]
            \includegraphics[width=\linewidth]{example-image-a}
        \end{subplace}
        %
        \sublegend[figure]{Título de figura} \label{fig1}
        \subsrc{\me}
        \hfil
        \begin{subplace}[.25\linewidth]
            \includegraphics[width=\linewidth]{example-image-b}
        \end{subplace}
        %
        \sublegend[figure]{Outra figura} \label{fig2}
        \subsrc{\me}
        \vadjust{\kern10pt}\\ \null \hfil
        \begin{subplace}[.25\linewidth]
            \includegraphics[width=\linewidth]{example-image-c}
        \end{subplace}
    \end{multiplace}
\end{lst}

\src{\me}
\begin{place}\vbox\bgroup
\legend{output}{Resultado do \ref{cd:multialt}} \label{out:alt}
\begin{multiplace}
    \sublegend[table]{Tabela de exemplo} \label{tab1}
    \subsrc{\me}
    \hfil
    \begin{subplace}[.25\linewidth]
        \includegraphics[width=\linewidth]{example-image-a}
    \end{subplace}
    \sublegend[figure]{Título de figura} \label{fig1}
    \subsrc{\me}
    \hfil
    \begin{subplace}[.25\linewidth]
        \includegraphics[width=\linewidth]{example-image-b}
    \end{subplace}
    \sublegend[figure]{Outra figura} \label{fig2}
    \subsrc{\me}
    \vadjust{\kern10pt}\\ \null \hfil
    \begin{subplace}[.25\linewidth]
        \includegraphics[width=\linewidth]{example-image-c}
    \end{subplace}
\end{multiplace}
\egroup\end{place}

A \ref{out:alt} possui a \ref{tab1}, \ref{fig1} e \ref{fig2}.

\subsection{Tabelas que se partem entre páginas}

No momento não existe suporte para legendar tabelas e outras estruturas que se partem entre páginas. Se o leitor tiver de usar `longtable.sty` \cite{longtable-package} ou `tabularray.sty` \cite{tabularray-package} (este último é mais flexível), será preciso legendar estruturas tabulares com as ferramentas fornecidas por esses pacotes.

\section{Formatação de \type{\\cite} e Referências}

O `abntexto` não possui macros para formatação de `\cite` e Referências. Utilize `\usepackage[style=abnt]{biblatex}`.

\section{Citações}

Visando seguir a norma \textcite{abnt-2023-10520-citacoes}, são fornecidos `\enquote` do pacote `csquotes`, carregado pela classe, e \^`\Enquote`. Para inserir aspas duplas em citações diretas, de até três linhas, use `\enquote{<texto>}\cite{<chave>}`. Porém, se tratando de citações diretas, com mais de três linhas, de parágrafo único, alterne para `\Enquote{<texto> \cite{<chave>}.}`.

\section{Apêndice e Anexo}

Com os comandos \^`\appendix` e \^`\annex` você pode inserir apêndices e anexos junto dos indicativos desejados: \enquote{APÊNDICE 1 --- \dots} ou \enquote{ANEXO 1 --- \dots} (ver \ref{cd:appendix}).

\legend{code}{Anexos e apêndices}
\src{\me}
\label{cd:appendix}

\begin{lst}
    \begin{document}
        \appendix{Lorem}
        \appendix{Lipsum}

        \annex{Lorem}
        \annex{Lipsum}
    \end{document}
\end{lst}

Até a versão anterior desta classe, era necessário ajustar a largura dos rótulos no Sumário manualmente redefinindo os comandos \^`\appendixlabelwidth` e \^`\annexlabelwidth`. Todavia, nesta versão, 4.0.6-beta, eles se ajustam automaticamente com o auxílio das novas macroinstruções \^`\eqbox` e \^`\eqboxsize`.

\section{Macros para identificação do documento}

\subsection{Ficha catalográfica}

Para produzir fichas catalográficas foi criado o ambiente \^`indexcard`. Ele sempre esteve disponível na forma `\indexcard{<conteúdo>}`, agora obsoleta, mas nunca foi devidamente documentado neste manual. Eis uma amostra de uso:
\begin{xp}
    \documentclass{abntexto}
    \usepackage{lipsum}
    \usepackage{xcolor}
    \begin{document}
        \def\icardWidth{12.5cm}
        \def\icardHeight{7.5cm}
        \def\icardPadding{1cm}
        \def\icardBorder{2pt}
        \appto\icardFont{\color{magenta}}
        \icardPrintGrid

        \noindent\hfil
        \begin{indexcard}
            \lipsum[1][1-4]
        \end{indexcard}
    \end{document}
\end{xp}
Que produz:

\bgroup
\def\icardWidth{12.5cm}
\def\icardHeight{7.5cm}
\def\icardPadding{1cm}
\def\icardBorder{2pt}
\appto\icardFont{\color{magenta}}
\icardPrintGrid
\noindent\hfil
\begin{indexcard}
    \lipsum[1][1-4]
\end{indexcard}
\egroup

Perceba que a instrução \^`\icardPrintGrid` desenha uma grade de quatro linhas para delimitar os contornos do bloco de texto, no intuito de verificar se ele não ultrapassou as margens internas. A definição original de `\icardFont` é
\begin{xp}
    \def\icardFont{\raggedright\parindent=2em\ttfamily}
\end{xp}

\subsection{Utilitário \type{\\Enter}}

O `abntexto` criou o comando \^`\Enter``[<número>]` para produzir um espaço vertical equivalente à tecla `Enter` do teclado. Ele é usado com frequência em `abntexto-exemplo.tex`. O argumento opcional `<número>` é o número de `Enters` que será impresso, por exemplo, `\Enter[2]` imprime um espaço vertical equivalente a dois `Enters` consecutivos. Como `\Enter` é, basicamente, `\vskip\baselineskip`, precedê-lo com `\par` não é necessário, porque `\vskip` possui, a grosso modo, um `\par` em seu funcionamento.

\subsection{Desativando a hifenização}

O comando \^`\nohyph` impede a hifenização dos parágrafos integralmente, sendo útil em títulos e parágrafos curtos. Seu uso pode ser feito assim:
\begin{xp}
    {\nohyph Parágrafo curto ou título\par} % Esse \par é necessário para
                                            % que a formatação do parágrafo
                                            % entre em ação.
\end{xp}

\subsection{Assinatura da banca}

Use \^`\judgeline``{<arg>}` para assinaturas na Folha de Aprovação. Exemplo:
\begin{xp}
    \judgeline{Prof. Dr. Nome Sobrenome \\ Instituição}\Enter
    \judgeline{Prof. Dr. Nome Sobrenome \\ Instituição}\Enter
    \judgeline{Prof. Dr. Nome Sobrenome \\ Instituição}
\end{xp}
Resultado:
\judgeline{Prof. Dr. Nome Sobrenome \\ Instituição}\Enter
\judgeline{Prof. Dr. Nome Sobrenome \\ Instituição}\Enter
\judgeline{Prof. Dr. Nome Sobrenome \\ Instituição}

Veja outro exemplo de `\judgeline` em `abntexto-exemplo.tex`.

\newpage
\begin{corrprint}
    \printbibliography
\end{corrprint}

\newrefsection

\appendix{Cleveref\label{app:cleveref}}

Até a versão passada, 3.2.1-beta, o `abntexto` tinha o poder de inferir o tipo de `\ref` e inserir o substantivo correspondente com auxílio do terceiro parâmetro de `\definelegendplace`. No entanto, esse parâmetro foi removido deixando o referido comando com três argumentos apenas. Como resgatar essa funcionalidade de inferência em `\ref`? Usando o pacote `cleveref.sty` e sua macro `\cref` \cite{cleveref-package}. Veja o exemplo em \ref{cd:cleveref}, cujo resultado é \enquote{Ler Seção 1}.

\legend{code}{Exemplo de uso do \type{cleveref.sty}}
\src\me
\label{cd:cleveref}
\begin{lst}
    \documentclass{abntexto}
    \usepackage[brazil]{babel}
    \usepackage[brazilian,capitalize,noabbrev]{cleveref}
    \begin{document}
        \section{Um título de exemplo\label{sec:ex}}
        Ler \cref{sec:ex}.
    \end{document}
\end{lst}

Percebe-se que a chave `brazil` não é reconhecida pelo `cleveref` compelindo-nos a usar `brazilian`. Além disso, ele deve ser carregado depois do `hyperref.sty`. Veja outras funcionalidades no manual do pacote.

\appendix{Lista de Abreviaturas}

Elemento opcional. O `glossaries-extra.sty` \cite{glossaries-extra-package} consegue criar uma Lista de Abreviaturas, conforme \ref{cd:abbrev}.

\legend{code}{Abreviaturas}
\src\me
\label{cd:abbrev}

\begin{lst}
    \documentclass{abntexto}

    \usepackage[brazil]{babel}
    \usepackage{hyperref}
    \usepackage[automake,shortcuts=all,abbreviations]{glossaries-extra}\makeglossaries
    \usepackage[brazilian,capitalize,noabbrev]{cleveref}

    \newabbreviation[category=abbreviations]{fil}{Fil.}{Filosofia}
    \newabbreviation[category=abbreviations]{el sf}{el.\,sf.}{elemento substantivo feminino}
    \newabbreviation[category=abbreviations]{eng el}{Eng.\,Elétr.}{Engenharia Elétrica}
    \newabbreviation[category=abbreviations]{med legal}{Med.\,Leg.}{Medicina Legal}

    \begin{document}
        Abreviaturas: \ab{fil},
                      \ab{el sf},
                      \ab{eng el},
                      \ab{med legal}.

        \begin{corrprint}\notoc % "Lista de Abreviaturas" não deve constar no Sumário.
            \printglossary[type=abbreviations]
        \end{corrprint}
    \end{document}
\end{lst}

Note que o `cleveref.sty` deve ser carregado depois de `glossaries-extra.sty`, e o `hyperref.sty`, antes.

\appendix{Lista de Siglas}

Elemento opcional. O `glossaries-extra.sty` \cite{glossaries-extra-package} também consegue criar uma Lista de Siglas, conforme \ref{cd:acro}.

\legend{code}{Siglas}
\src\me
\label{cd:acro}

\begin{lst}
    \documentclass{abntexto}

    \usepackage[brazil]{babel}
    \usepackage{hyperref}
    \usepackage[automake,shortcuts=all,acronym]{glossaries-extra}\makeglossaries
    \usepackage[brazilian,capitalize,noabbrev]{cleveref}

    \setabbreviationstyle[acronym]{long-short}

    \newacronym{ibge}{IBGE}{Instituto Brasileiro de Geografia e Estatística}
    \newacronym{html}{HTML}{HyperText Markup Language}
    \newacronym{ce}{CE}{Ceará}
    \newacronym{edo}{EDO}{Equação Diferencial Ordinária}

    \begin{document}
        1ª chamada: \ac{ibge}, \ac{html}, \ac{ce}, \ac{edo}.

        2ª chamada: \ac{ibge}, \ac{html}, \ac{ce}, \ac{edo}.

        \begin{corrprint}\notoc % "Lista de Siglas" não deve constar no Sumário.
            \printglossary[type=acronym]
        \end{corrprint}
    \end{document}
\end{lst}

Note que o `cleveref.sty` deve ser carregado depois de `glossaries-extra.sty`, e o `hyperref.sty`, antes.

\appendix{Lista de Símbolos}

Elemento opcional. O `glossaries-extra.sty` \cite{glossaries-extra-package} também consegue criar uma Lista de Símbolos, conforme \ref{cd:simbolos}.

\legend{code}{Símbolos}
\src\me
\label{cd:simbolos}

\begin{lst}
    \documentclass{abntexto}

    \usepackage[brazil]{babel}
    \usepackage{hyperref}
    \usepackage[automake,shortcuts=all,symbols]{glossaries-extra}\makeglossaries
    \usepackage[brazilian,capitalize,noabbrev]{cleveref}

    \def\newsymbol#1#2#3#4{\newglossaryentry{#1}{
        name={#2},sort={#3},description={#4},type=symbols,category=symbol}}

    \newsymbol{alpha}{\ensuremath{\alpha}}{a}{Letra grega alfa}
    \newsymbol{o(n)}{\ensuremath{O(n)}}{On}{Ordem de um algoritmo}
    \newsymbol{i}{\ensuremath{i}}{i}{Unidade imaginária}
    \newsymbol{A}{\ensuremath{\mathbf{A}}}{A}{Matriz $\mathbf{A}$}
    \newsymbol{N}{N}{N}{Newton}

    \def\symb#1{\gls{#1}}

    \begin{document}
        Símbolos: \symb{alpha}, \symb{o(n)}, \symb{i}, \symb{A}, \symb{N}.

        \begin{corrprint}\notoc % "Lista de Símbolos" não deve constar no Sumário.
            \printglossary[type=symbols]
        \end{corrprint}
    \end{document}
\end{lst}

Note que o `cleveref.sty` deve ser carregado depois de `glossaries-extra.sty`, e o `hyperref.sty`, antes.

\appendix{Glossário}

Elemento opcional. Você pode criar glossários com `glossaries-extra` \cite{glossaries-extra-package}, como em \ref{cd:gloss}.

\legend{code}{Glossário}
\src\me
\label{cd:gloss}

\begin{lst}
    \documentclass{abntexto}

    \usepackage[brazil]{babel}
    \usepackage{hyperref}
    \usepackage[automake,shortcuts=all]{glossaries-extra}\makeglossaries
    \usepackage[brazilian,capitalize,noabbrev]{cleveref}

    \newglossaryentry{metalismo}{
        name        = metalismo,
        description = {Tipo de sistema monetário cujo valor da moeda,
                       geralmente de uma nação ou de um país, é determinado
                       pela quantidade constante de certo metal (ouro ou
                       prata), desta forma se estabelece um valor fixo de
                       troca entre esses metais e o próprio dinheiro.}
    }
    \newglossaryentry{backend}{
        name        = \emph{backend},
        sort        = backend,
        description = {A parte de um site ou aplicação que não é visível
                       para o usuário e que gerencia a lógica de negócios
                       e a comunicação com o banco de dados.}
    }
    \newglossaryentry{cache}{
        name        = cache,
        description = {Um armazenamento temporário de dados para acesso
                       rápido. Ajuda a reduzir o tempo de carregamento.}
    }
    \newglossaryentry{autovetor}{
        name        = autovetor,
        description = {Um vetor não nulo que, quando multiplicado por uma
                       matriz, resulta em um vetor que é um múltiplo escalar do
                       vetor original.}
    }

    \begin{document}
        Glossário: \gls{metalismo}, \gls{backend}, \gls{cache}, \gls{autovetor}.

        \begin{corrprint}
            \printglossary[type=main]
        \end{corrprint}
    \end{document}
\end{lst}

Note que o `cleveref.sty` deve ser carregado depois de `glossaries-extra.sty`, e o `hyperref.sty`, antes.

\appendix{Ambiente \type{place} e o pacote \type{listings}}

O `listings` é um impressor de código-fonte para LaTeX \cite{listings-package}. Você pode usar `\legend` em conjunto com `listings` como mostrado em \ref{cd:listings}. Há algumas observações quanto a essa solução: primeiro, o `listings` não permite que `place` esteja na forma convencional, mas na forma `\place`-`\endplace`. Segundo, a região em \emph{source} terá largura fixa de `\linewidth`, que é a largura padrão de uma `\vbox`.

\legend{code}{\type{\\legend} em conjunto com \type{listings}}
\src\me
\label{cd:listings}

\lstnewenvironment{lstn}{%
    \let\medskipamount=\smallskipamount
    \place\vbox\bgroup
}{%
    \egroup\endplace
}

\begin{lstn}
    % No preâmbulo:
    \usepackage{listings}

    \lstnewenvironment{lst}[1][\placepos]{%
        \let\medskipamount=\smallskipamount
        \place[#1]\vbox\bgroup
    }{%
        \egroup\endplace
    }
    \definelegendplace{code}{Código}{loc}

    % No documento:
    \nonum\notoc\section{Lista de códigos}
    \makeloc

    \legend{code}{Título}
    \src{Fonte}
    \label{Chave}

    \begin{lst}
        int x = 1;
    \end{lst}
\end{lstn}

Na solução acima o ambiente `lst` é inquebrável, portanto quebras de páginas não são possíveis. Se você precisa de um ambiente que se parta entre páginas, veja a solução abaixo (\ref{cd:longlst}):

\legend{code}{O ambiente \type{longlst} quebrável como alternativa ao \type{lst}}
\src\me
\label{cd:longlst}

\begin{lst}
    \lstnewenvironment{longlst}[1][\placepos]{%
        \begingroup
        \processplacearg{#1}%
        \parskip=0pt \parindent=0pt
        \initplace
        \printlegendbox
        \endgraf\nointerlineskip \nobreak\vskip\medskipamount
        \resetORIlabel \ignorespaces
    }{%
        \unskip
        \savedplacewidth=\hsize
        \printsrcbox
        \finishplace
        \resetplace \gresetORIlabel
        \endgroup
    }
\end{lst}

\appendix{Ambiente \type{place} e o pacote \type{algorithm2e}}

O pacote `algorithm2e` foi criado para geração de algoritmos em LaTeX \cite{algorithm2e-package}. É possível usar `place` em conjunto com `algorithm2e` (\ref{cd:algorithm2e}):

\legend{code}{\type{algorithm2e} em \type{place}}
\src\me
\label{cd:algorithm2e}

\begin{lst}
    % No preâmbulo
    \usepackage{algorithm2e}

    \definelegendplace{alg}{Algoritmo}{loa}

    % No documento
    \nonum\notoc\section{Lista de algoritmos}
    \makeloa

    \legend{alg}{How to write algorithms}
    \src{Elaboração própria}
    \label{alg:chave}

    \begin{place}
    \begin{algorithm}[H]
    \SetAlgoLined
    \KwData{this text}
    \KwResult{how to write algorithm with \LaTeX2e }
    initialization\;
    \While{not at end of this document}{
        read current\;
        \eIf{understand}{
            go to next section\;
            current section becomes this one\;
        }{
            go back to the beginning of current section\;
        }
    }
    \end{algorithm}
    \end{place}
\end{lst}

A especificação `H` é obrigatória.

% MACROS PARA O CÓDIGO-FONTE
% ================================================

\pdfstringdefDisableCommands{\def\marks#1#2{}}
\appendix{\marks0{}\marks1{}Código-fonte\label{ap:codigo-fonte}}

\makeatletter
\newcounter{codesection}
\newcounter{codesubsection}[codesection]
\let\c@section=\c@codesection
\let\c@subsection=\c@codesubsection
\def\thecodesection{\arabic{codesection}}
\def\thecodesubsection{\thecodesection.\arabic{codesubsection}}
\edef\toclevel@codesection{\mainseclevel}
\edef\toclevel@codesubsection{\the\numexpr\mainseclevel-1\relax}

\appto\hooktocsection{\def\extleaders{}%
    \edef\savedskips{\rightskip=\the\rightskip \parfillskip=\the\parfillskip\relax}}
\appto\hooktocsubsection{\def\extleaders{}%
    \edef\savedskips{\rightskip=\the\rightskip \parfillskip=\the\parfillskip\relax}}
\appto\hookextline{\savedskips}

\def\toclabelbox{\lowercase{\eqbox{codetoc}}}
\def\tocsectionfont{\itshape\MakeUppercase}
\def\tocsubsectionfont{}
\def\sectionfont{\tocsectionfont}
\def\subsectionfont{\itshape}

% O \addcontentsline{#1}{#2}{#3} está programado para
% enviar marcadores apenas se #1 = toc. Podemos
% mudar isso com \hypersetup{bookmarkstype=<outro>}.
\hypersetup{bookmarkstype=codetoc}
\patchcmd{\targetsection}{\refstepcounter{section}}{\refstepcounter{codesection}}{}{}
\patchcmd{\targetsection}{\addcontentsline{toc}}{\addcontentsline{codetoc}}{}{}
\patchcmd{\targetsubsection}{\refstepcounter{subsection}}{\refstepcounter{codesubsection}}{}{}
\patchcmd{\targetsubsection}{\addcontentsline{toc}}{\addcontentsline{codetoc}}{}{}

\def\hooksection{}
\def\codesection{\section}
\def\codesubsection{\subsection}

\colorlet{C}{.}
\definecolor{tmp}{RGB}{60, 90, 110}  \colorlet{P}{tmp}
                                     \colorlet{B}{blue!7}

\expanded{\noexpand\lstdefinelanguage{mylang}{
    language    = [latex]tex,
    texcs       = {\directlua{for k,v in pairs(tex.primitives())
                   do tex.print(k, v.. ',')end}},
    deletetexcs = {end}
}}

\lstset{
    language    = mylang,
    frame       = single,
    numbers     = left,
    numbersep   = \dimexpr-5pt+\addnumbersep\relax,
    numberstyle = \fontsize{8bp}{9bp}\selectfont,
    literate    = {\^^M\%.}{}1,
}
\def\addnumbersep{9pt}
\mylstset{\linewidth}{0pt}{0pt}{3pt}{0pt}{3pt}{13pt+\addnumbersep}{0pt}{0pt}

\newcount\codecount \codecount=1
\newbox\mybox
\setbox\mybox=\hbox to\hsize{\hfil\abntsmall\ttfamily\bfseries\color{gray}abntexto.cls}
\lstnewenvironment{code}[1][]{%
    \lstset{
        title       = \copy\mybox\nobreak,
        firstnumber = \codecount,
        #1
    }%
}{%
    \global\codecount=\value{lstnumber}\relax
}
\appto\hookXP{
    frame     = none,
    numbers   = none,
    numbersep = 0pt,
    gobble    = 4,
}

\parskip=0pt plus.2ex minus.1ex
\@starttoc{codetoc}
{\def\excl{!}\catcode\1\!=14 \catcode\1\%=9 %! Public Domain Software 2022-2026
%! 2026-03-04
%! 1.0
%! Elayson Abreu
%! abntexto.classe@gmail.com

% !TeX TS-program = lualatex
\documentclass{abntexto}

\usepackage[german,english,brazil,shorthands=off]{babel}
\usepackage[cmyk]{xcolor}
\usepackage{lipsum}
\usepackage[
    style       = abnt,
    maxbibnames = 100,
    extrayear   = true,
]{biblatex} \addbibresource{abntexto.bib}
\usepackage{unicode-math}
\usepackage[colorlinks,linktoc=page]{hyperref}
\usepackage{microtype}
\usepackage[brazilian,capitalize,noabbrev]{cleveref}\AtBeginDocument{\def\ref#1{\cref{#1}}}
\usepackage{fontspec}
\usepackage{listings}
\usepackage{lua-widow-control} \clubpenalty=10000\widowpenalty=10000

\makeatletter

% VERBATIM
% ================================================

\def\adef#1{\catcode`#1=13 \begingroup \lccode`\~=`#1\lowercase{\endgroup\def~}}
\def\setverb{\def\do##1{\catcode`##1=12}\dospecials\space}
\def\verbchar#1{%
    \ifx\savedttchar\undefined\else \catcode\savedttchar=\savedttcharc \fi
    \chardef\savedttchar=`#1\relax
    \chardef\savedttcharc=\catcode`#1\relax
    \adef{#1}{\leavevmode\hbox\bgroup\setverbchar\readverb}%
    \bgroup\lccode`\~=`#1\lowercase{\egroup\def\readverb ##1~}{##1\egroup}%
    \catcode`#1=13
}
\def\setverbchar{\setverb
    \adef{ }{\ }%
    \adef{\{}{\bgroup\normalcolor\char`\{}%
    \adef{\}}{\char`\}\egroup}%
    \adef{[}{\bgroup\normalcolor\char`[}%
    \adef{]}{\char`]\egroup}%
    \ttfamily
}

\edef\bslash{\csstring\\}
\def\mytarget#1{\vbox to0pt{\kern-12pt \hypertarget{#1}{}\vss}}
\def\1{`}\def\2{``}
{\catcode`\`=13 \AtBeginDocument{%
    \verbchar{`}%
    \def\^`{\bgroup \docpoint}%
    \def\docpoint#1`{\egroup\leavevmode\edef\tmp{\csstring#1}%
        \expandafter\parsecn \string#1\relax
        \ifcsname \cn s/^\tmp\endcsname \else
            \mytarget{\cn s/^\tmp}\csgdef{\cn s/^\tmp}{}%
        \fi
        \hyperlink{\cn s/\tmp}{\color{blue}\ttfamily\if\cn c\bslash\fi\tmp}%
    }%
    \def\parsecn #1#2\relax {\edef\cn{\if\bslash#1c\else n\fi}}%
    \def\`{\bgroup \mainpoint}%
    \def\mainpoint #1`{\egroup\leavevmode\edef\tmp{\csstring#1}%
        \expandafter\parsecn \string#1\relax
        \ifcsname \cn s/\tmp\endcsname \moremainpoints \else
            \mytarget{\cn s/\tmp}\csgdef{\cn s/\tmp}{}%
        \fi
        \ifcsname \cn s/^\tmp\endcsname
            \hyperlink{\cn s/^\tmp}{\ttfamily\if\cn c\bslash\fi\tmp}\else
            {\color{red}\ttfamily\string#1}%
        \fi
    }%
    \def\moremainpoints{\errmessage{Second main documentation point \if\cn c\bslash\fi\tmp}}
}}

\let\lchevron=<
{\catcode`<=13 \AtBeginDocument{\catcode`<=13
    \def<#1>{\ifhmode\else\expandafter\hbox\fi\bgroup\normalcolor\ttfamily
             $\color{gray}\langle$\/{\rmfamily\itshape #1\/}$\color{gray}\rangle$\egroup}}%
}

\def\type#1{\NoCaseChange{\protect\typeA{#1}}}
\def\typeA#1{\bgroup\normalshape\ttfamily{\escapechar=-1\relax\expandafter}\detokenize{#1}\egroup}
% Precisamos de \m@ne aqui, \m@ne=-1.
\pdfstringdefDisableCommands{%
    \def\type#1{\bgroup\escapechar\m@ne\expandafter\egroup\detokenize{#1}}%
    \let\escapechar=\relax
    \let\m@ne=\relax
}

\definecolor{P}{RGB}{242, 121, 0}   % Primitivas.
\definecolor{C}{RGB}{23, 230, 23}   % Comentários.
\definecolor{B}{RGB}{245, 250, 247} % Fundo.

\lstset{
    language        = [latex]tex,
    basicstyle      = \abntsmall\singlesp\ttfamily,
    texcsstyle      = *\color{P}, % O listings não permite ":"
    commentstyle    = \color{C},  % no nome dessas cores. Bug.
    backgroundcolor = \color{B},
    frame           = leftline,
    rulecolor       = \color{lightgray},
    columns         = fullflexible,
    breaklines      = false,
    keepspaces      = true,
    aboveskip       = 0pt plus1.2ex minus.67ex,
    belowskip       = 0pt plus1.2ex minus.67ex,
    abovecaptionskip = 0pt,
    belowcaptionskip = 0pt,
%    tabsize         = 1, % Sem efeito, porque o documento usa Espaço na tabulação.
}

% #1 = linewidth
% #2 = border
% #3 = padding
% #4 = padding-top
% #5 = padding-right
% #6 = padding-bottom
% #7 = padding-left
% #8 = margin-left
% #9 = margin-right
\def\mylstset#1#2#3#4#5#6#7#8#9{\lstset{
    linewidth          = \dimexpr#1 -#2 -#3 -#5\relax,
    framerule          = \dimexpr#2\relax,
    framesep           = \dimexpr#3\relax,
    xleftmargin        = \dimexpr#2 +#3 +#7 +#8\relax,
    xrightmargin       = \dimexpr#9\relax,
    framextopmargin    = \dimexpr#4\relax,
    framexrightmargin  = \dimexpr#5\relax,
    framexbottommargin = \dimexpr#6\relax,
    framexleftmargin   = \dimexpr#7\relax,
}}
\mylstset{\linewidth}{1pt}{0pt}{0pt}{0pt}{0pt}{0pt}{0pt}{0pt}

\lstnewenvironment{lst}[1][\placepos]{%
    \let\medskipamount=\smallskipamount
    \place[#1]\vbox\bgroup
}{%
    \egroup\endplace
}

\def\hookXP{}
\lstnewenvironment{xp}[1][]{%
    \lstset{
        frame           = leftline,
        basicstyle      = \normalsize\ttfamily\singlesp\vskip-\baselineskip,
        texcsstyle      = \normalcolor,
        commentstyle    = \normalcolor,
        backgroundcolor = \color{white},
        gobble          = 4,
    }%
    \mylstset{\linewidth}{0pt}{0pt}{0pt}{0pt}{0pt}{\parindent}{0pt}{0pt}%
    \expandafter\lstset \expandafter{\hookXP,#1}%
    \endgraf\null
}{}

% ETC.
% ================================================

\setmainfont{XITS}[
    UprightFont    = *-Regular,
    BoldFont       = *-Bold,
    ItalicFont     = *-Italic,
    BoldItalicFont = *-BoldItalic,
    Extension      = .otf
]
\setmathfont{XITSMath-Regular.otf}
\setmonofont{InconsolataN}[
    UprightFont = *-Regular,
    BoldFont    = *-Bold,
    Extension   = .otf
]

\definelegendplace{code}{Código}{loc}
\definelegendplace{output}{Saída}{loo}
\definelegendplace{scr}{\emph{Screenshot}}{loscr}

\Crefname{code}{Código}{Códigos}
\Crefname{output}{Saída}{Saídas}
\Crefname{appendix}{Apêndice}{Apêndices}
\Crefname{codesection}{Seção}{Seções}
\Crefname{codesubsection}{Seção}{Seções}
\Crefname{scr}{\emph{Screenshot}}{\emph{Screenshots}}

\let\over=\@@over
\AtBeginDocument{\def\refname{Referências}} % Removido \^ da definição.
\def\hooksection{}
\def\tocsubsectionfont{\itshape}
\def\subsectionfont   {\itshape}
\AtBeginDocument{\pretocmd\appendix{\clearpage}{}{}}
\pretocmd\annex{\clearpage}{}{}
\def\etex{\leavevmode\hbox{$\varepsilon$-\TeX}}
\def\pdflatex{\leavevmode\hbox{PDF\LaTeX}}
\def\xelatex{\leavevmode\hbox{Xe\LaTeX}}
\def\lualatex{\leavevmode\hbox{Lua\LaTeX}}
\def\me{Elaboração própria.}
\def\bibfont{\raggedright\interlinepenalty=10000\singlesp\bibitemsep=\baselineskip}
\appto\textual{
    \def\@evenhead{\abntsmall \firstmarks0\hfil\thepage}
    \def\@oddhead {\abntsmall \botmarks1\hfil\thepage}
    \def\sectionmark   #1{\marks0{Seção \thesection: \unexpanded{#1}}}
    \def\subsectionmark#1{\marks1{Subseção \thesubsection: \unexpanded{#1}}}
    \preto\@oddhead {\lower4pt\rlap{\vrule width\hsize height.4pt}}
    \preto\@evenhead{\lower4pt\rlap{\vrule width\hsize height.4pt}}
}
\patchcmd{\eletroniclayout}{false}{true}{}{}
\let\onesidelayout=\eletroniclayout
\let\twosidelayout=\eletroniclayout
\hypersetup{
    pdfauthor   = Elayson Abreu,
    pdftitle    = Classe ABNTexto: normas ABNT em LaTeX,
    pdfsubject  = Manual da classe ABNTexto,
    pdfcreator  = LaTeX with abntexto,
    pdfkeywords = ABNTexto. LaTeX. Classe. ABNT,
}
\makeatother



\begin{document}
\leavevmode\vskip1.9in
\begingroup \centering \LARGE\sffamily\bfseries
    Classe ABNTexto\vskip-3pt \Large\itshape\mdseries
    Normas ABNT em {\upshape\LaTeX}\vskip2em \normalfont\large
    Elayson Abreu\vskip-3pt
    \href{mailto:abntexto.classe@gmail.com}%
                {abntexto.classe@gmail.com}\vskip2ex
    \today
\vskip2in\endgroup

\noindent\hfil {\Large\ttfamily 1.0}
\newpage

\nonum\notoc\section{Sumário}
\maketoc
\newpage

\section{Introdução}

O `abntexto` é uma classe do \LaTeX\ criada para facilitar a utilização das normas da Associação Brasileira de Normas Técnicas (ABNT) em trabalhos acadêmicos (monografias, dissertações, teses). Ela carrega os recursos necessários para a elaboração de trabalhos nessas normas.

Este \emph{não} é um manual sobre LaTeX (ou TeX, a ferramenta que o constitui) nem tampouco uma descrição detalhada sobre as normas ABNT. Para aprender sobre LaTeX leia \textcite{massago-2022-latex2e-via-exemplos}, em português, ou \textcite{martinsen-2022-latex2e-an-unnoficial}, disponível em inglês, espanhol e francês. Além disso, também se pode ler \textcite{olsak-2012-tex-in-nutshell} e \textcite{eijkhout-1992-tex-by-topic}, disponíveis em inglês, para aprender sobre o sistema \TeX.

Já existe uma classe para aplicação das normas ABNT cujo nome é `abntex2` \cite{araujo-2012-a-classe-abntex2}. Mas, diferente dessa, o `abntexto` não depende do `memoir`, uma classe de código-fonte e manual extensos e que implementa funções as quais, em sua maioria, não são pertinentes a um documento que use as diretrizes da ABNT.

As normas ABNT as quais o `abntexto` se propõe a seguir seguem listadas:

\begin{topics}
    \item NBR 14724:2024. Informação e documentação --- Trabalhos acadêmicos --- Apresentação \cite{abnt-2024-14724-trabalhos-academicos};
%    \item NBR 12225:2004. Informação e documentação --- Lombada --- Apresentação \cite{abnt-2004-12225-lombada};
%    \item NBR 6028:2021. Informação e documentação --- Resumo --- Apresentação \cite{abnt-2021-6028-resumo-resenha-recensao};
    \item NBR 6027:2012. Informação e documentação --- Sumário --- Apresentação \cite{abnt-2012-6027-sumario};
    \item NBR 10520:2023. Informação e documentação --- Citações em documentos --- Apresentação \cite{abnt-2023-10520-citacoes};
%    \item NBR 6023:2018. Informação e documentação --- Referências --- Elaboração \cite{abnt-2018-6023-referencias};
%    \item NBR 5892:2019. Informação e documentação --- Representação e formatos de tempo --- Datas e horas --- Apresentação \cite{abnt-2019-5892-representacao-e-formatos-de-tempo};
    \item NBR 6024:2012. Informação e documentação --- Numeração progressiva das seções de um documento --- Apresentação \cite{abnt-2012-6024-numeracao-progressiva-das};
    \item IBGE. Normas de apresentação tabular \cite{ibge-1993-normas-de-apresentacao}.
\end{topics}

\subsection{Nota de compatibilidade}

No salto da versão anterior, 3.2.1-beta, para esta, 4.0.6-beta, muito do código-fonte foi reescrito; funcionalidades foram adicionadas; comandos obsolesceram. Por consequência, documentos antigos irão certamente quebrar em \emph{layout} e aparência.

Para que os usuários tenham tempo de atualizar para esta nova versão, o arquivo `abntexto-3-2-1-beta.cls` foi deixado no repositório da classe no CTAN\footnote{\url{https://ctan.org/pkg/abntexto}.}. Todavia, para documentos antigos que já foram concluídos, não faz sentido atualizá-los para a versão atual. Nesse caso, seria pertinente copiar o arquivo `abntexto-3-2-1-beta.cls` para o mesmo diretório do documento cuja fase de escrita já foi concluída, chamá-lo por meio de `\documentclass{abntexto-3-2-1-beta}`, e, dessa forma, garantir que o resultado da compilação seja sempre o mesmo.

Para usuários que desejam atualizar seus documentos, o `abntexto` preservou a funcionalidade de comandos obsoletados até onde foi possível e muitas mensagens de alerta/erro foram criadas.

\subsection{Mudanças e novidades}

A sintaxe de `\place{<conteúdo>}` foi substituída por `\begin{place}[<opcional>]`. Além disso, após requisições, a classe ganhou suporte para objetos flutuantes. O posicionamento de `place` pode ser alterado em `<opcional>` com as palavras-chave de objetos flutuantes: `!, t, b, p, h`. Além destas, a classe criou uma nova: `here`\footnote{Semelhante a especificação `H` do pacote `float` \cite{float-package}.}. Com ela, o usuário impõe que `<conteúdo>` seja exibido no PDF na mesma posição em que foi inserido no documento.

Comandos da forma `\<ext>labelwidth` tais como `\toclabelwidth`, `\appendixlabelwidth`, entre outros, não precisam mais ser configurados, porque o `abntexto` alinha os rótulos de suas respectivas entradas automaticamente com a ajuda dos novos comandos `\eqbox` e `\eqboxsize`.

O pacote `enumitem` é carregado pela classe e o ambiente `topics` foi reescrito em função do referido pacote.

A partir de agora, use `\abntsmall` ao invés de `\small`. Por enquanto, o significado deste será igual ao daquele.

\subsection{\emph{Feedback}}

Sugestões, dúvidas, relatos de \emph{bugs} e doações são bem-vindas:\hfil\break {\ttfamily abntexto.classe@gmail.com}.

Chave PIX: {\ttfamily abntexto.classe@gmail.com}

\section{Chamando a classe}

O autor da classe recomenda o uso do \lualatex\ na compilação, embora o \pdflatex\ e o \xelatex\ também sejam suportados. Digite as linhas do \ref{cd:inicio} em conjunto com o \lualatex\ e use `abntexto`.

\legend{code}{Início}
\src\me
\label{cd:inicio}
\begin{lst}
    \documentclass{abntexto}

    \usepackage[english,brazil]{babel}
    \usepackage{fontspec}
    \usepackage{unicode-math}

    \setmainfont{XITS}[
        UprightFont    = *-Regular,
        BoldFont       = *-Bold,
        ItalicFont     = *-Italic,
        BoldItalicFont = *-BoldItalic,
        Extension      = .otf
    ]
    \setmathfont{XITSMath-Regular.otf}

    \begin{document}
        Texto
    \end{document}
\end{lst}

O arquivo `abntexto-exemplo.tex`, disponível no diretório da classe no CTAN, mostra como elaborar os elementos pré-textuais.

\section{Fonte tipográfica}

O tamanho da fonte é 12\,pt de imediato e pode ser configurado redefinindo `\normalsize` por meio de \^`\definesize`, cuja sintaxe é
\begin{xp}
    \definesize\<seletor>{<tamanho de fonte>}{<entrel. simples>}%
                         {<configs. extras>}
\end{xp}
O primeiro argumento é o nome do seletor (`\normalsize`, por exemplo). O segundo argumento é o tamanho da fonte. O terceiro é a entrelinha simples que deve ser selecionada levando em conta o valor do segundo argumento. Por exemplo, uma fonte com tamanho 10\,pt deve ter por volta de 12\,pt de entrelinha simples. O quarto argumento é para ajuste de valores dependentes do tamanho da fonte, por exemplo `\abovedisplayskip` ou `\bigskipamount`.

A classe `article`, na qual se baseia o `abntexto`, fornece os seletores habituais `\tiny`, `\footnotesize`, `\Large` etc. No entanto, esta classe se limitou a redefinir apenas dois: `\normalsize` para o tamanho padrão do texto e \^`\abntsmall` para satisfazer a exigência da \textcite{abnt-2024-14724-trabalhos-academicos} acerca do tamanho menor e uniforme. Até a versão anterior da classe, 3.2.1-beta, o comando `\small` era usado, mas foi substituído por `\abntsmall` para não haver conflito com o significado original de `\small` da classe `article`.

O espaçamento é de 1,5 por padrão e pode ser alterado com os atalhos \^`\singlesp` e \^`\onehalfsp` que alternam para espaçamento simples e de 1,5, respectivamente\footnote{O \type{\\doublesp} está obsoleto.}. A fórmula desses comandos foi atualizada, portanto documentos com versões antigas do `abntexto` terão o \emph{layout} alterado.

Nesta versão, `\spacing` está obsoleto porque esse comando ainda usa a fórmula antiga. Além disso, ele tem uma limitação: só pode ser usado após um seletor de tamanho, porque tais seletores redefinem `\baselineskip` diretamente. Por exemplo, `\abntsmall\spacing{<fator>}` é correto, porém `\spacing{<fator>}\abntsmall` é ineficaz. Essa limitação não existe ao usar os reimplementados `\singlesp` e `\onehalf` de forma que tanto `\singlesp\onehalfsp` quanto `\onehalfsp\singlesp` produzem o mesmo efeito.

\section{\emph{Layout}}

A folha é A4 por padrão e as margens esquerda e direita são de, respectivamente, 3\,cm, 2\,cm para páginas ímpares e de, respectivamente, 2\,cm, 3\,cm para páginas pares. Além disso, as margens superior e inferior têm, nessa ordem, 3\,cm, 2\,cm para todo o documento. Essas margens foram ajustadas com o pacote `geometry`, carregado na classe.

Existem dois comandos que configuram o \emph{layout} da página: \^`\pretextual` e \^`\textual`. O primeiro é iniciado automaticamente durante `\begin{document}`, invoca \^`\onesidelayout` e remove a numeração de página. O segundo, `\textual`, também é executado automaticamente pelo primeiro `\section` numerado, invoca \^`\twosidelayout` e ativa a numeração de página.

O comando `\onesidelayout` configura as páginas para impressão de somente um lado do papel (anverso) ao passo que `\twosidelayout` espelha as margens do documento, habilitando a impressão para os dois lados do papel (anverso e verso).

Os trabalhos acadêmicos podem vir a ser disponiblizados num repositório \emph{online}. Nesse caso, o usuário pode querer reconfigurar as margens de forma a se ter uma leitura mais confortável em dispositivos eletrônicos. Para esse propósito foi criado o comando \^`\eletroniclayout`, apesar dessa configuração não estar prevista na ABNT.

Há duas formas de ativá-lo. A primeira é, no preâmbulo, redefinir `\pretextual` e `\textual` substituindo as instruções `\onesidelayout` e `\twosidelayout` por `\eletroniclayout`. A segunda possibilidade é, também no preâmbulo, redefinir `\onesidelayout` e `\twosidelayout` diretamente conforme \ref{cd:eletroniclayout}. O comando `\eletroniclayout` centraliza o corpo do texto configurando 2,5\,cm nas margens esquerda e direita.

\legend{code}{\emph{Layout} para dispositivos eletrônicos}
\src\me
\label{cd:eletroniclayout}
\begin{lst}
    \documentclass{abntexto}
    \let\onesidelayout=\eletroniclayout
    \let\twosidelayout=\eletroniclayout
    \begin{document}
        texto
    \end{document}
\end{lst}

Outros comandos de estilo de página tais como `\pagestyle`\{{\ttfamily headings, empty, myheadings, plain}\} herdados da classe `article` não fazem sentido no `abntexto`. Ao invés de usá-los, é aconselhável redefinir `\textual`, `\pretextual` ou ambos.

\section{Sumário}

A classe `article` criou o comando `\tableofcontents` para impressão do Sumário, cuja definição será mantida por compatibilidade. No `abntexto`, porém, é recomendado usar \^`\maketoc`, cuja definição é ligeiramente diferente.

Os comandos que configuram a fonte tipográfica das entradas do Sumário são \^`\tocsectionfont`, \^`\tocsubsectionfont` etc. (ver \ref{cd:contents}). O último token nessas definições pode conter parâmetro, por exemplo:
\begin{xp}
    \def\tocsectionfont{\bfseries\MakeUppercase}
    \def\tocparagraphfont{\fbox}
\end{xp}

\legend{code}{Aparência do Sumário}
\src\me
\label{cd:contents}

\begin{lst}
    \documentclass{abntexto}
    \usepackage{xcolor}
    \def \tocsectionfont    {\color{red}\itshape\MakeUppercase}
    \def \tocsubsectionfont {\color{blue}\bfseries}

    \begin{document}
        \nonum\notoc\section{Sumário}
        \maketoc
        \newpage

        \section{Título}
        \subsection{Título}
    \end{document}
\end{lst}

\section{Secionamento}\label{sec:secionamento}

A personalização das macros de seção (\^`\section`, \^`\subsection`, \^`\subsubsection`, \^`\paragraph` e \^`\subparagraph`) é análoga a personalização do Sumário e deve ser feita redefinindo \^`\sectionfont`, \^`\subsectionfont`, \^`\subsubsectionfont`, \^`\paragraphfont` e \^`\subparagraphfont` como exemplificado em \ref{cd:sec} abaixo.

\src\me
\legend{code}{Aparência das seções}
\label{cd:sec}
\begin{lst}
    \documentclass{abntexto}
    \usepackage{xcolor}
    \def \sectionfont       {\color{blue}\scshape}
    \def \subsectionfont    {\color{red}\itshape}
    \def \subsubsectionfont {\color{yellow}\sffamily}

    \begin{document}
        \nonum\notoc\section{Sumário}
        \maketoc

        \section{Texto}
        \subsection{Texto}
        \subsubsection{Texto}
    \end{document}
\end{lst}

Também há como alterar o espaço vertical acima e abaixo das seções redefinindo \^`\abovesection`, \^`\belowsection`. Por exemplo, `\def\abovesubsection{2cm}` insere um espaço de 2\,cm acima de `\subsection`. Se o leitor deseja zerar essas instruções, deve escrever `\def\abovesubsection{0pt}` ao invés de `\def\abovesubsection{}`.

\subsection{O uso de \type{\\chapter} e \type{\\part}}

De acordo com a \textcite{abnt-2024-14724-trabalhos-academicos}: \enquote{[\dots] O trabalho acadêmico não pode ser dividido em capítulos; deve ser organizado em seções [\dots]}.

O usuário pode querer utilizar capítulos no trabalho, apesar de a unidade secional `\section` ser suficiente para um trabalho acadêmico. Para esse propósito, o `abntexto` disponibiliza a instrução \^`\chapter` para a utilização de capítulos no documento, porém não em conformidade com a \textcite{abnt-2024-14724-trabalhos-academicos}, mas no estilo da classe `book`. Para ativá-los é preciso chamar a instrução \^`\usechapters` no preâmbulo.

Também é possível usar \^`\part` em conjunto com o comando \^`\useparts`. Esse útlimo contém `\usechapters` em sua definição, portanto, ao chamar `\useparts`, não se deve chamar `\usechapters` para que este não sobrescreva as definições daquele.

Os níveis secionais `\section` e `\chapter` podem ser prefixados com \^`\nonum` que remove a numeração do título: `\nonum\section`. Os títulos prefixados com `\nonum` serão centralizados \cite{abnt-2024-14724-trabalhos-academicos}. Veja um exemplo de uso em \ref{cd:nonum}.

\legend{code}{Exemplo de título sem numeração e sem ingresso ao Sumário}
\src\me
\label{cd:nonum}

\begin{lst}
    \nonum\notoc\section{Resumo}
    Texto

    \nonum\notoc\section{Sumário}
    \maketoc
\end{lst}

As unidades secionais supracitadas também podem ser prefixadas com \^`\notoc` que instrui o título a não ingressar no Sumário: `\notoc\section`. Além disso, é possível usar ambos os prefixos ao mesmo tempo, conforme mostrado em \ref{cd:nonum}.

O primeiro parágrafo após uma seção é indentado por padrão, dispensando, assim, o uso do pacote `indentfirst`.

Infelizmente, comandos como `\printbibliography`, `\printglossary` e `\printindex` inserem o título internamente. Para que esse título seja exibido corretamente no `abntexto`, é necessário usar o ambiente \^`corrprint`, como demonstrado em \ref{cd:corrprint}.

\legend{code}{Ambiente \type{corrprint}}
\src\me
\label{cd:corrprint}

\begin{lst}
    \begin{corrprint}
        \printbibliography % ou \printglossary, \printindex e semelhantes
    \end{corrprint}
\end{lst}

Existe a possibilidade de o usuário desejar criar um Índice usando o ambiente `theindex` --- embora alternativas automatizadas sejam preferíveis \cite{martinsen-2022-latex2e-an-unnoficial}. Nesse caso, o ambiente `corrprint` funcionará como esperado conforme explicitado em \ref{cd:corrprint-theindex}. O ambiente `corrprint` também pode ser utilizado com o ambiente `thebibliography`, assim como no exemplo com `theindex`.

\legend{code}{Ambientes \type{corrprint} e \type{theindex}}
\src\me
\label{cd:corrprint-theindex}

\begin{lst}
    \begin{corrprint}
        \begin{theindex}
            \item carro, 1
                \subitem árvore, 2
            \indexspace
            \item mar, 3
                \subitem continente, 4
        \end{theindex}
    \end{corrprint}
\end{lst}

\section{Alíneas}

Use o ambiente \^`topics` para criar alíneas. Ele dispõe de dois níveis sendo que o segundo usa um travessão como marcador conforme o exemplo em \ref{cd:topics}. Naturalmente, é permitido usar `\label` após um item.

\legend{code}{Alíneas}
\src\me
\label{cd:topics}

\begin{lst}
    \begin{topics}
        \item \label{al:x} 1
            \begin{topics}
                \item \label{subal:x} x
                \item y
                \item z
            \end{topics}
        \item \label{al:y} 2
        \item 3
    \end{topics}
\end{lst}

\section{Áreas de legenda\label{sec:area-de-legenda}}

Use o ambiente \^`place` para inserir áreas de legenda. Ele substitui integralmente os ambientes `figure` e `table`, bem como o comando `\caption`, não suportados pela classe. Além disso, tem-se \^`\legend`\nobreak`{<tipo>}{<legenda>}` para titular; \^`\src`\nobreak`{<fonte>}` (\emph{source}) para inserir a fonte como mostrado em \ref{cd:place}. Porém, existem condições: `place` deve constar em último lugar e `\label` deve suceder `\legend`.

\legend{code}{Áreas de legenda}
\src\me
\label{cd:place}

\begin{lst}
    \begin{document}
        \legend{figure}{Um título}
        \src{\texttt{www.site.com.br}. Acesso em 01 jan. 2020.}
        \label{fig:teste}
        \begin{place}
            \includegraphics[width=0.4\linewidth]{example-image}
        \end{place}

        Ver \ref{fig:teste}.
    \end{document}
\end{lst}

Além disso, existe a possibilidade de alterar o espaço vertical acima e abaixo de uma área de legenda com `\def`\^`\aboveplace``{<valor>}` e `\def`\^`\belowplace``{<valor>}`. Todavia, tais comandos não têm efeito em objetos flutuantes. Para modificar o espaçamento em `floats` o leitor pode fornecer valores diferentes para `\floatsep`, `\intextsep` e `\textfloatsep`, definidos no LaTeX.

A partir desta versão, o `abntexto` passou a suportar objetos flutuantes. Por exemplo, em
\begin{xp}
    \begin{place}[<especificação>]
\end{xp}
A `<especificação>` é composta pelas palavras-chave habituais: `!`, `t`, `b`, `p`, e `h`. Alternativamente, é possível usar a palavra-chave `here`, definida nesta classe. Com ela, o conteúdo de `place` é mostrado na página exatamente onde foi escrito no documento.

A especificação padrão é `here` caso nenhuma seja explicitamente assumida, mas o usuário pode mudar o padrão digitando, por exemplo, `\def`\nobreak\^`\placepos``{tbp}` no preâmbulo do documento.

\subsection{Criando áreas de legenda}

Até a versão anterior, 3.2.1-beta, o `abntexto` fornecia \^`\definelegendplace` para criação de novos espaços de legenda sob a sintaxe
\begin{xp}
    \definelegendplace{<tipo>}{<nome>}{<ref rótulo>}{<extensão>}
\end{xp}
Nesta versão, no entanto, o terceiro argumento foi excluído em favor do suporte mais extensivo do pacote `cleveref.sty` (não carregado pela classe). Veja o \ref{app:cleveref} para mais detalhes. A partir de agora, a nova sintaxe é
\begin{xp}
    \definelegendplace{<tipo>}{<nome>}{<extensão>}
\end{xp}
Esta classe continuará fornecendo suporte sintático, porém sem funcionalidade, para a versão anterior do referido comando, mas vai exibir uma mensagem de alerta pedindo para que o terceiro argumento seja removido.

Acerca do novo `\definelegendplace`: a instrução
\begin{xp}
    \definelegendplace{diagram}{Diagrama}{lod},
\end{xp}
por exemplo, criará o comando `\makelod` para escrever no documento as entradas do arquivo auxiliar `.lod`. Além disso, o tipo de legenda `diagram` estará disponível para ser usado em `\legend`: `\legend{diagram}{<título>}`.

\subsection{Figuras lado a lado}

O `abntexto` tem suporte nativo para inserção de figuras lado a lado com o ambiente \^`multiplace`. Elas se comportam como caracteres então podemos usar `\hfil` para centralizá-las. Veja o exemplo no \ref{cd:multiplace} e o resultado em \ref{out:result}.

\legend{code}{Figuras lado a lado}
\src\me
\label{cd:multiplace}
\begin{lst}
    \legend{figure}{Duas figuras}
    \src{Elaboração do autor.}

    \begin{multiplace}
        \sublegend{Um título} \label{fig:teste1}
        \subsrc{\texttt{www.site.com.br}. Acesso em 01 jan. 2020.}
        \hfil
        \begin{subplace}[.25\linewidth]
            \includegraphics[width=\linewidth]{example-image-a}
        \end{subplace}
        %
        \sublegend{Outro título} \label{fig:teste2}
        \subsrc{\texttt{www.site.com.br}. Acesso em 01 jan. 2020.}
        \hfil
        \begin{subplace}[.25\linewidth]
            \includegraphics[width=\linewidth]{example-image-b}
        \end{subplace}
    \end{multiplace}

    \parindent=1.5cm Ver \ref{fig:teste1} e \ref{fig:teste2}.
\end{lst}

\legend{output}{Resultado do \ref{cd:multiplace}}
\label{out:result}
\src\me
\begin{place}\vbox\bgroup
    \legend{figure}{Duas figuras}
    \src{Elaboração do autor.}
    %
    \begin{multiplace}
        \sublegend{Um título} \label{fig:teste1}
        \subsrc{\texttt{www.site.com.br}. Acesso em 01 jan. 2020.}
        \hfil
        \begin{subplace}[.25\linewidth]
            \includegraphics[width=\linewidth]{example-image-a}
        \end{subplace}
        %
        \sublegend{Outro título} \label{fig:teste2}
        \subsrc{\texttt{www.site.com.br}. Acesso em 01 jan. 2020.}
        \hfil
        \begin{subplace}[.25\linewidth]
            \includegraphics[width=\linewidth]{example-image-b}
        \end{subplace}
    \end{multiplace}

    \parindent=1.5cm Ver \ref{fig:teste1} e \ref{fig:teste2}.
\egroup\end{place}

Há uma legenda principal impressa com `\legend` e duas sublegendas para cada objeto inseridas com \^`\sublegend`\nobreak`{<título>}`, bem como um par de ambientes \^`subplace`\footnotemark. Além disso, a ordem das instruções é similar àquela aplicada no ambiente `place`: `subplace` deve ser posto após ambos `\sublegend` e \^`\subsrc``{<fonte>}`, e `\label` deve suceder `\legend`.

\footnotetext{A sintaxe \type{\\subplace{<conteúdo>}} está obsoleta. Use `subplace` na forma de ambiente: \type{\\begin{subplace}}.}

Linhas em branco produzem quebra de parágrafo em `multiplace` e, provavelmente, não são desejadas.

Normalmente, espaços em branco depois de `\end{<env>}` são significativos, mas, no caso do ambiente `subplace`, isso não é verdade, porque existe um `\ignorespacesafterend` em sua definição, um comando definido no LaTeX.

Existe, ainda, a possibilidade de um arranjo diferente no `multiplace` usando o parâmetro opcional de `\sublegend`, como se pode ver no \ref{cd:multialt} em conjunto com o resultado em \ref{out:alt}.

\legend{code}{Uso alternativo de \type{multiplace}}
\src{Elaborado pelo autor.}
\label{cd:multialt}
\begin{lst}
    \begin{multiplace}
        \sublegend[table]{Tabela de exemplo} \label{tab1}
        \subsrc{\me}
        \hfil
        \begin{subplace}[.25\linewidth]
            \includegraphics[width=\linewidth]{example-image-a}
        \end{subplace}
        %
        \sublegend[figure]{Título de figura} \label{fig1}
        \subsrc{\me}
        \hfil
        \begin{subplace}[.25\linewidth]
            \includegraphics[width=\linewidth]{example-image-b}
        \end{subplace}
        %
        \sublegend[figure]{Outra figura} \label{fig2}
        \subsrc{\me}
        \vadjust{\kern10pt}\\ \null \hfil
        \begin{subplace}[.25\linewidth]
            \includegraphics[width=\linewidth]{example-image-c}
        \end{subplace}
    \end{multiplace}
\end{lst}

\src{\me}
\begin{place}\vbox\bgroup
\legend{output}{Resultado do \ref{cd:multialt}} \label{out:alt}
\begin{multiplace}
    \sublegend[table]{Tabela de exemplo} \label{tab1}
    \subsrc{\me}
    \hfil
    \begin{subplace}[.25\linewidth]
        \includegraphics[width=\linewidth]{example-image-a}
    \end{subplace}
    \sublegend[figure]{Título de figura} \label{fig1}
    \subsrc{\me}
    \hfil
    \begin{subplace}[.25\linewidth]
        \includegraphics[width=\linewidth]{example-image-b}
    \end{subplace}
    \sublegend[figure]{Outra figura} \label{fig2}
    \subsrc{\me}
    \vadjust{\kern10pt}\\ \null \hfil
    \begin{subplace}[.25\linewidth]
        \includegraphics[width=\linewidth]{example-image-c}
    \end{subplace}
\end{multiplace}
\egroup\end{place}

A \ref{out:alt} possui a \ref{tab1}, \ref{fig1} e \ref{fig2}.

\subsection{Tabelas que se partem entre páginas}

No momento não existe suporte para legendar tabelas e outras estruturas que se partem entre páginas. Se o leitor tiver de usar `longtable.sty` \cite{longtable-package} ou `tabularray.sty` \cite{tabularray-package} (este último é mais flexível), será preciso legendar estruturas tabulares com as ferramentas fornecidas por esses pacotes.

\section{Formatação de \type{\\cite} e Referências}

O `abntexto` não possui macros para formatação de `\cite` e Referências. Utilize `\usepackage[style=abnt]{biblatex}`.

\section{Citações}

Visando seguir a norma \textcite{abnt-2023-10520-citacoes}, são fornecidos `\enquote` do pacote `csquotes`, carregado pela classe, e \^`\Enquote`. Para inserir aspas duplas em citações diretas, de até três linhas, use `\enquote{<texto>}\cite{<chave>}`. Porém, se tratando de citações diretas, com mais de três linhas, de parágrafo único, alterne para `\Enquote{<texto> \cite{<chave>}.}`.

\section{Apêndice e Anexo}

Com os comandos \^`\appendix` e \^`\annex` você pode inserir apêndices e anexos junto dos indicativos desejados: \enquote{APÊNDICE 1 --- \dots} ou \enquote{ANEXO 1 --- \dots} (ver \ref{cd:appendix}).

\legend{code}{Anexos e apêndices}
\src{\me}
\label{cd:appendix}

\begin{lst}
    \begin{document}
        \appendix{Lorem}
        \appendix{Lipsum}

        \annex{Lorem}
        \annex{Lipsum}
    \end{document}
\end{lst}

Até a versão anterior desta classe, era necessário ajustar a largura dos rótulos no Sumário manualmente redefinindo os comandos \^`\appendixlabelwidth` e \^`\annexlabelwidth`. Todavia, nesta versão, 4.0.6-beta, eles se ajustam automaticamente com o auxílio das novas macroinstruções \^`\eqbox` e \^`\eqboxsize`.

\section{Macros para identificação do documento}

\subsection{Ficha catalográfica}

Para produzir fichas catalográficas foi criado o ambiente \^`indexcard`. Ele sempre esteve disponível na forma `\indexcard{<conteúdo>}`, agora obsoleta, mas nunca foi devidamente documentado neste manual. Eis uma amostra de uso:
\begin{xp}
    \documentclass{abntexto}
    \usepackage{lipsum}
    \usepackage{xcolor}
    \begin{document}
        \def\icardWidth{12.5cm}
        \def\icardHeight{7.5cm}
        \def\icardPadding{1cm}
        \def\icardBorder{2pt}
        \appto\icardFont{\color{magenta}}
        \icardPrintGrid

        \noindent\hfil
        \begin{indexcard}
            \lipsum[1][1-4]
        \end{indexcard}
    \end{document}
\end{xp}
Que produz:

\bgroup
\def\icardWidth{12.5cm}
\def\icardHeight{7.5cm}
\def\icardPadding{1cm}
\def\icardBorder{2pt}
\appto\icardFont{\color{magenta}}
\icardPrintGrid
\noindent\hfil
\begin{indexcard}
    \lipsum[1][1-4]
\end{indexcard}
\egroup

Perceba que a instrução \^`\icardPrintGrid` desenha uma grade de quatro linhas para delimitar os contornos do bloco de texto, no intuito de verificar se ele não ultrapassou as margens internas. A definição original de `\icardFont` é
\begin{xp}
    \def\icardFont{\raggedright\parindent=2em\ttfamily}
\end{xp}

\subsection{Utilitário \type{\\Enter}}

O `abntexto` criou o comando \^`\Enter``[<número>]` para produzir um espaço vertical equivalente à tecla `Enter` do teclado. Ele é usado com frequência em `abntexto-exemplo.tex`. O argumento opcional `<número>` é o número de `Enters` que será impresso, por exemplo, `\Enter[2]` imprime um espaço vertical equivalente a dois `Enters` consecutivos. Como `\Enter` é, basicamente, `\vskip\baselineskip`, precedê-lo com `\par` não é necessário, porque `\vskip` possui, a grosso modo, um `\par` em seu funcionamento.

\subsection{Desativando a hifenização}

O comando \^`\nohyph` impede a hifenização dos parágrafos integralmente, sendo útil em títulos e parágrafos curtos. Seu uso pode ser feito assim:
\begin{xp}
    {\nohyph Parágrafo curto ou título\par} % Esse \par é necessário para
                                            % que a formatação do parágrafo
                                            % entre em ação.
\end{xp}

\subsection{Assinatura da banca}

Use \^`\judgeline``{<arg>}` para assinaturas na Folha de Aprovação. Exemplo:
\begin{xp}
    \judgeline{Prof. Dr. Nome Sobrenome \\ Instituição}\Enter
    \judgeline{Prof. Dr. Nome Sobrenome \\ Instituição}\Enter
    \judgeline{Prof. Dr. Nome Sobrenome \\ Instituição}
\end{xp}
Resultado:
\judgeline{Prof. Dr. Nome Sobrenome \\ Instituição}\Enter
\judgeline{Prof. Dr. Nome Sobrenome \\ Instituição}\Enter
\judgeline{Prof. Dr. Nome Sobrenome \\ Instituição}

Veja outro exemplo de `\judgeline` em `abntexto-exemplo.tex`.

\newpage
\begin{corrprint}
    \printbibliography
\end{corrprint}

\newrefsection

\appendix{Cleveref\label{app:cleveref}}

Até a versão passada, 3.2.1-beta, o `abntexto` tinha o poder de inferir o tipo de `\ref` e inserir o substantivo correspondente com auxílio do terceiro parâmetro de `\definelegendplace`. No entanto, esse parâmetro foi removido deixando o referido comando com três argumentos apenas. Como resgatar essa funcionalidade de inferência em `\ref`? Usando o pacote `cleveref.sty` e sua macro `\cref` \cite{cleveref-package}. Veja o exemplo em \ref{cd:cleveref}, cujo resultado é \enquote{Ler Seção 1}.

\legend{code}{Exemplo de uso do \type{cleveref.sty}}
\src\me
\label{cd:cleveref}
\begin{lst}
    \documentclass{abntexto}
    \usepackage[brazil]{babel}
    \usepackage[brazilian,capitalize,noabbrev]{cleveref}
    \begin{document}
        \section{Um título de exemplo\label{sec:ex}}
        Ler \cref{sec:ex}.
    \end{document}
\end{lst}

Percebe-se que a chave `brazil` não é reconhecida pelo `cleveref` compelindo-nos a usar `brazilian`. Além disso, ele deve ser carregado depois do `hyperref.sty`. Veja outras funcionalidades no manual do pacote.

\appendix{Lista de Abreviaturas}

Elemento opcional. O `glossaries-extra.sty` \cite{glossaries-extra-package} consegue criar uma Lista de Abreviaturas, conforme \ref{cd:abbrev}.

\legend{code}{Abreviaturas}
\src\me
\label{cd:abbrev}

\begin{lst}
    \documentclass{abntexto}

    \usepackage[brazil]{babel}
    \usepackage{hyperref}
    \usepackage[automake,shortcuts=all,abbreviations]{glossaries-extra}\makeglossaries
    \usepackage[brazilian,capitalize,noabbrev]{cleveref}

    \newabbreviation[category=abbreviations]{fil}{Fil.}{Filosofia}
    \newabbreviation[category=abbreviations]{el sf}{el.\,sf.}{elemento substantivo feminino}
    \newabbreviation[category=abbreviations]{eng el}{Eng.\,Elétr.}{Engenharia Elétrica}
    \newabbreviation[category=abbreviations]{med legal}{Med.\,Leg.}{Medicina Legal}

    \begin{document}
        Abreviaturas: \ab{fil},
                      \ab{el sf},
                      \ab{eng el},
                      \ab{med legal}.

        \begin{corrprint}\notoc % "Lista de Abreviaturas" não deve constar no Sumário.
            \printglossary[type=abbreviations]
        \end{corrprint}
    \end{document}
\end{lst}

Note que o `cleveref.sty` deve ser carregado depois de `glossaries-extra.sty`, e o `hyperref.sty`, antes.

\appendix{Lista de Siglas}

Elemento opcional. O `glossaries-extra.sty` \cite{glossaries-extra-package} também consegue criar uma Lista de Siglas, conforme \ref{cd:acro}.

\legend{code}{Siglas}
\src\me
\label{cd:acro}

\begin{lst}
    \documentclass{abntexto}

    \usepackage[brazil]{babel}
    \usepackage{hyperref}
    \usepackage[automake,shortcuts=all,acronym]{glossaries-extra}\makeglossaries
    \usepackage[brazilian,capitalize,noabbrev]{cleveref}

    \setabbreviationstyle[acronym]{long-short}

    \newacronym{ibge}{IBGE}{Instituto Brasileiro de Geografia e Estatística}
    \newacronym{html}{HTML}{HyperText Markup Language}
    \newacronym{ce}{CE}{Ceará}
    \newacronym{edo}{EDO}{Equação Diferencial Ordinária}

    \begin{document}
        1ª chamada: \ac{ibge}, \ac{html}, \ac{ce}, \ac{edo}.

        2ª chamada: \ac{ibge}, \ac{html}, \ac{ce}, \ac{edo}.

        \begin{corrprint}\notoc % "Lista de Siglas" não deve constar no Sumário.
            \printglossary[type=acronym]
        \end{corrprint}
    \end{document}
\end{lst}

Note que o `cleveref.sty` deve ser carregado depois de `glossaries-extra.sty`, e o `hyperref.sty`, antes.

\appendix{Lista de Símbolos}

Elemento opcional. O `glossaries-extra.sty` \cite{glossaries-extra-package} também consegue criar uma Lista de Símbolos, conforme \ref{cd:simbolos}.

\legend{code}{Símbolos}
\src\me
\label{cd:simbolos}

\begin{lst}
    \documentclass{abntexto}

    \usepackage[brazil]{babel}
    \usepackage{hyperref}
    \usepackage[automake,shortcuts=all,symbols]{glossaries-extra}\makeglossaries
    \usepackage[brazilian,capitalize,noabbrev]{cleveref}

    \def\newsymbol#1#2#3#4{\newglossaryentry{#1}{
        name={#2},sort={#3},description={#4},type=symbols,category=symbol}}

    \newsymbol{alpha}{\ensuremath{\alpha}}{a}{Letra grega alfa}
    \newsymbol{o(n)}{\ensuremath{O(n)}}{On}{Ordem de um algoritmo}
    \newsymbol{i}{\ensuremath{i}}{i}{Unidade imaginária}
    \newsymbol{A}{\ensuremath{\mathbf{A}}}{A}{Matriz $\mathbf{A}$}
    \newsymbol{N}{N}{N}{Newton}

    \def\symb#1{\gls{#1}}

    \begin{document}
        Símbolos: \symb{alpha}, \symb{o(n)}, \symb{i}, \symb{A}, \symb{N}.

        \begin{corrprint}\notoc % "Lista de Símbolos" não deve constar no Sumário.
            \printglossary[type=symbols]
        \end{corrprint}
    \end{document}
\end{lst}

Note que o `cleveref.sty` deve ser carregado depois de `glossaries-extra.sty`, e o `hyperref.sty`, antes.

\appendix{Glossário}

Elemento opcional. Você pode criar glossários com `glossaries-extra` \cite{glossaries-extra-package}, como em \ref{cd:gloss}.

\legend{code}{Glossário}
\src\me
\label{cd:gloss}

\begin{lst}
    \documentclass{abntexto}

    \usepackage[brazil]{babel}
    \usepackage{hyperref}
    \usepackage[automake,shortcuts=all]{glossaries-extra}\makeglossaries
    \usepackage[brazilian,capitalize,noabbrev]{cleveref}

    \newglossaryentry{metalismo}{
        name        = metalismo,
        description = {Tipo de sistema monetário cujo valor da moeda,
                       geralmente de uma nação ou de um país, é determinado
                       pela quantidade constante de certo metal (ouro ou
                       prata), desta forma se estabelece um valor fixo de
                       troca entre esses metais e o próprio dinheiro.}
    }
    \newglossaryentry{backend}{
        name        = \emph{backend},
        sort        = backend,
        description = {A parte de um site ou aplicação que não é visível
                       para o usuário e que gerencia a lógica de negócios
                       e a comunicação com o banco de dados.}
    }
    \newglossaryentry{cache}{
        name        = cache,
        description = {Um armazenamento temporário de dados para acesso
                       rápido. Ajuda a reduzir o tempo de carregamento.}
    }
    \newglossaryentry{autovetor}{
        name        = autovetor,
        description = {Um vetor não nulo que, quando multiplicado por uma
                       matriz, resulta em um vetor que é um múltiplo escalar do
                       vetor original.}
    }

    \begin{document}
        Glossário: \gls{metalismo}, \gls{backend}, \gls{cache}, \gls{autovetor}.

        \begin{corrprint}
            \printglossary[type=main]
        \end{corrprint}
    \end{document}
\end{lst}

Note que o `cleveref.sty` deve ser carregado depois de `glossaries-extra.sty`, e o `hyperref.sty`, antes.

\appendix{Ambiente \type{place} e o pacote \type{listings}}

O `listings` é um impressor de código-fonte para LaTeX \cite{listings-package}. Você pode usar `\legend` em conjunto com `listings` como mostrado em \ref{cd:listings}. Há algumas observações quanto a essa solução: primeiro, o `listings` não permite que `place` esteja na forma convencional, mas na forma `\place`-`\endplace`. Segundo, a região em \emph{source} terá largura fixa de `\linewidth`, que é a largura padrão de uma `\vbox`.

\legend{code}{\type{\\legend} em conjunto com \type{listings}}
\src\me
\label{cd:listings}

\lstnewenvironment{lstn}{%
    \let\medskipamount=\smallskipamount
    \place\vbox\bgroup
}{%
    \egroup\endplace
}

\begin{lstn}
    % No preâmbulo:
    \usepackage{listings}

    \lstnewenvironment{lst}[1][\placepos]{%
        \let\medskipamount=\smallskipamount
        \place[#1]\vbox\bgroup
    }{%
        \egroup\endplace
    }
    \definelegendplace{code}{Código}{loc}

    % No documento:
    \nonum\notoc\section{Lista de códigos}
    \makeloc

    \legend{code}{Título}
    \src{Fonte}
    \label{Chave}

    \begin{lst}
        int x = 1;
    \end{lst}
\end{lstn}

Na solução acima o ambiente `lst` é inquebrável, portanto quebras de páginas não são possíveis. Se você precisa de um ambiente que se parta entre páginas, veja a solução abaixo (\ref{cd:longlst}):

\legend{code}{O ambiente \type{longlst} quebrável como alternativa ao \type{lst}}
\src\me
\label{cd:longlst}

\begin{lst}
    \lstnewenvironment{longlst}[1][\placepos]{%
        \begingroup
        \processplacearg{#1}%
        \parskip=0pt \parindent=0pt
        \initplace
        \printlegendbox
        \endgraf\nointerlineskip \nobreak\vskip\medskipamount
        \resetORIlabel \ignorespaces
    }{%
        \unskip
        \savedplacewidth=\hsize
        \printsrcbox
        \finishplace
        \resetplace \gresetORIlabel
        \endgroup
    }
\end{lst}

\appendix{Ambiente \type{place} e o pacote \type{algorithm2e}}

O pacote `algorithm2e` foi criado para geração de algoritmos em LaTeX \cite{algorithm2e-package}. É possível usar `place` em conjunto com `algorithm2e` (\ref{cd:algorithm2e}):

\legend{code}{\type{algorithm2e} em \type{place}}
\src\me
\label{cd:algorithm2e}

\begin{lst}
    % No preâmbulo
    \usepackage{algorithm2e}

    \definelegendplace{alg}{Algoritmo}{loa}

    % No documento
    \nonum\notoc\section{Lista de algoritmos}
    \makeloa

    \legend{alg}{How to write algorithms}
    \src{Elaboração própria}
    \label{alg:chave}

    \begin{place}
    \begin{algorithm}[H]
    \SetAlgoLined
    \KwData{this text}
    \KwResult{how to write algorithm with \LaTeX2e }
    initialization\;
    \While{not at end of this document}{
        read current\;
        \eIf{understand}{
            go to next section\;
            current section becomes this one\;
        }{
            go back to the beginning of current section\;
        }
    }
    \end{algorithm}
    \end{place}
\end{lst}

A especificação `H` é obrigatória.

% MACROS PARA O CÓDIGO-FONTE
% ================================================

\pdfstringdefDisableCommands{\def\marks#1#2{}}
\appendix{\marks0{}\marks1{}Código-fonte\label{ap:codigo-fonte}}

\makeatletter
\newcounter{codesection}
\newcounter{codesubsection}[codesection]
\let\c@section=\c@codesection
\let\c@subsection=\c@codesubsection
\def\thecodesection{\arabic{codesection}}
\def\thecodesubsection{\thecodesection.\arabic{codesubsection}}
\edef\toclevel@codesection{\mainseclevel}
\edef\toclevel@codesubsection{\the\numexpr\mainseclevel-1\relax}

\appto\hooktocsection{\def\extleaders{}%
    \edef\savedskips{\rightskip=\the\rightskip \parfillskip=\the\parfillskip\relax}}
\appto\hooktocsubsection{\def\extleaders{}%
    \edef\savedskips{\rightskip=\the\rightskip \parfillskip=\the\parfillskip\relax}}
\appto\hookextline{\savedskips}

\def\toclabelbox{\lowercase{\eqbox{codetoc}}}
\def\tocsectionfont{\itshape\MakeUppercase}
\def\tocsubsectionfont{}
\def\sectionfont{\tocsectionfont}
\def\subsectionfont{\itshape}

% O \addcontentsline{#1}{#2}{#3} está programado para
% enviar marcadores apenas se #1 = toc. Podemos
% mudar isso com \hypersetup{bookmarkstype=<outro>}.
\hypersetup{bookmarkstype=codetoc}
\patchcmd{\targetsection}{\refstepcounter{section}}{\refstepcounter{codesection}}{}{}
\patchcmd{\targetsection}{\addcontentsline{toc}}{\addcontentsline{codetoc}}{}{}
\patchcmd{\targetsubsection}{\refstepcounter{subsection}}{\refstepcounter{codesubsection}}{}{}
\patchcmd{\targetsubsection}{\addcontentsline{toc}}{\addcontentsline{codetoc}}{}{}

\def\hooksection{}
\def\codesection{\section}
\def\codesubsection{\subsection}

\colorlet{C}{.}
\definecolor{tmp}{RGB}{60, 90, 110}  \colorlet{P}{tmp}
                                     \colorlet{B}{blue!7}

\expanded{\noexpand\lstdefinelanguage{mylang}{
    language    = [latex]tex,
    texcs       = {\directlua{for k,v in pairs(tex.primitives())
                   do tex.print(k, v.. ',')end}},
    deletetexcs = {end}
}}

\lstset{
    language    = mylang,
    frame       = single,
    numbers     = left,
    numbersep   = \dimexpr-5pt+\addnumbersep\relax,
    numberstyle = \fontsize{8bp}{9bp}\selectfont,
    literate    = {\^^M\%.}{}1,
}
\def\addnumbersep{9pt}
\mylstset{\linewidth}{0pt}{0pt}{3pt}{0pt}{3pt}{13pt+\addnumbersep}{0pt}{0pt}

\newcount\codecount \codecount=1
\newbox\mybox
\setbox\mybox=\hbox to\hsize{\hfil\abntsmall\ttfamily\bfseries\color{gray}abntexto.cls}
\lstnewenvironment{code}[1][]{%
    \lstset{
        title       = \copy\mybox\nobreak,
        firstnumber = \codecount,
        #1
    }%
}{%
    \global\codecount=\value{lstnumber}\relax
}
\appto\hookXP{
    frame     = none,
    numbers   = none,
    numbersep = 0pt,
    gobble    = 4,
}

\parskip=0pt plus.2ex minus.1ex
\@starttoc{codetoc}
{\def\excl{!}\catcode\1\!=14 \catcode\1\%=9 %! Public Domain Software 2022-2026
%! 2026-03-04
%! 1.0
%! Elayson Abreu
%! abntexto.classe@gmail.com

% !TeX TS-program = lualatex
\documentclass{abntexto}

\usepackage[german,english,brazil,shorthands=off]{babel}
\usepackage[cmyk]{xcolor}
\usepackage{lipsum}
\usepackage[
    style       = abnt,
    maxbibnames = 100,
    extrayear   = true,
]{biblatex} \addbibresource{abntexto.bib}
\usepackage{unicode-math}
\usepackage[colorlinks,linktoc=page]{hyperref}
\usepackage{microtype}
\usepackage[brazilian,capitalize,noabbrev]{cleveref}\AtBeginDocument{\def\ref#1{\cref{#1}}}
\usepackage{fontspec}
\usepackage{listings}
\usepackage{lua-widow-control} \clubpenalty=10000\widowpenalty=10000

\makeatletter

% VERBATIM
% ================================================

\def\adef#1{\catcode`#1=13 \begingroup \lccode`\~=`#1\lowercase{\endgroup\def~}}
\def\setverb{\def\do##1{\catcode`##1=12}\dospecials\space}
\def\verbchar#1{%
    \ifx\savedttchar\undefined\else \catcode\savedttchar=\savedttcharc \fi
    \chardef\savedttchar=`#1\relax
    \chardef\savedttcharc=\catcode`#1\relax
    \adef{#1}{\leavevmode\hbox\bgroup\setverbchar\readverb}%
    \bgroup\lccode`\~=`#1\lowercase{\egroup\def\readverb ##1~}{##1\egroup}%
    \catcode`#1=13
}
\def\setverbchar{\setverb
    \adef{ }{\ }%
    \adef{\{}{\bgroup\normalcolor\char`\{}%
    \adef{\}}{\char`\}\egroup}%
    \adef{[}{\bgroup\normalcolor\char`[}%
    \adef{]}{\char`]\egroup}%
    \ttfamily
}

\edef\bslash{\csstring\\}
\def\mytarget#1{\vbox to0pt{\kern-12pt \hypertarget{#1}{}\vss}}
\def\1{`}\def\2{``}
{\catcode`\`=13 \AtBeginDocument{%
    \verbchar{`}%
    \def\^`{\bgroup \docpoint}%
    \def\docpoint#1`{\egroup\leavevmode\edef\tmp{\csstring#1}%
        \expandafter\parsecn \string#1\relax
        \ifcsname \cn s/^\tmp\endcsname \else
            \mytarget{\cn s/^\tmp}\csgdef{\cn s/^\tmp}{}%
        \fi
        \hyperlink{\cn s/\tmp}{\color{blue}\ttfamily\if\cn c\bslash\fi\tmp}%
    }%
    \def\parsecn #1#2\relax {\edef\cn{\if\bslash#1c\else n\fi}}%
    \def\`{\bgroup \mainpoint}%
    \def\mainpoint #1`{\egroup\leavevmode\edef\tmp{\csstring#1}%
        \expandafter\parsecn \string#1\relax
        \ifcsname \cn s/\tmp\endcsname \moremainpoints \else
            \mytarget{\cn s/\tmp}\csgdef{\cn s/\tmp}{}%
        \fi
        \ifcsname \cn s/^\tmp\endcsname
            \hyperlink{\cn s/^\tmp}{\ttfamily\if\cn c\bslash\fi\tmp}\else
            {\color{red}\ttfamily\string#1}%
        \fi
    }%
    \def\moremainpoints{\errmessage{Second main documentation point \if\cn c\bslash\fi\tmp}}
}}

\let\lchevron=<
{\catcode`<=13 \AtBeginDocument{\catcode`<=13
    \def<#1>{\ifhmode\else\expandafter\hbox\fi\bgroup\normalcolor\ttfamily
             $\color{gray}\langle$\/{\rmfamily\itshape #1\/}$\color{gray}\rangle$\egroup}}%
}

\def\type#1{\NoCaseChange{\protect\typeA{#1}}}
\def\typeA#1{\bgroup\normalshape\ttfamily{\escapechar=-1\relax\expandafter}\detokenize{#1}\egroup}
% Precisamos de \m@ne aqui, \m@ne=-1.
\pdfstringdefDisableCommands{%
    \def\type#1{\bgroup\escapechar\m@ne\expandafter\egroup\detokenize{#1}}%
    \let\escapechar=\relax
    \let\m@ne=\relax
}

\definecolor{P}{RGB}{242, 121, 0}   % Primitivas.
\definecolor{C}{RGB}{23, 230, 23}   % Comentários.
\definecolor{B}{RGB}{245, 250, 247} % Fundo.

\lstset{
    language        = [latex]tex,
    basicstyle      = \abntsmall\singlesp\ttfamily,
    texcsstyle      = *\color{P}, % O listings não permite ":"
    commentstyle    = \color{C},  % no nome dessas cores. Bug.
    backgroundcolor = \color{B},
    frame           = leftline,
    rulecolor       = \color{lightgray},
    columns         = fullflexible,
    breaklines      = false,
    keepspaces      = true,
    aboveskip       = 0pt plus1.2ex minus.67ex,
    belowskip       = 0pt plus1.2ex minus.67ex,
    abovecaptionskip = 0pt,
    belowcaptionskip = 0pt,
%    tabsize         = 1, % Sem efeito, porque o documento usa Espaço na tabulação.
}

% #1 = linewidth
% #2 = border
% #3 = padding
% #4 = padding-top
% #5 = padding-right
% #6 = padding-bottom
% #7 = padding-left
% #8 = margin-left
% #9 = margin-right
\def\mylstset#1#2#3#4#5#6#7#8#9{\lstset{
    linewidth          = \dimexpr#1 -#2 -#3 -#5\relax,
    framerule          = \dimexpr#2\relax,
    framesep           = \dimexpr#3\relax,
    xleftmargin        = \dimexpr#2 +#3 +#7 +#8\relax,
    xrightmargin       = \dimexpr#9\relax,
    framextopmargin    = \dimexpr#4\relax,
    framexrightmargin  = \dimexpr#5\relax,
    framexbottommargin = \dimexpr#6\relax,
    framexleftmargin   = \dimexpr#7\relax,
}}
\mylstset{\linewidth}{1pt}{0pt}{0pt}{0pt}{0pt}{0pt}{0pt}{0pt}

\lstnewenvironment{lst}[1][\placepos]{%
    \let\medskipamount=\smallskipamount
    \place[#1]\vbox\bgroup
}{%
    \egroup\endplace
}

\def\hookXP{}
\lstnewenvironment{xp}[1][]{%
    \lstset{
        frame           = leftline,
        basicstyle      = \normalsize\ttfamily\singlesp\vskip-\baselineskip,
        texcsstyle      = \normalcolor,
        commentstyle    = \normalcolor,
        backgroundcolor = \color{white},
        gobble          = 4,
    }%
    \mylstset{\linewidth}{0pt}{0pt}{0pt}{0pt}{0pt}{\parindent}{0pt}{0pt}%
    \expandafter\lstset \expandafter{\hookXP,#1}%
    \endgraf\null
}{}

% ETC.
% ================================================

\setmainfont{XITS}[
    UprightFont    = *-Regular,
    BoldFont       = *-Bold,
    ItalicFont     = *-Italic,
    BoldItalicFont = *-BoldItalic,
    Extension      = .otf
]
\setmathfont{XITSMath-Regular.otf}
\setmonofont{InconsolataN}[
    UprightFont = *-Regular,
    BoldFont    = *-Bold,
    Extension   = .otf
]

\definelegendplace{code}{Código}{loc}
\definelegendplace{output}{Saída}{loo}
\definelegendplace{scr}{\emph{Screenshot}}{loscr}

\Crefname{code}{Código}{Códigos}
\Crefname{output}{Saída}{Saídas}
\Crefname{appendix}{Apêndice}{Apêndices}
\Crefname{codesection}{Seção}{Seções}
\Crefname{codesubsection}{Seção}{Seções}
\Crefname{scr}{\emph{Screenshot}}{\emph{Screenshots}}

\let\over=\@@over
\AtBeginDocument{\def\refname{Referências}} % Removido \^ da definição.
\def\hooksection{}
\def\tocsubsectionfont{\itshape}
\def\subsectionfont   {\itshape}
\AtBeginDocument{\pretocmd\appendix{\clearpage}{}{}}
\pretocmd\annex{\clearpage}{}{}
\def\etex{\leavevmode\hbox{$\varepsilon$-\TeX}}
\def\pdflatex{\leavevmode\hbox{PDF\LaTeX}}
\def\xelatex{\leavevmode\hbox{Xe\LaTeX}}
\def\lualatex{\leavevmode\hbox{Lua\LaTeX}}
\def\me{Elaboração própria.}
\def\bibfont{\raggedright\interlinepenalty=10000\singlesp\bibitemsep=\baselineskip}
\appto\textual{
    \def\@evenhead{\abntsmall \firstmarks0\hfil\thepage}
    \def\@oddhead {\abntsmall \botmarks1\hfil\thepage}
    \def\sectionmark   #1{\marks0{Seção \thesection: \unexpanded{#1}}}
    \def\subsectionmark#1{\marks1{Subseção \thesubsection: \unexpanded{#1}}}
    \preto\@oddhead {\lower4pt\rlap{\vrule width\hsize height.4pt}}
    \preto\@evenhead{\lower4pt\rlap{\vrule width\hsize height.4pt}}
}
\patchcmd{\eletroniclayout}{false}{true}{}{}
\let\onesidelayout=\eletroniclayout
\let\twosidelayout=\eletroniclayout
\hypersetup{
    pdfauthor   = Elayson Abreu,
    pdftitle    = Classe ABNTexto: normas ABNT em LaTeX,
    pdfsubject  = Manual da classe ABNTexto,
    pdfcreator  = LaTeX with abntexto,
    pdfkeywords = ABNTexto. LaTeX. Classe. ABNT,
}
\makeatother



\begin{document}
\leavevmode\vskip1.9in
\begingroup \centering \LARGE\sffamily\bfseries
    Classe ABNTexto\vskip-3pt \Large\itshape\mdseries
    Normas ABNT em {\upshape\LaTeX}\vskip2em \normalfont\large
    Elayson Abreu\vskip-3pt
    \href{mailto:abntexto.classe@gmail.com}%
                {abntexto.classe@gmail.com}\vskip2ex
    \today
\vskip2in\endgroup

\noindent\hfil {\Large\ttfamily 1.0}
\newpage

\nonum\notoc\section{Sumário}
\maketoc
\newpage

\section{Introdução}

O `abntexto` é uma classe do \LaTeX\ criada para facilitar a utilização das normas da Associação Brasileira de Normas Técnicas (ABNT) em trabalhos acadêmicos (monografias, dissertações, teses). Ela carrega os recursos necessários para a elaboração de trabalhos nessas normas.

Este \emph{não} é um manual sobre LaTeX (ou TeX, a ferramenta que o constitui) nem tampouco uma descrição detalhada sobre as normas ABNT. Para aprender sobre LaTeX leia \textcite{massago-2022-latex2e-via-exemplos}, em português, ou \textcite{martinsen-2022-latex2e-an-unnoficial}, disponível em inglês, espanhol e francês. Além disso, também se pode ler \textcite{olsak-2012-tex-in-nutshell} e \textcite{eijkhout-1992-tex-by-topic}, disponíveis em inglês, para aprender sobre o sistema \TeX.

Já existe uma classe para aplicação das normas ABNT cujo nome é `abntex2` \cite{araujo-2012-a-classe-abntex2}. Mas, diferente dessa, o `abntexto` não depende do `memoir`, uma classe de código-fonte e manual extensos e que implementa funções as quais, em sua maioria, não são pertinentes a um documento que use as diretrizes da ABNT.

As normas ABNT as quais o `abntexto` se propõe a seguir seguem listadas:

\begin{topics}
    \item NBR 14724:2024. Informação e documentação --- Trabalhos acadêmicos --- Apresentação \cite{abnt-2024-14724-trabalhos-academicos};
%    \item NBR 12225:2004. Informação e documentação --- Lombada --- Apresentação \cite{abnt-2004-12225-lombada};
%    \item NBR 6028:2021. Informação e documentação --- Resumo --- Apresentação \cite{abnt-2021-6028-resumo-resenha-recensao};
    \item NBR 6027:2012. Informação e documentação --- Sumário --- Apresentação \cite{abnt-2012-6027-sumario};
    \item NBR 10520:2023. Informação e documentação --- Citações em documentos --- Apresentação \cite{abnt-2023-10520-citacoes};
%    \item NBR 6023:2018. Informação e documentação --- Referências --- Elaboração \cite{abnt-2018-6023-referencias};
%    \item NBR 5892:2019. Informação e documentação --- Representação e formatos de tempo --- Datas e horas --- Apresentação \cite{abnt-2019-5892-representacao-e-formatos-de-tempo};
    \item NBR 6024:2012. Informação e documentação --- Numeração progressiva das seções de um documento --- Apresentação \cite{abnt-2012-6024-numeracao-progressiva-das};
    \item IBGE. Normas de apresentação tabular \cite{ibge-1993-normas-de-apresentacao}.
\end{topics}

\subsection{Nota de compatibilidade}

No salto da versão anterior, 3.2.1-beta, para esta, 4.0.6-beta, muito do código-fonte foi reescrito; funcionalidades foram adicionadas; comandos obsolesceram. Por consequência, documentos antigos irão certamente quebrar em \emph{layout} e aparência.

Para que os usuários tenham tempo de atualizar para esta nova versão, o arquivo `abntexto-3-2-1-beta.cls` foi deixado no repositório da classe no CTAN\footnote{\url{https://ctan.org/pkg/abntexto}.}. Todavia, para documentos antigos que já foram concluídos, não faz sentido atualizá-los para a versão atual. Nesse caso, seria pertinente copiar o arquivo `abntexto-3-2-1-beta.cls` para o mesmo diretório do documento cuja fase de escrita já foi concluída, chamá-lo por meio de `\documentclass{abntexto-3-2-1-beta}`, e, dessa forma, garantir que o resultado da compilação seja sempre o mesmo.

Para usuários que desejam atualizar seus documentos, o `abntexto` preservou a funcionalidade de comandos obsoletados até onde foi possível e muitas mensagens de alerta/erro foram criadas.

\subsection{Mudanças e novidades}

A sintaxe de `\place{<conteúdo>}` foi substituída por `\begin{place}[<opcional>]`. Além disso, após requisições, a classe ganhou suporte para objetos flutuantes. O posicionamento de `place` pode ser alterado em `<opcional>` com as palavras-chave de objetos flutuantes: `!, t, b, p, h`. Além destas, a classe criou uma nova: `here`\footnote{Semelhante a especificação `H` do pacote `float` \cite{float-package}.}. Com ela, o usuário impõe que `<conteúdo>` seja exibido no PDF na mesma posição em que foi inserido no documento.

Comandos da forma `\<ext>labelwidth` tais como `\toclabelwidth`, `\appendixlabelwidth`, entre outros, não precisam mais ser configurados, porque o `abntexto` alinha os rótulos de suas respectivas entradas automaticamente com a ajuda dos novos comandos `\eqbox` e `\eqboxsize`.

O pacote `enumitem` é carregado pela classe e o ambiente `topics` foi reescrito em função do referido pacote.

A partir de agora, use `\abntsmall` ao invés de `\small`. Por enquanto, o significado deste será igual ao daquele.

\subsection{\emph{Feedback}}

Sugestões, dúvidas, relatos de \emph{bugs} e doações são bem-vindas:\hfil\break {\ttfamily abntexto.classe@gmail.com}.

Chave PIX: {\ttfamily abntexto.classe@gmail.com}

\section{Chamando a classe}

O autor da classe recomenda o uso do \lualatex\ na compilação, embora o \pdflatex\ e o \xelatex\ também sejam suportados. Digite as linhas do \ref{cd:inicio} em conjunto com o \lualatex\ e use `abntexto`.

\legend{code}{Início}
\src\me
\label{cd:inicio}
\begin{lst}
    \documentclass{abntexto}

    \usepackage[english,brazil]{babel}
    \usepackage{fontspec}
    \usepackage{unicode-math}

    \setmainfont{XITS}[
        UprightFont    = *-Regular,
        BoldFont       = *-Bold,
        ItalicFont     = *-Italic,
        BoldItalicFont = *-BoldItalic,
        Extension      = .otf
    ]
    \setmathfont{XITSMath-Regular.otf}

    \begin{document}
        Texto
    \end{document}
\end{lst}

O arquivo `abntexto-exemplo.tex`, disponível no diretório da classe no CTAN, mostra como elaborar os elementos pré-textuais.

\section{Fonte tipográfica}

O tamanho da fonte é 12\,pt de imediato e pode ser configurado redefinindo `\normalsize` por meio de \^`\definesize`, cuja sintaxe é
\begin{xp}
    \definesize\<seletor>{<tamanho de fonte>}{<entrel. simples>}%
                         {<configs. extras>}
\end{xp}
O primeiro argumento é o nome do seletor (`\normalsize`, por exemplo). O segundo argumento é o tamanho da fonte. O terceiro é a entrelinha simples que deve ser selecionada levando em conta o valor do segundo argumento. Por exemplo, uma fonte com tamanho 10\,pt deve ter por volta de 12\,pt de entrelinha simples. O quarto argumento é para ajuste de valores dependentes do tamanho da fonte, por exemplo `\abovedisplayskip` ou `\bigskipamount`.

A classe `article`, na qual se baseia o `abntexto`, fornece os seletores habituais `\tiny`, `\footnotesize`, `\Large` etc. No entanto, esta classe se limitou a redefinir apenas dois: `\normalsize` para o tamanho padrão do texto e \^`\abntsmall` para satisfazer a exigência da \textcite{abnt-2024-14724-trabalhos-academicos} acerca do tamanho menor e uniforme. Até a versão anterior da classe, 3.2.1-beta, o comando `\small` era usado, mas foi substituído por `\abntsmall` para não haver conflito com o significado original de `\small` da classe `article`.

O espaçamento é de 1,5 por padrão e pode ser alterado com os atalhos \^`\singlesp` e \^`\onehalfsp` que alternam para espaçamento simples e de 1,5, respectivamente\footnote{O \type{\\doublesp} está obsoleto.}. A fórmula desses comandos foi atualizada, portanto documentos com versões antigas do `abntexto` terão o \emph{layout} alterado.

Nesta versão, `\spacing` está obsoleto porque esse comando ainda usa a fórmula antiga. Além disso, ele tem uma limitação: só pode ser usado após um seletor de tamanho, porque tais seletores redefinem `\baselineskip` diretamente. Por exemplo, `\abntsmall\spacing{<fator>}` é correto, porém `\spacing{<fator>}\abntsmall` é ineficaz. Essa limitação não existe ao usar os reimplementados `\singlesp` e `\onehalf` de forma que tanto `\singlesp\onehalfsp` quanto `\onehalfsp\singlesp` produzem o mesmo efeito.

\section{\emph{Layout}}

A folha é A4 por padrão e as margens esquerda e direita são de, respectivamente, 3\,cm, 2\,cm para páginas ímpares e de, respectivamente, 2\,cm, 3\,cm para páginas pares. Além disso, as margens superior e inferior têm, nessa ordem, 3\,cm, 2\,cm para todo o documento. Essas margens foram ajustadas com o pacote `geometry`, carregado na classe.

Existem dois comandos que configuram o \emph{layout} da página: \^`\pretextual` e \^`\textual`. O primeiro é iniciado automaticamente durante `\begin{document}`, invoca \^`\onesidelayout` e remove a numeração de página. O segundo, `\textual`, também é executado automaticamente pelo primeiro `\section` numerado, invoca \^`\twosidelayout` e ativa a numeração de página.

O comando `\onesidelayout` configura as páginas para impressão de somente um lado do papel (anverso) ao passo que `\twosidelayout` espelha as margens do documento, habilitando a impressão para os dois lados do papel (anverso e verso).

Os trabalhos acadêmicos podem vir a ser disponiblizados num repositório \emph{online}. Nesse caso, o usuário pode querer reconfigurar as margens de forma a se ter uma leitura mais confortável em dispositivos eletrônicos. Para esse propósito foi criado o comando \^`\eletroniclayout`, apesar dessa configuração não estar prevista na ABNT.

Há duas formas de ativá-lo. A primeira é, no preâmbulo, redefinir `\pretextual` e `\textual` substituindo as instruções `\onesidelayout` e `\twosidelayout` por `\eletroniclayout`. A segunda possibilidade é, também no preâmbulo, redefinir `\onesidelayout` e `\twosidelayout` diretamente conforme \ref{cd:eletroniclayout}. O comando `\eletroniclayout` centraliza o corpo do texto configurando 2,5\,cm nas margens esquerda e direita.

\legend{code}{\emph{Layout} para dispositivos eletrônicos}
\src\me
\label{cd:eletroniclayout}
\begin{lst}
    \documentclass{abntexto}
    \let\onesidelayout=\eletroniclayout
    \let\twosidelayout=\eletroniclayout
    \begin{document}
        texto
    \end{document}
\end{lst}

Outros comandos de estilo de página tais como `\pagestyle`\{{\ttfamily headings, empty, myheadings, plain}\} herdados da classe `article` não fazem sentido no `abntexto`. Ao invés de usá-los, é aconselhável redefinir `\textual`, `\pretextual` ou ambos.

\section{Sumário}

A classe `article` criou o comando `\tableofcontents` para impressão do Sumário, cuja definição será mantida por compatibilidade. No `abntexto`, porém, é recomendado usar \^`\maketoc`, cuja definição é ligeiramente diferente.

Os comandos que configuram a fonte tipográfica das entradas do Sumário são \^`\tocsectionfont`, \^`\tocsubsectionfont` etc. (ver \ref{cd:contents}). O último token nessas definições pode conter parâmetro, por exemplo:
\begin{xp}
    \def\tocsectionfont{\bfseries\MakeUppercase}
    \def\tocparagraphfont{\fbox}
\end{xp}

\legend{code}{Aparência do Sumário}
\src\me
\label{cd:contents}

\begin{lst}
    \documentclass{abntexto}
    \usepackage{xcolor}
    \def \tocsectionfont    {\color{red}\itshape\MakeUppercase}
    \def \tocsubsectionfont {\color{blue}\bfseries}

    \begin{document}
        \nonum\notoc\section{Sumário}
        \maketoc
        \newpage

        \section{Título}
        \subsection{Título}
    \end{document}
\end{lst}

\section{Secionamento}\label{sec:secionamento}

A personalização das macros de seção (\^`\section`, \^`\subsection`, \^`\subsubsection`, \^`\paragraph` e \^`\subparagraph`) é análoga a personalização do Sumário e deve ser feita redefinindo \^`\sectionfont`, \^`\subsectionfont`, \^`\subsubsectionfont`, \^`\paragraphfont` e \^`\subparagraphfont` como exemplificado em \ref{cd:sec} abaixo.

\src\me
\legend{code}{Aparência das seções}
\label{cd:sec}
\begin{lst}
    \documentclass{abntexto}
    \usepackage{xcolor}
    \def \sectionfont       {\color{blue}\scshape}
    \def \subsectionfont    {\color{red}\itshape}
    \def \subsubsectionfont {\color{yellow}\sffamily}

    \begin{document}
        \nonum\notoc\section{Sumário}
        \maketoc

        \section{Texto}
        \subsection{Texto}
        \subsubsection{Texto}
    \end{document}
\end{lst}

Também há como alterar o espaço vertical acima e abaixo das seções redefinindo \^`\abovesection`, \^`\belowsection`. Por exemplo, `\def\abovesubsection{2cm}` insere um espaço de 2\,cm acima de `\subsection`. Se o leitor deseja zerar essas instruções, deve escrever `\def\abovesubsection{0pt}` ao invés de `\def\abovesubsection{}`.

\subsection{O uso de \type{\\chapter} e \type{\\part}}

De acordo com a \textcite{abnt-2024-14724-trabalhos-academicos}: \enquote{[\dots] O trabalho acadêmico não pode ser dividido em capítulos; deve ser organizado em seções [\dots]}.

O usuário pode querer utilizar capítulos no trabalho, apesar de a unidade secional `\section` ser suficiente para um trabalho acadêmico. Para esse propósito, o `abntexto` disponibiliza a instrução \^`\chapter` para a utilização de capítulos no documento, porém não em conformidade com a \textcite{abnt-2024-14724-trabalhos-academicos}, mas no estilo da classe `book`. Para ativá-los é preciso chamar a instrução \^`\usechapters` no preâmbulo.

Também é possível usar \^`\part` em conjunto com o comando \^`\useparts`. Esse útlimo contém `\usechapters` em sua definição, portanto, ao chamar `\useparts`, não se deve chamar `\usechapters` para que este não sobrescreva as definições daquele.

Os níveis secionais `\section` e `\chapter` podem ser prefixados com \^`\nonum` que remove a numeração do título: `\nonum\section`. Os títulos prefixados com `\nonum` serão centralizados \cite{abnt-2024-14724-trabalhos-academicos}. Veja um exemplo de uso em \ref{cd:nonum}.

\legend{code}{Exemplo de título sem numeração e sem ingresso ao Sumário}
\src\me
\label{cd:nonum}

\begin{lst}
    \nonum\notoc\section{Resumo}
    Texto

    \nonum\notoc\section{Sumário}
    \maketoc
\end{lst}

As unidades secionais supracitadas também podem ser prefixadas com \^`\notoc` que instrui o título a não ingressar no Sumário: `\notoc\section`. Além disso, é possível usar ambos os prefixos ao mesmo tempo, conforme mostrado em \ref{cd:nonum}.

O primeiro parágrafo após uma seção é indentado por padrão, dispensando, assim, o uso do pacote `indentfirst`.

Infelizmente, comandos como `\printbibliography`, `\printglossary` e `\printindex` inserem o título internamente. Para que esse título seja exibido corretamente no `abntexto`, é necessário usar o ambiente \^`corrprint`, como demonstrado em \ref{cd:corrprint}.

\legend{code}{Ambiente \type{corrprint}}
\src\me
\label{cd:corrprint}

\begin{lst}
    \begin{corrprint}
        \printbibliography % ou \printglossary, \printindex e semelhantes
    \end{corrprint}
\end{lst}

Existe a possibilidade de o usuário desejar criar um Índice usando o ambiente `theindex` --- embora alternativas automatizadas sejam preferíveis \cite{martinsen-2022-latex2e-an-unnoficial}. Nesse caso, o ambiente `corrprint` funcionará como esperado conforme explicitado em \ref{cd:corrprint-theindex}. O ambiente `corrprint` também pode ser utilizado com o ambiente `thebibliography`, assim como no exemplo com `theindex`.

\legend{code}{Ambientes \type{corrprint} e \type{theindex}}
\src\me
\label{cd:corrprint-theindex}

\begin{lst}
    \begin{corrprint}
        \begin{theindex}
            \item carro, 1
                \subitem árvore, 2
            \indexspace
            \item mar, 3
                \subitem continente, 4
        \end{theindex}
    \end{corrprint}
\end{lst}

\section{Alíneas}

Use o ambiente \^`topics` para criar alíneas. Ele dispõe de dois níveis sendo que o segundo usa um travessão como marcador conforme o exemplo em \ref{cd:topics}. Naturalmente, é permitido usar `\label` após um item.

\legend{code}{Alíneas}
\src\me
\label{cd:topics}

\begin{lst}
    \begin{topics}
        \item \label{al:x} 1
            \begin{topics}
                \item \label{subal:x} x
                \item y
                \item z
            \end{topics}
        \item \label{al:y} 2
        \item 3
    \end{topics}
\end{lst}

\section{Áreas de legenda\label{sec:area-de-legenda}}

Use o ambiente \^`place` para inserir áreas de legenda. Ele substitui integralmente os ambientes `figure` e `table`, bem como o comando `\caption`, não suportados pela classe. Além disso, tem-se \^`\legend`\nobreak`{<tipo>}{<legenda>}` para titular; \^`\src`\nobreak`{<fonte>}` (\emph{source}) para inserir a fonte como mostrado em \ref{cd:place}. Porém, existem condições: `place` deve constar em último lugar e `\label` deve suceder `\legend`.

\legend{code}{Áreas de legenda}
\src\me
\label{cd:place}

\begin{lst}
    \begin{document}
        \legend{figure}{Um título}
        \src{\texttt{www.site.com.br}. Acesso em 01 jan. 2020.}
        \label{fig:teste}
        \begin{place}
            \includegraphics[width=0.4\linewidth]{example-image}
        \end{place}

        Ver \ref{fig:teste}.
    \end{document}
\end{lst}

Além disso, existe a possibilidade de alterar o espaço vertical acima e abaixo de uma área de legenda com `\def`\^`\aboveplace``{<valor>}` e `\def`\^`\belowplace``{<valor>}`. Todavia, tais comandos não têm efeito em objetos flutuantes. Para modificar o espaçamento em `floats` o leitor pode fornecer valores diferentes para `\floatsep`, `\intextsep` e `\textfloatsep`, definidos no LaTeX.

A partir desta versão, o `abntexto` passou a suportar objetos flutuantes. Por exemplo, em
\begin{xp}
    \begin{place}[<especificação>]
\end{xp}
A `<especificação>` é composta pelas palavras-chave habituais: `!`, `t`, `b`, `p`, e `h`. Alternativamente, é possível usar a palavra-chave `here`, definida nesta classe. Com ela, o conteúdo de `place` é mostrado na página exatamente onde foi escrito no documento.

A especificação padrão é `here` caso nenhuma seja explicitamente assumida, mas o usuário pode mudar o padrão digitando, por exemplo, `\def`\nobreak\^`\placepos``{tbp}` no preâmbulo do documento.

\subsection{Criando áreas de legenda}

Até a versão anterior, 3.2.1-beta, o `abntexto` fornecia \^`\definelegendplace` para criação de novos espaços de legenda sob a sintaxe
\begin{xp}
    \definelegendplace{<tipo>}{<nome>}{<ref rótulo>}{<extensão>}
\end{xp}
Nesta versão, no entanto, o terceiro argumento foi excluído em favor do suporte mais extensivo do pacote `cleveref.sty` (não carregado pela classe). Veja o \ref{app:cleveref} para mais detalhes. A partir de agora, a nova sintaxe é
\begin{xp}
    \definelegendplace{<tipo>}{<nome>}{<extensão>}
\end{xp}
Esta classe continuará fornecendo suporte sintático, porém sem funcionalidade, para a versão anterior do referido comando, mas vai exibir uma mensagem de alerta pedindo para que o terceiro argumento seja removido.

Acerca do novo `\definelegendplace`: a instrução
\begin{xp}
    \definelegendplace{diagram}{Diagrama}{lod},
\end{xp}
por exemplo, criará o comando `\makelod` para escrever no documento as entradas do arquivo auxiliar `.lod`. Além disso, o tipo de legenda `diagram` estará disponível para ser usado em `\legend`: `\legend{diagram}{<título>}`.

\subsection{Figuras lado a lado}

O `abntexto` tem suporte nativo para inserção de figuras lado a lado com o ambiente \^`multiplace`. Elas se comportam como caracteres então podemos usar `\hfil` para centralizá-las. Veja o exemplo no \ref{cd:multiplace} e o resultado em \ref{out:result}.

\legend{code}{Figuras lado a lado}
\src\me
\label{cd:multiplace}
\begin{lst}
    \legend{figure}{Duas figuras}
    \src{Elaboração do autor.}

    \begin{multiplace}
        \sublegend{Um título} \label{fig:teste1}
        \subsrc{\texttt{www.site.com.br}. Acesso em 01 jan. 2020.}
        \hfil
        \begin{subplace}[.25\linewidth]
            \includegraphics[width=\linewidth]{example-image-a}
        \end{subplace}
        %
        \sublegend{Outro título} \label{fig:teste2}
        \subsrc{\texttt{www.site.com.br}. Acesso em 01 jan. 2020.}
        \hfil
        \begin{subplace}[.25\linewidth]
            \includegraphics[width=\linewidth]{example-image-b}
        \end{subplace}
    \end{multiplace}

    \parindent=1.5cm Ver \ref{fig:teste1} e \ref{fig:teste2}.
\end{lst}

\legend{output}{Resultado do \ref{cd:multiplace}}
\label{out:result}
\src\me
\begin{place}\vbox\bgroup
    \legend{figure}{Duas figuras}
    \src{Elaboração do autor.}
    %
    \begin{multiplace}
        \sublegend{Um título} \label{fig:teste1}
        \subsrc{\texttt{www.site.com.br}. Acesso em 01 jan. 2020.}
        \hfil
        \begin{subplace}[.25\linewidth]
            \includegraphics[width=\linewidth]{example-image-a}
        \end{subplace}
        %
        \sublegend{Outro título} \label{fig:teste2}
        \subsrc{\texttt{www.site.com.br}. Acesso em 01 jan. 2020.}
        \hfil
        \begin{subplace}[.25\linewidth]
            \includegraphics[width=\linewidth]{example-image-b}
        \end{subplace}
    \end{multiplace}

    \parindent=1.5cm Ver \ref{fig:teste1} e \ref{fig:teste2}.
\egroup\end{place}

Há uma legenda principal impressa com `\legend` e duas sublegendas para cada objeto inseridas com \^`\sublegend`\nobreak`{<título>}`, bem como um par de ambientes \^`subplace`\footnotemark. Além disso, a ordem das instruções é similar àquela aplicada no ambiente `place`: `subplace` deve ser posto após ambos `\sublegend` e \^`\subsrc``{<fonte>}`, e `\label` deve suceder `\legend`.

\footnotetext{A sintaxe \type{\\subplace{<conteúdo>}} está obsoleta. Use `subplace` na forma de ambiente: \type{\\begin{subplace}}.}

Linhas em branco produzem quebra de parágrafo em `multiplace` e, provavelmente, não são desejadas.

Normalmente, espaços em branco depois de `\end{<env>}` são significativos, mas, no caso do ambiente `subplace`, isso não é verdade, porque existe um `\ignorespacesafterend` em sua definição, um comando definido no LaTeX.

Existe, ainda, a possibilidade de um arranjo diferente no `multiplace` usando o parâmetro opcional de `\sublegend`, como se pode ver no \ref{cd:multialt} em conjunto com o resultado em \ref{out:alt}.

\legend{code}{Uso alternativo de \type{multiplace}}
\src{Elaborado pelo autor.}
\label{cd:multialt}
\begin{lst}
    \begin{multiplace}
        \sublegend[table]{Tabela de exemplo} \label{tab1}
        \subsrc{\me}
        \hfil
        \begin{subplace}[.25\linewidth]
            \includegraphics[width=\linewidth]{example-image-a}
        \end{subplace}
        %
        \sublegend[figure]{Título de figura} \label{fig1}
        \subsrc{\me}
        \hfil
        \begin{subplace}[.25\linewidth]
            \includegraphics[width=\linewidth]{example-image-b}
        \end{subplace}
        %
        \sublegend[figure]{Outra figura} \label{fig2}
        \subsrc{\me}
        \vadjust{\kern10pt}\\ \null \hfil
        \begin{subplace}[.25\linewidth]
            \includegraphics[width=\linewidth]{example-image-c}
        \end{subplace}
    \end{multiplace}
\end{lst}

\src{\me}
\begin{place}\vbox\bgroup
\legend{output}{Resultado do \ref{cd:multialt}} \label{out:alt}
\begin{multiplace}
    \sublegend[table]{Tabela de exemplo} \label{tab1}
    \subsrc{\me}
    \hfil
    \begin{subplace}[.25\linewidth]
        \includegraphics[width=\linewidth]{example-image-a}
    \end{subplace}
    \sublegend[figure]{Título de figura} \label{fig1}
    \subsrc{\me}
    \hfil
    \begin{subplace}[.25\linewidth]
        \includegraphics[width=\linewidth]{example-image-b}
    \end{subplace}
    \sublegend[figure]{Outra figura} \label{fig2}
    \subsrc{\me}
    \vadjust{\kern10pt}\\ \null \hfil
    \begin{subplace}[.25\linewidth]
        \includegraphics[width=\linewidth]{example-image-c}
    \end{subplace}
\end{multiplace}
\egroup\end{place}

A \ref{out:alt} possui a \ref{tab1}, \ref{fig1} e \ref{fig2}.

\subsection{Tabelas que se partem entre páginas}

No momento não existe suporte para legendar tabelas e outras estruturas que se partem entre páginas. Se o leitor tiver de usar `longtable.sty` \cite{longtable-package} ou `tabularray.sty` \cite{tabularray-package} (este último é mais flexível), será preciso legendar estruturas tabulares com as ferramentas fornecidas por esses pacotes.

\section{Formatação de \type{\\cite} e Referências}

O `abntexto` não possui macros para formatação de `\cite` e Referências. Utilize `\usepackage[style=abnt]{biblatex}`.

\section{Citações}

Visando seguir a norma \textcite{abnt-2023-10520-citacoes}, são fornecidos `\enquote` do pacote `csquotes`, carregado pela classe, e \^`\Enquote`. Para inserir aspas duplas em citações diretas, de até três linhas, use `\enquote{<texto>}\cite{<chave>}`. Porém, se tratando de citações diretas, com mais de três linhas, de parágrafo único, alterne para `\Enquote{<texto> \cite{<chave>}.}`.

\section{Apêndice e Anexo}

Com os comandos \^`\appendix` e \^`\annex` você pode inserir apêndices e anexos junto dos indicativos desejados: \enquote{APÊNDICE 1 --- \dots} ou \enquote{ANEXO 1 --- \dots} (ver \ref{cd:appendix}).

\legend{code}{Anexos e apêndices}
\src{\me}
\label{cd:appendix}

\begin{lst}
    \begin{document}
        \appendix{Lorem}
        \appendix{Lipsum}

        \annex{Lorem}
        \annex{Lipsum}
    \end{document}
\end{lst}

Até a versão anterior desta classe, era necessário ajustar a largura dos rótulos no Sumário manualmente redefinindo os comandos \^`\appendixlabelwidth` e \^`\annexlabelwidth`. Todavia, nesta versão, 4.0.6-beta, eles se ajustam automaticamente com o auxílio das novas macroinstruções \^`\eqbox` e \^`\eqboxsize`.

\section{Macros para identificação do documento}

\subsection{Ficha catalográfica}

Para produzir fichas catalográficas foi criado o ambiente \^`indexcard`. Ele sempre esteve disponível na forma `\indexcard{<conteúdo>}`, agora obsoleta, mas nunca foi devidamente documentado neste manual. Eis uma amostra de uso:
\begin{xp}
    \documentclass{abntexto}
    \usepackage{lipsum}
    \usepackage{xcolor}
    \begin{document}
        \def\icardWidth{12.5cm}
        \def\icardHeight{7.5cm}
        \def\icardPadding{1cm}
        \def\icardBorder{2pt}
        \appto\icardFont{\color{magenta}}
        \icardPrintGrid

        \noindent\hfil
        \begin{indexcard}
            \lipsum[1][1-4]
        \end{indexcard}
    \end{document}
\end{xp}
Que produz:

\bgroup
\def\icardWidth{12.5cm}
\def\icardHeight{7.5cm}
\def\icardPadding{1cm}
\def\icardBorder{2pt}
\appto\icardFont{\color{magenta}}
\icardPrintGrid
\noindent\hfil
\begin{indexcard}
    \lipsum[1][1-4]
\end{indexcard}
\egroup

Perceba que a instrução \^`\icardPrintGrid` desenha uma grade de quatro linhas para delimitar os contornos do bloco de texto, no intuito de verificar se ele não ultrapassou as margens internas. A definição original de `\icardFont` é
\begin{xp}
    \def\icardFont{\raggedright\parindent=2em\ttfamily}
\end{xp}

\subsection{Utilitário \type{\\Enter}}

O `abntexto` criou o comando \^`\Enter``[<número>]` para produzir um espaço vertical equivalente à tecla `Enter` do teclado. Ele é usado com frequência em `abntexto-exemplo.tex`. O argumento opcional `<número>` é o número de `Enters` que será impresso, por exemplo, `\Enter[2]` imprime um espaço vertical equivalente a dois `Enters` consecutivos. Como `\Enter` é, basicamente, `\vskip\baselineskip`, precedê-lo com `\par` não é necessário, porque `\vskip` possui, a grosso modo, um `\par` em seu funcionamento.

\subsection{Desativando a hifenização}

O comando \^`\nohyph` impede a hifenização dos parágrafos integralmente, sendo útil em títulos e parágrafos curtos. Seu uso pode ser feito assim:
\begin{xp}
    {\nohyph Parágrafo curto ou título\par} % Esse \par é necessário para
                                            % que a formatação do parágrafo
                                            % entre em ação.
\end{xp}

\subsection{Assinatura da banca}

Use \^`\judgeline``{<arg>}` para assinaturas na Folha de Aprovação. Exemplo:
\begin{xp}
    \judgeline{Prof. Dr. Nome Sobrenome \\ Instituição}\Enter
    \judgeline{Prof. Dr. Nome Sobrenome \\ Instituição}\Enter
    \judgeline{Prof. Dr. Nome Sobrenome \\ Instituição}
\end{xp}
Resultado:
\judgeline{Prof. Dr. Nome Sobrenome \\ Instituição}\Enter
\judgeline{Prof. Dr. Nome Sobrenome \\ Instituição}\Enter
\judgeline{Prof. Dr. Nome Sobrenome \\ Instituição}

Veja outro exemplo de `\judgeline` em `abntexto-exemplo.tex`.

\newpage
\begin{corrprint}
    \printbibliography
\end{corrprint}

\newrefsection

\appendix{Cleveref\label{app:cleveref}}

Até a versão passada, 3.2.1-beta, o `abntexto` tinha o poder de inferir o tipo de `\ref` e inserir o substantivo correspondente com auxílio do terceiro parâmetro de `\definelegendplace`. No entanto, esse parâmetro foi removido deixando o referido comando com três argumentos apenas. Como resgatar essa funcionalidade de inferência em `\ref`? Usando o pacote `cleveref.sty` e sua macro `\cref` \cite{cleveref-package}. Veja o exemplo em \ref{cd:cleveref}, cujo resultado é \enquote{Ler Seção 1}.

\legend{code}{Exemplo de uso do \type{cleveref.sty}}
\src\me
\label{cd:cleveref}
\begin{lst}
    \documentclass{abntexto}
    \usepackage[brazil]{babel}
    \usepackage[brazilian,capitalize,noabbrev]{cleveref}
    \begin{document}
        \section{Um título de exemplo\label{sec:ex}}
        Ler \cref{sec:ex}.
    \end{document}
\end{lst}

Percebe-se que a chave `brazil` não é reconhecida pelo `cleveref` compelindo-nos a usar `brazilian`. Além disso, ele deve ser carregado depois do `hyperref.sty`. Veja outras funcionalidades no manual do pacote.

\appendix{Lista de Abreviaturas}

Elemento opcional. O `glossaries-extra.sty` \cite{glossaries-extra-package} consegue criar uma Lista de Abreviaturas, conforme \ref{cd:abbrev}.

\legend{code}{Abreviaturas}
\src\me
\label{cd:abbrev}

\begin{lst}
    \documentclass{abntexto}

    \usepackage[brazil]{babel}
    \usepackage{hyperref}
    \usepackage[automake,shortcuts=all,abbreviations]{glossaries-extra}\makeglossaries
    \usepackage[brazilian,capitalize,noabbrev]{cleveref}

    \newabbreviation[category=abbreviations]{fil}{Fil.}{Filosofia}
    \newabbreviation[category=abbreviations]{el sf}{el.\,sf.}{elemento substantivo feminino}
    \newabbreviation[category=abbreviations]{eng el}{Eng.\,Elétr.}{Engenharia Elétrica}
    \newabbreviation[category=abbreviations]{med legal}{Med.\,Leg.}{Medicina Legal}

    \begin{document}
        Abreviaturas: \ab{fil},
                      \ab{el sf},
                      \ab{eng el},
                      \ab{med legal}.

        \begin{corrprint}\notoc % "Lista de Abreviaturas" não deve constar no Sumário.
            \printglossary[type=abbreviations]
        \end{corrprint}
    \end{document}
\end{lst}

Note que o `cleveref.sty` deve ser carregado depois de `glossaries-extra.sty`, e o `hyperref.sty`, antes.

\appendix{Lista de Siglas}

Elemento opcional. O `glossaries-extra.sty` \cite{glossaries-extra-package} também consegue criar uma Lista de Siglas, conforme \ref{cd:acro}.

\legend{code}{Siglas}
\src\me
\label{cd:acro}

\begin{lst}
    \documentclass{abntexto}

    \usepackage[brazil]{babel}
    \usepackage{hyperref}
    \usepackage[automake,shortcuts=all,acronym]{glossaries-extra}\makeglossaries
    \usepackage[brazilian,capitalize,noabbrev]{cleveref}

    \setabbreviationstyle[acronym]{long-short}

    \newacronym{ibge}{IBGE}{Instituto Brasileiro de Geografia e Estatística}
    \newacronym{html}{HTML}{HyperText Markup Language}
    \newacronym{ce}{CE}{Ceará}
    \newacronym{edo}{EDO}{Equação Diferencial Ordinária}

    \begin{document}
        1ª chamada: \ac{ibge}, \ac{html}, \ac{ce}, \ac{edo}.

        2ª chamada: \ac{ibge}, \ac{html}, \ac{ce}, \ac{edo}.

        \begin{corrprint}\notoc % "Lista de Siglas" não deve constar no Sumário.
            \printglossary[type=acronym]
        \end{corrprint}
    \end{document}
\end{lst}

Note que o `cleveref.sty` deve ser carregado depois de `glossaries-extra.sty`, e o `hyperref.sty`, antes.

\appendix{Lista de Símbolos}

Elemento opcional. O `glossaries-extra.sty` \cite{glossaries-extra-package} também consegue criar uma Lista de Símbolos, conforme \ref{cd:simbolos}.

\legend{code}{Símbolos}
\src\me
\label{cd:simbolos}

\begin{lst}
    \documentclass{abntexto}

    \usepackage[brazil]{babel}
    \usepackage{hyperref}
    \usepackage[automake,shortcuts=all,symbols]{glossaries-extra}\makeglossaries
    \usepackage[brazilian,capitalize,noabbrev]{cleveref}

    \def\newsymbol#1#2#3#4{\newglossaryentry{#1}{
        name={#2},sort={#3},description={#4},type=symbols,category=symbol}}

    \newsymbol{alpha}{\ensuremath{\alpha}}{a}{Letra grega alfa}
    \newsymbol{o(n)}{\ensuremath{O(n)}}{On}{Ordem de um algoritmo}
    \newsymbol{i}{\ensuremath{i}}{i}{Unidade imaginária}
    \newsymbol{A}{\ensuremath{\mathbf{A}}}{A}{Matriz $\mathbf{A}$}
    \newsymbol{N}{N}{N}{Newton}

    \def\symb#1{\gls{#1}}

    \begin{document}
        Símbolos: \symb{alpha}, \symb{o(n)}, \symb{i}, \symb{A}, \symb{N}.

        \begin{corrprint}\notoc % "Lista de Símbolos" não deve constar no Sumário.
            \printglossary[type=symbols]
        \end{corrprint}
    \end{document}
\end{lst}

Note que o `cleveref.sty` deve ser carregado depois de `glossaries-extra.sty`, e o `hyperref.sty`, antes.

\appendix{Glossário}

Elemento opcional. Você pode criar glossários com `glossaries-extra` \cite{glossaries-extra-package}, como em \ref{cd:gloss}.

\legend{code}{Glossário}
\src\me
\label{cd:gloss}

\begin{lst}
    \documentclass{abntexto}

    \usepackage[brazil]{babel}
    \usepackage{hyperref}
    \usepackage[automake,shortcuts=all]{glossaries-extra}\makeglossaries
    \usepackage[brazilian,capitalize,noabbrev]{cleveref}

    \newglossaryentry{metalismo}{
        name        = metalismo,
        description = {Tipo de sistema monetário cujo valor da moeda,
                       geralmente de uma nação ou de um país, é determinado
                       pela quantidade constante de certo metal (ouro ou
                       prata), desta forma se estabelece um valor fixo de
                       troca entre esses metais e o próprio dinheiro.}
    }
    \newglossaryentry{backend}{
        name        = \emph{backend},
        sort        = backend,
        description = {A parte de um site ou aplicação que não é visível
                       para o usuário e que gerencia a lógica de negócios
                       e a comunicação com o banco de dados.}
    }
    \newglossaryentry{cache}{
        name        = cache,
        description = {Um armazenamento temporário de dados para acesso
                       rápido. Ajuda a reduzir o tempo de carregamento.}
    }
    \newglossaryentry{autovetor}{
        name        = autovetor,
        description = {Um vetor não nulo que, quando multiplicado por uma
                       matriz, resulta em um vetor que é um múltiplo escalar do
                       vetor original.}
    }

    \begin{document}
        Glossário: \gls{metalismo}, \gls{backend}, \gls{cache}, \gls{autovetor}.

        \begin{corrprint}
            \printglossary[type=main]
        \end{corrprint}
    \end{document}
\end{lst}

Note que o `cleveref.sty` deve ser carregado depois de `glossaries-extra.sty`, e o `hyperref.sty`, antes.

\appendix{Ambiente \type{place} e o pacote \type{listings}}

O `listings` é um impressor de código-fonte para LaTeX \cite{listings-package}. Você pode usar `\legend` em conjunto com `listings` como mostrado em \ref{cd:listings}. Há algumas observações quanto a essa solução: primeiro, o `listings` não permite que `place` esteja na forma convencional, mas na forma `\place`-`\endplace`. Segundo, a região em \emph{source} terá largura fixa de `\linewidth`, que é a largura padrão de uma `\vbox`.

\legend{code}{\type{\\legend} em conjunto com \type{listings}}
\src\me
\label{cd:listings}

\lstnewenvironment{lstn}{%
    \let\medskipamount=\smallskipamount
    \place\vbox\bgroup
}{%
    \egroup\endplace
}

\begin{lstn}
    % No preâmbulo:
    \usepackage{listings}

    \lstnewenvironment{lst}[1][\placepos]{%
        \let\medskipamount=\smallskipamount
        \place[#1]\vbox\bgroup
    }{%
        \egroup\endplace
    }
    \definelegendplace{code}{Código}{loc}

    % No documento:
    \nonum\notoc\section{Lista de códigos}
    \makeloc

    \legend{code}{Título}
    \src{Fonte}
    \label{Chave}

    \begin{lst}
        int x = 1;
    \end{lst}
\end{lstn}

Na solução acima o ambiente `lst` é inquebrável, portanto quebras de páginas não são possíveis. Se você precisa de um ambiente que se parta entre páginas, veja a solução abaixo (\ref{cd:longlst}):

\legend{code}{O ambiente \type{longlst} quebrável como alternativa ao \type{lst}}
\src\me
\label{cd:longlst}

\begin{lst}
    \lstnewenvironment{longlst}[1][\placepos]{%
        \begingroup
        \processplacearg{#1}%
        \parskip=0pt \parindent=0pt
        \initplace
        \printlegendbox
        \endgraf\nointerlineskip \nobreak\vskip\medskipamount
        \resetORIlabel \ignorespaces
    }{%
        \unskip
        \savedplacewidth=\hsize
        \printsrcbox
        \finishplace
        \resetplace \gresetORIlabel
        \endgroup
    }
\end{lst}

\appendix{Ambiente \type{place} e o pacote \type{algorithm2e}}

O pacote `algorithm2e` foi criado para geração de algoritmos em LaTeX \cite{algorithm2e-package}. É possível usar `place` em conjunto com `algorithm2e` (\ref{cd:algorithm2e}):

\legend{code}{\type{algorithm2e} em \type{place}}
\src\me
\label{cd:algorithm2e}

\begin{lst}
    % No preâmbulo
    \usepackage{algorithm2e}

    \definelegendplace{alg}{Algoritmo}{loa}

    % No documento
    \nonum\notoc\section{Lista de algoritmos}
    \makeloa

    \legend{alg}{How to write algorithms}
    \src{Elaboração própria}
    \label{alg:chave}

    \begin{place}
    \begin{algorithm}[H]
    \SetAlgoLined
    \KwData{this text}
    \KwResult{how to write algorithm with \LaTeX2e }
    initialization\;
    \While{not at end of this document}{
        read current\;
        \eIf{understand}{
            go to next section\;
            current section becomes this one\;
        }{
            go back to the beginning of current section\;
        }
    }
    \end{algorithm}
    \end{place}
\end{lst}

A especificação `H` é obrigatória.

% MACROS PARA O CÓDIGO-FONTE
% ================================================

\pdfstringdefDisableCommands{\def\marks#1#2{}}
\appendix{\marks0{}\marks1{}Código-fonte\label{ap:codigo-fonte}}

\makeatletter
\newcounter{codesection}
\newcounter{codesubsection}[codesection]
\let\c@section=\c@codesection
\let\c@subsection=\c@codesubsection
\def\thecodesection{\arabic{codesection}}
\def\thecodesubsection{\thecodesection.\arabic{codesubsection}}
\edef\toclevel@codesection{\mainseclevel}
\edef\toclevel@codesubsection{\the\numexpr\mainseclevel-1\relax}

\appto\hooktocsection{\def\extleaders{}%
    \edef\savedskips{\rightskip=\the\rightskip \parfillskip=\the\parfillskip\relax}}
\appto\hooktocsubsection{\def\extleaders{}%
    \edef\savedskips{\rightskip=\the\rightskip \parfillskip=\the\parfillskip\relax}}
\appto\hookextline{\savedskips}

\def\toclabelbox{\lowercase{\eqbox{codetoc}}}
\def\tocsectionfont{\itshape\MakeUppercase}
\def\tocsubsectionfont{}
\def\sectionfont{\tocsectionfont}
\def\subsectionfont{\itshape}

% O \addcontentsline{#1}{#2}{#3} está programado para
% enviar marcadores apenas se #1 = toc. Podemos
% mudar isso com \hypersetup{bookmarkstype=<outro>}.
\hypersetup{bookmarkstype=codetoc}
\patchcmd{\targetsection}{\refstepcounter{section}}{\refstepcounter{codesection}}{}{}
\patchcmd{\targetsection}{\addcontentsline{toc}}{\addcontentsline{codetoc}}{}{}
\patchcmd{\targetsubsection}{\refstepcounter{subsection}}{\refstepcounter{codesubsection}}{}{}
\patchcmd{\targetsubsection}{\addcontentsline{toc}}{\addcontentsline{codetoc}}{}{}

\def\hooksection{}
\def\codesection{\section}
\def\codesubsection{\subsection}

\colorlet{C}{.}
\definecolor{tmp}{RGB}{60, 90, 110}  \colorlet{P}{tmp}
                                     \colorlet{B}{blue!7}

\expanded{\noexpand\lstdefinelanguage{mylang}{
    language    = [latex]tex,
    texcs       = {\directlua{for k,v in pairs(tex.primitives())
                   do tex.print(k, v.. ',')end}},
    deletetexcs = {end}
}}

\lstset{
    language    = mylang,
    frame       = single,
    numbers     = left,
    numbersep   = \dimexpr-5pt+\addnumbersep\relax,
    numberstyle = \fontsize{8bp}{9bp}\selectfont,
    literate    = {\^^M\%.}{}1,
}
\def\addnumbersep{9pt}
\mylstset{\linewidth}{0pt}{0pt}{3pt}{0pt}{3pt}{13pt+\addnumbersep}{0pt}{0pt}

\newcount\codecount \codecount=1
\newbox\mybox
\setbox\mybox=\hbox to\hsize{\hfil\abntsmall\ttfamily\bfseries\color{gray}abntexto.cls}
\lstnewenvironment{code}[1][]{%
    \lstset{
        title       = \copy\mybox\nobreak,
        firstnumber = \codecount,
        #1
    }%
}{%
    \global\codecount=\value{lstnumber}\relax
}
\appto\hookXP{
    frame     = none,
    numbers   = none,
    numbersep = 0pt,
    gobble    = 4,
}

\parskip=0pt plus.2ex minus.1ex
\@starttoc{codetoc}
{\def\excl{!}\catcode\1\!=14 \catcode\1\%=9 %! Public Domain Software 2022-2026
%! 2026-03-04
%! 1.0
%! Elayson Abreu
%! abntexto.classe@gmail.com

% !TeX TS-program = lualatex
\documentclass{abntexto}

\usepackage[german,english,brazil,shorthands=off]{babel}
\usepackage[cmyk]{xcolor}
\usepackage{lipsum}
\usepackage[
    style       = abnt,
    maxbibnames = 100,
    extrayear   = true,
]{biblatex} \addbibresource{abntexto.bib}
\usepackage{unicode-math}
\usepackage[colorlinks,linktoc=page]{hyperref}
\usepackage{microtype}
\usepackage[brazilian,capitalize,noabbrev]{cleveref}\AtBeginDocument{\def\ref#1{\cref{#1}}}
\usepackage{fontspec}
\usepackage{listings}
\usepackage{lua-widow-control} \clubpenalty=10000\widowpenalty=10000

\makeatletter

% VERBATIM
% ================================================

\def\adef#1{\catcode`#1=13 \begingroup \lccode`\~=`#1\lowercase{\endgroup\def~}}
\def\setverb{\def\do##1{\catcode`##1=12}\dospecials\space}
\def\verbchar#1{%
    \ifx\savedttchar\undefined\else \catcode\savedttchar=\savedttcharc \fi
    \chardef\savedttchar=`#1\relax
    \chardef\savedttcharc=\catcode`#1\relax
    \adef{#1}{\leavevmode\hbox\bgroup\setverbchar\readverb}%
    \bgroup\lccode`\~=`#1\lowercase{\egroup\def\readverb ##1~}{##1\egroup}%
    \catcode`#1=13
}
\def\setverbchar{\setverb
    \adef{ }{\ }%
    \adef{\{}{\bgroup\normalcolor\char`\{}%
    \adef{\}}{\char`\}\egroup}%
    \adef{[}{\bgroup\normalcolor\char`[}%
    \adef{]}{\char`]\egroup}%
    \ttfamily
}

\edef\bslash{\csstring\\}
\def\mytarget#1{\vbox to0pt{\kern-12pt \hypertarget{#1}{}\vss}}
\def\1{`}\def\2{``}
{\catcode`\`=13 \AtBeginDocument{%
    \verbchar{`}%
    \def\^`{\bgroup \docpoint}%
    \def\docpoint#1`{\egroup\leavevmode\edef\tmp{\csstring#1}%
        \expandafter\parsecn \string#1\relax
        \ifcsname \cn s/^\tmp\endcsname \else
            \mytarget{\cn s/^\tmp}\csgdef{\cn s/^\tmp}{}%
        \fi
        \hyperlink{\cn s/\tmp}{\color{blue}\ttfamily\if\cn c\bslash\fi\tmp}%
    }%
    \def\parsecn #1#2\relax {\edef\cn{\if\bslash#1c\else n\fi}}%
    \def\`{\bgroup \mainpoint}%
    \def\mainpoint #1`{\egroup\leavevmode\edef\tmp{\csstring#1}%
        \expandafter\parsecn \string#1\relax
        \ifcsname \cn s/\tmp\endcsname \moremainpoints \else
            \mytarget{\cn s/\tmp}\csgdef{\cn s/\tmp}{}%
        \fi
        \ifcsname \cn s/^\tmp\endcsname
            \hyperlink{\cn s/^\tmp}{\ttfamily\if\cn c\bslash\fi\tmp}\else
            {\color{red}\ttfamily\string#1}%
        \fi
    }%
    \def\moremainpoints{\errmessage{Second main documentation point \if\cn c\bslash\fi\tmp}}
}}

\let\lchevron=<
{\catcode`<=13 \AtBeginDocument{\catcode`<=13
    \def<#1>{\ifhmode\else\expandafter\hbox\fi\bgroup\normalcolor\ttfamily
             $\color{gray}\langle$\/{\rmfamily\itshape #1\/}$\color{gray}\rangle$\egroup}}%
}

\def\type#1{\NoCaseChange{\protect\typeA{#1}}}
\def\typeA#1{\bgroup\normalshape\ttfamily{\escapechar=-1\relax\expandafter}\detokenize{#1}\egroup}
% Precisamos de \m@ne aqui, \m@ne=-1.
\pdfstringdefDisableCommands{%
    \def\type#1{\bgroup\escapechar\m@ne\expandafter\egroup\detokenize{#1}}%
    \let\escapechar=\relax
    \let\m@ne=\relax
}

\definecolor{P}{RGB}{242, 121, 0}   % Primitivas.
\definecolor{C}{RGB}{23, 230, 23}   % Comentários.
\definecolor{B}{RGB}{245, 250, 247} % Fundo.

\lstset{
    language        = [latex]tex,
    basicstyle      = \abntsmall\singlesp\ttfamily,
    texcsstyle      = *\color{P}, % O listings não permite ":"
    commentstyle    = \color{C},  % no nome dessas cores. Bug.
    backgroundcolor = \color{B},
    frame           = leftline,
    rulecolor       = \color{lightgray},
    columns         = fullflexible,
    breaklines      = false,
    keepspaces      = true,
    aboveskip       = 0pt plus1.2ex minus.67ex,
    belowskip       = 0pt plus1.2ex minus.67ex,
    abovecaptionskip = 0pt,
    belowcaptionskip = 0pt,
%    tabsize         = 1, % Sem efeito, porque o documento usa Espaço na tabulação.
}

% #1 = linewidth
% #2 = border
% #3 = padding
% #4 = padding-top
% #5 = padding-right
% #6 = padding-bottom
% #7 = padding-left
% #8 = margin-left
% #9 = margin-right
\def\mylstset#1#2#3#4#5#6#7#8#9{\lstset{
    linewidth          = \dimexpr#1 -#2 -#3 -#5\relax,
    framerule          = \dimexpr#2\relax,
    framesep           = \dimexpr#3\relax,
    xleftmargin        = \dimexpr#2 +#3 +#7 +#8\relax,
    xrightmargin       = \dimexpr#9\relax,
    framextopmargin    = \dimexpr#4\relax,
    framexrightmargin  = \dimexpr#5\relax,
    framexbottommargin = \dimexpr#6\relax,
    framexleftmargin   = \dimexpr#7\relax,
}}
\mylstset{\linewidth}{1pt}{0pt}{0pt}{0pt}{0pt}{0pt}{0pt}{0pt}

\lstnewenvironment{lst}[1][\placepos]{%
    \let\medskipamount=\smallskipamount
    \place[#1]\vbox\bgroup
}{%
    \egroup\endplace
}

\def\hookXP{}
\lstnewenvironment{xp}[1][]{%
    \lstset{
        frame           = leftline,
        basicstyle      = \normalsize\ttfamily\singlesp\vskip-\baselineskip,
        texcsstyle      = \normalcolor,
        commentstyle    = \normalcolor,
        backgroundcolor = \color{white},
        gobble          = 4,
    }%
    \mylstset{\linewidth}{0pt}{0pt}{0pt}{0pt}{0pt}{\parindent}{0pt}{0pt}%
    \expandafter\lstset \expandafter{\hookXP,#1}%
    \endgraf\null
}{}

% ETC.
% ================================================

\setmainfont{XITS}[
    UprightFont    = *-Regular,
    BoldFont       = *-Bold,
    ItalicFont     = *-Italic,
    BoldItalicFont = *-BoldItalic,
    Extension      = .otf
]
\setmathfont{XITSMath-Regular.otf}
\setmonofont{InconsolataN}[
    UprightFont = *-Regular,
    BoldFont    = *-Bold,
    Extension   = .otf
]

\definelegendplace{code}{Código}{loc}
\definelegendplace{output}{Saída}{loo}
\definelegendplace{scr}{\emph{Screenshot}}{loscr}

\Crefname{code}{Código}{Códigos}
\Crefname{output}{Saída}{Saídas}
\Crefname{appendix}{Apêndice}{Apêndices}
\Crefname{codesection}{Seção}{Seções}
\Crefname{codesubsection}{Seção}{Seções}
\Crefname{scr}{\emph{Screenshot}}{\emph{Screenshots}}

\let\over=\@@over
\AtBeginDocument{\def\refname{Referências}} % Removido \^ da definição.
\def\hooksection{}
\def\tocsubsectionfont{\itshape}
\def\subsectionfont   {\itshape}
\AtBeginDocument{\pretocmd\appendix{\clearpage}{}{}}
\pretocmd\annex{\clearpage}{}{}
\def\etex{\leavevmode\hbox{$\varepsilon$-\TeX}}
\def\pdflatex{\leavevmode\hbox{PDF\LaTeX}}
\def\xelatex{\leavevmode\hbox{Xe\LaTeX}}
\def\lualatex{\leavevmode\hbox{Lua\LaTeX}}
\def\me{Elaboração própria.}
\def\bibfont{\raggedright\interlinepenalty=10000\singlesp\bibitemsep=\baselineskip}
\appto\textual{
    \def\@evenhead{\abntsmall \firstmarks0\hfil\thepage}
    \def\@oddhead {\abntsmall \botmarks1\hfil\thepage}
    \def\sectionmark   #1{\marks0{Seção \thesection: \unexpanded{#1}}}
    \def\subsectionmark#1{\marks1{Subseção \thesubsection: \unexpanded{#1}}}
    \preto\@oddhead {\lower4pt\rlap{\vrule width\hsize height.4pt}}
    \preto\@evenhead{\lower4pt\rlap{\vrule width\hsize height.4pt}}
}
\patchcmd{\eletroniclayout}{false}{true}{}{}
\let\onesidelayout=\eletroniclayout
\let\twosidelayout=\eletroniclayout
\hypersetup{
    pdfauthor   = Elayson Abreu,
    pdftitle    = Classe ABNTexto: normas ABNT em LaTeX,
    pdfsubject  = Manual da classe ABNTexto,
    pdfcreator  = LaTeX with abntexto,
    pdfkeywords = ABNTexto. LaTeX. Classe. ABNT,
}
\makeatother



\begin{document}
\leavevmode\vskip1.9in
\begingroup \centering \LARGE\sffamily\bfseries
    Classe ABNTexto\vskip-3pt \Large\itshape\mdseries
    Normas ABNT em {\upshape\LaTeX}\vskip2em \normalfont\large
    Elayson Abreu\vskip-3pt
    \href{mailto:abntexto.classe@gmail.com}%
                {abntexto.classe@gmail.com}\vskip2ex
    \today
\vskip2in\endgroup

\noindent\hfil {\Large\ttfamily 1.0}
\newpage

\nonum\notoc\section{Sumário}
\maketoc
\newpage

\section{Introdução}

O `abntexto` é uma classe do \LaTeX\ criada para facilitar a utilização das normas da Associação Brasileira de Normas Técnicas (ABNT) em trabalhos acadêmicos (monografias, dissertações, teses). Ela carrega os recursos necessários para a elaboração de trabalhos nessas normas.

Este \emph{não} é um manual sobre LaTeX (ou TeX, a ferramenta que o constitui) nem tampouco uma descrição detalhada sobre as normas ABNT. Para aprender sobre LaTeX leia \textcite{massago-2022-latex2e-via-exemplos}, em português, ou \textcite{martinsen-2022-latex2e-an-unnoficial}, disponível em inglês, espanhol e francês. Além disso, também se pode ler \textcite{olsak-2012-tex-in-nutshell} e \textcite{eijkhout-1992-tex-by-topic}, disponíveis em inglês, para aprender sobre o sistema \TeX.

Já existe uma classe para aplicação das normas ABNT cujo nome é `abntex2` \cite{araujo-2012-a-classe-abntex2}. Mas, diferente dessa, o `abntexto` não depende do `memoir`, uma classe de código-fonte e manual extensos e que implementa funções as quais, em sua maioria, não são pertinentes a um documento que use as diretrizes da ABNT.

As normas ABNT as quais o `abntexto` se propõe a seguir seguem listadas:

\begin{topics}
    \item NBR 14724:2024. Informação e documentação --- Trabalhos acadêmicos --- Apresentação \cite{abnt-2024-14724-trabalhos-academicos};
%    \item NBR 12225:2004. Informação e documentação --- Lombada --- Apresentação \cite{abnt-2004-12225-lombada};
%    \item NBR 6028:2021. Informação e documentação --- Resumo --- Apresentação \cite{abnt-2021-6028-resumo-resenha-recensao};
    \item NBR 6027:2012. Informação e documentação --- Sumário --- Apresentação \cite{abnt-2012-6027-sumario};
    \item NBR 10520:2023. Informação e documentação --- Citações em documentos --- Apresentação \cite{abnt-2023-10520-citacoes};
%    \item NBR 6023:2018. Informação e documentação --- Referências --- Elaboração \cite{abnt-2018-6023-referencias};
%    \item NBR 5892:2019. Informação e documentação --- Representação e formatos de tempo --- Datas e horas --- Apresentação \cite{abnt-2019-5892-representacao-e-formatos-de-tempo};
    \item NBR 6024:2012. Informação e documentação --- Numeração progressiva das seções de um documento --- Apresentação \cite{abnt-2012-6024-numeracao-progressiva-das};
    \item IBGE. Normas de apresentação tabular \cite{ibge-1993-normas-de-apresentacao}.
\end{topics}

\subsection{Nota de compatibilidade}

No salto da versão anterior, 3.2.1-beta, para esta, 4.0.6-beta, muito do código-fonte foi reescrito; funcionalidades foram adicionadas; comandos obsolesceram. Por consequência, documentos antigos irão certamente quebrar em \emph{layout} e aparência.

Para que os usuários tenham tempo de atualizar para esta nova versão, o arquivo `abntexto-3-2-1-beta.cls` foi deixado no repositório da classe no CTAN\footnote{\url{https://ctan.org/pkg/abntexto}.}. Todavia, para documentos antigos que já foram concluídos, não faz sentido atualizá-los para a versão atual. Nesse caso, seria pertinente copiar o arquivo `abntexto-3-2-1-beta.cls` para o mesmo diretório do documento cuja fase de escrita já foi concluída, chamá-lo por meio de `\documentclass{abntexto-3-2-1-beta}`, e, dessa forma, garantir que o resultado da compilação seja sempre o mesmo.

Para usuários que desejam atualizar seus documentos, o `abntexto` preservou a funcionalidade de comandos obsoletados até onde foi possível e muitas mensagens de alerta/erro foram criadas.

\subsection{Mudanças e novidades}

A sintaxe de `\place{<conteúdo>}` foi substituída por `\begin{place}[<opcional>]`. Além disso, após requisições, a classe ganhou suporte para objetos flutuantes. O posicionamento de `place` pode ser alterado em `<opcional>` com as palavras-chave de objetos flutuantes: `!, t, b, p, h`. Além destas, a classe criou uma nova: `here`\footnote{Semelhante a especificação `H` do pacote `float` \cite{float-package}.}. Com ela, o usuário impõe que `<conteúdo>` seja exibido no PDF na mesma posição em que foi inserido no documento.

Comandos da forma `\<ext>labelwidth` tais como `\toclabelwidth`, `\appendixlabelwidth`, entre outros, não precisam mais ser configurados, porque o `abntexto` alinha os rótulos de suas respectivas entradas automaticamente com a ajuda dos novos comandos `\eqbox` e `\eqboxsize`.

O pacote `enumitem` é carregado pela classe e o ambiente `topics` foi reescrito em função do referido pacote.

A partir de agora, use `\abntsmall` ao invés de `\small`. Por enquanto, o significado deste será igual ao daquele.

\subsection{\emph{Feedback}}

Sugestões, dúvidas, relatos de \emph{bugs} e doações são bem-vindas:\hfil\break {\ttfamily abntexto.classe@gmail.com}.

Chave PIX: {\ttfamily abntexto.classe@gmail.com}

\section{Chamando a classe}

O autor da classe recomenda o uso do \lualatex\ na compilação, embora o \pdflatex\ e o \xelatex\ também sejam suportados. Digite as linhas do \ref{cd:inicio} em conjunto com o \lualatex\ e use `abntexto`.

\legend{code}{Início}
\src\me
\label{cd:inicio}
\begin{lst}
    \documentclass{abntexto}

    \usepackage[english,brazil]{babel}
    \usepackage{fontspec}
    \usepackage{unicode-math}

    \setmainfont{XITS}[
        UprightFont    = *-Regular,
        BoldFont       = *-Bold,
        ItalicFont     = *-Italic,
        BoldItalicFont = *-BoldItalic,
        Extension      = .otf
    ]
    \setmathfont{XITSMath-Regular.otf}

    \begin{document}
        Texto
    \end{document}
\end{lst}

O arquivo `abntexto-exemplo.tex`, disponível no diretório da classe no CTAN, mostra como elaborar os elementos pré-textuais.

\section{Fonte tipográfica}

O tamanho da fonte é 12\,pt de imediato e pode ser configurado redefinindo `\normalsize` por meio de \^`\definesize`, cuja sintaxe é
\begin{xp}
    \definesize\<seletor>{<tamanho de fonte>}{<entrel. simples>}%
                         {<configs. extras>}
\end{xp}
O primeiro argumento é o nome do seletor (`\normalsize`, por exemplo). O segundo argumento é o tamanho da fonte. O terceiro é a entrelinha simples que deve ser selecionada levando em conta o valor do segundo argumento. Por exemplo, uma fonte com tamanho 10\,pt deve ter por volta de 12\,pt de entrelinha simples. O quarto argumento é para ajuste de valores dependentes do tamanho da fonte, por exemplo `\abovedisplayskip` ou `\bigskipamount`.

A classe `article`, na qual se baseia o `abntexto`, fornece os seletores habituais `\tiny`, `\footnotesize`, `\Large` etc. No entanto, esta classe se limitou a redefinir apenas dois: `\normalsize` para o tamanho padrão do texto e \^`\abntsmall` para satisfazer a exigência da \textcite{abnt-2024-14724-trabalhos-academicos} acerca do tamanho menor e uniforme. Até a versão anterior da classe, 3.2.1-beta, o comando `\small` era usado, mas foi substituído por `\abntsmall` para não haver conflito com o significado original de `\small` da classe `article`.

O espaçamento é de 1,5 por padrão e pode ser alterado com os atalhos \^`\singlesp` e \^`\onehalfsp` que alternam para espaçamento simples e de 1,5, respectivamente\footnote{O \type{\\doublesp} está obsoleto.}. A fórmula desses comandos foi atualizada, portanto documentos com versões antigas do `abntexto` terão o \emph{layout} alterado.

Nesta versão, `\spacing` está obsoleto porque esse comando ainda usa a fórmula antiga. Além disso, ele tem uma limitação: só pode ser usado após um seletor de tamanho, porque tais seletores redefinem `\baselineskip` diretamente. Por exemplo, `\abntsmall\spacing{<fator>}` é correto, porém `\spacing{<fator>}\abntsmall` é ineficaz. Essa limitação não existe ao usar os reimplementados `\singlesp` e `\onehalf` de forma que tanto `\singlesp\onehalfsp` quanto `\onehalfsp\singlesp` produzem o mesmo efeito.

\section{\emph{Layout}}

A folha é A4 por padrão e as margens esquerda e direita são de, respectivamente, 3\,cm, 2\,cm para páginas ímpares e de, respectivamente, 2\,cm, 3\,cm para páginas pares. Além disso, as margens superior e inferior têm, nessa ordem, 3\,cm, 2\,cm para todo o documento. Essas margens foram ajustadas com o pacote `geometry`, carregado na classe.

Existem dois comandos que configuram o \emph{layout} da página: \^`\pretextual` e \^`\textual`. O primeiro é iniciado automaticamente durante `\begin{document}`, invoca \^`\onesidelayout` e remove a numeração de página. O segundo, `\textual`, também é executado automaticamente pelo primeiro `\section` numerado, invoca \^`\twosidelayout` e ativa a numeração de página.

O comando `\onesidelayout` configura as páginas para impressão de somente um lado do papel (anverso) ao passo que `\twosidelayout` espelha as margens do documento, habilitando a impressão para os dois lados do papel (anverso e verso).

Os trabalhos acadêmicos podem vir a ser disponiblizados num repositório \emph{online}. Nesse caso, o usuário pode querer reconfigurar as margens de forma a se ter uma leitura mais confortável em dispositivos eletrônicos. Para esse propósito foi criado o comando \^`\eletroniclayout`, apesar dessa configuração não estar prevista na ABNT.

Há duas formas de ativá-lo. A primeira é, no preâmbulo, redefinir `\pretextual` e `\textual` substituindo as instruções `\onesidelayout` e `\twosidelayout` por `\eletroniclayout`. A segunda possibilidade é, também no preâmbulo, redefinir `\onesidelayout` e `\twosidelayout` diretamente conforme \ref{cd:eletroniclayout}. O comando `\eletroniclayout` centraliza o corpo do texto configurando 2,5\,cm nas margens esquerda e direita.

\legend{code}{\emph{Layout} para dispositivos eletrônicos}
\src\me
\label{cd:eletroniclayout}
\begin{lst}
    \documentclass{abntexto}
    \let\onesidelayout=\eletroniclayout
    \let\twosidelayout=\eletroniclayout
    \begin{document}
        texto
    \end{document}
\end{lst}

Outros comandos de estilo de página tais como `\pagestyle`\{{\ttfamily headings, empty, myheadings, plain}\} herdados da classe `article` não fazem sentido no `abntexto`. Ao invés de usá-los, é aconselhável redefinir `\textual`, `\pretextual` ou ambos.

\section{Sumário}

A classe `article` criou o comando `\tableofcontents` para impressão do Sumário, cuja definição será mantida por compatibilidade. No `abntexto`, porém, é recomendado usar \^`\maketoc`, cuja definição é ligeiramente diferente.

Os comandos que configuram a fonte tipográfica das entradas do Sumário são \^`\tocsectionfont`, \^`\tocsubsectionfont` etc. (ver \ref{cd:contents}). O último token nessas definições pode conter parâmetro, por exemplo:
\begin{xp}
    \def\tocsectionfont{\bfseries\MakeUppercase}
    \def\tocparagraphfont{\fbox}
\end{xp}

\legend{code}{Aparência do Sumário}
\src\me
\label{cd:contents}

\begin{lst}
    \documentclass{abntexto}
    \usepackage{xcolor}
    \def \tocsectionfont    {\color{red}\itshape\MakeUppercase}
    \def \tocsubsectionfont {\color{blue}\bfseries}

    \begin{document}
        \nonum\notoc\section{Sumário}
        \maketoc
        \newpage

        \section{Título}
        \subsection{Título}
    \end{document}
\end{lst}

\section{Secionamento}\label{sec:secionamento}

A personalização das macros de seção (\^`\section`, \^`\subsection`, \^`\subsubsection`, \^`\paragraph` e \^`\subparagraph`) é análoga a personalização do Sumário e deve ser feita redefinindo \^`\sectionfont`, \^`\subsectionfont`, \^`\subsubsectionfont`, \^`\paragraphfont` e \^`\subparagraphfont` como exemplificado em \ref{cd:sec} abaixo.

\src\me
\legend{code}{Aparência das seções}
\label{cd:sec}
\begin{lst}
    \documentclass{abntexto}
    \usepackage{xcolor}
    \def \sectionfont       {\color{blue}\scshape}
    \def \subsectionfont    {\color{red}\itshape}
    \def \subsubsectionfont {\color{yellow}\sffamily}

    \begin{document}
        \nonum\notoc\section{Sumário}
        \maketoc

        \section{Texto}
        \subsection{Texto}
        \subsubsection{Texto}
    \end{document}
\end{lst}

Também há como alterar o espaço vertical acima e abaixo das seções redefinindo \^`\abovesection`, \^`\belowsection`. Por exemplo, `\def\abovesubsection{2cm}` insere um espaço de 2\,cm acima de `\subsection`. Se o leitor deseja zerar essas instruções, deve escrever `\def\abovesubsection{0pt}` ao invés de `\def\abovesubsection{}`.

\subsection{O uso de \type{\\chapter} e \type{\\part}}

De acordo com a \textcite{abnt-2024-14724-trabalhos-academicos}: \enquote{[\dots] O trabalho acadêmico não pode ser dividido em capítulos; deve ser organizado em seções [\dots]}.

O usuário pode querer utilizar capítulos no trabalho, apesar de a unidade secional `\section` ser suficiente para um trabalho acadêmico. Para esse propósito, o `abntexto` disponibiliza a instrução \^`\chapter` para a utilização de capítulos no documento, porém não em conformidade com a \textcite{abnt-2024-14724-trabalhos-academicos}, mas no estilo da classe `book`. Para ativá-los é preciso chamar a instrução \^`\usechapters` no preâmbulo.

Também é possível usar \^`\part` em conjunto com o comando \^`\useparts`. Esse útlimo contém `\usechapters` em sua definição, portanto, ao chamar `\useparts`, não se deve chamar `\usechapters` para que este não sobrescreva as definições daquele.

Os níveis secionais `\section` e `\chapter` podem ser prefixados com \^`\nonum` que remove a numeração do título: `\nonum\section`. Os títulos prefixados com `\nonum` serão centralizados \cite{abnt-2024-14724-trabalhos-academicos}. Veja um exemplo de uso em \ref{cd:nonum}.

\legend{code}{Exemplo de título sem numeração e sem ingresso ao Sumário}
\src\me
\label{cd:nonum}

\begin{lst}
    \nonum\notoc\section{Resumo}
    Texto

    \nonum\notoc\section{Sumário}
    \maketoc
\end{lst}

As unidades secionais supracitadas também podem ser prefixadas com \^`\notoc` que instrui o título a não ingressar no Sumário: `\notoc\section`. Além disso, é possível usar ambos os prefixos ao mesmo tempo, conforme mostrado em \ref{cd:nonum}.

O primeiro parágrafo após uma seção é indentado por padrão, dispensando, assim, o uso do pacote `indentfirst`.

Infelizmente, comandos como `\printbibliography`, `\printglossary` e `\printindex` inserem o título internamente. Para que esse título seja exibido corretamente no `abntexto`, é necessário usar o ambiente \^`corrprint`, como demonstrado em \ref{cd:corrprint}.

\legend{code}{Ambiente \type{corrprint}}
\src\me
\label{cd:corrprint}

\begin{lst}
    \begin{corrprint}
        \printbibliography % ou \printglossary, \printindex e semelhantes
    \end{corrprint}
\end{lst}

Existe a possibilidade de o usuário desejar criar um Índice usando o ambiente `theindex` --- embora alternativas automatizadas sejam preferíveis \cite{martinsen-2022-latex2e-an-unnoficial}. Nesse caso, o ambiente `corrprint` funcionará como esperado conforme explicitado em \ref{cd:corrprint-theindex}. O ambiente `corrprint` também pode ser utilizado com o ambiente `thebibliography`, assim como no exemplo com `theindex`.

\legend{code}{Ambientes \type{corrprint} e \type{theindex}}
\src\me
\label{cd:corrprint-theindex}

\begin{lst}
    \begin{corrprint}
        \begin{theindex}
            \item carro, 1
                \subitem árvore, 2
            \indexspace
            \item mar, 3
                \subitem continente, 4
        \end{theindex}
    \end{corrprint}
\end{lst}

\section{Alíneas}

Use o ambiente \^`topics` para criar alíneas. Ele dispõe de dois níveis sendo que o segundo usa um travessão como marcador conforme o exemplo em \ref{cd:topics}. Naturalmente, é permitido usar `\label` após um item.

\legend{code}{Alíneas}
\src\me
\label{cd:topics}

\begin{lst}
    \begin{topics}
        \item \label{al:x} 1
            \begin{topics}
                \item \label{subal:x} x
                \item y
                \item z
            \end{topics}
        \item \label{al:y} 2
        \item 3
    \end{topics}
\end{lst}

\section{Áreas de legenda\label{sec:area-de-legenda}}

Use o ambiente \^`place` para inserir áreas de legenda. Ele substitui integralmente os ambientes `figure` e `table`, bem como o comando `\caption`, não suportados pela classe. Além disso, tem-se \^`\legend`\nobreak`{<tipo>}{<legenda>}` para titular; \^`\src`\nobreak`{<fonte>}` (\emph{source}) para inserir a fonte como mostrado em \ref{cd:place}. Porém, existem condições: `place` deve constar em último lugar e `\label` deve suceder `\legend`.

\legend{code}{Áreas de legenda}
\src\me
\label{cd:place}

\begin{lst}
    \begin{document}
        \legend{figure}{Um título}
        \src{\texttt{www.site.com.br}. Acesso em 01 jan. 2020.}
        \label{fig:teste}
        \begin{place}
            \includegraphics[width=0.4\linewidth]{example-image}
        \end{place}

        Ver \ref{fig:teste}.
    \end{document}
\end{lst}

Além disso, existe a possibilidade de alterar o espaço vertical acima e abaixo de uma área de legenda com `\def`\^`\aboveplace``{<valor>}` e `\def`\^`\belowplace``{<valor>}`. Todavia, tais comandos não têm efeito em objetos flutuantes. Para modificar o espaçamento em `floats` o leitor pode fornecer valores diferentes para `\floatsep`, `\intextsep` e `\textfloatsep`, definidos no LaTeX.

A partir desta versão, o `abntexto` passou a suportar objetos flutuantes. Por exemplo, em
\begin{xp}
    \begin{place}[<especificação>]
\end{xp}
A `<especificação>` é composta pelas palavras-chave habituais: `!`, `t`, `b`, `p`, e `h`. Alternativamente, é possível usar a palavra-chave `here`, definida nesta classe. Com ela, o conteúdo de `place` é mostrado na página exatamente onde foi escrito no documento.

A especificação padrão é `here` caso nenhuma seja explicitamente assumida, mas o usuário pode mudar o padrão digitando, por exemplo, `\def`\nobreak\^`\placepos``{tbp}` no preâmbulo do documento.

\subsection{Criando áreas de legenda}

Até a versão anterior, 3.2.1-beta, o `abntexto` fornecia \^`\definelegendplace` para criação de novos espaços de legenda sob a sintaxe
\begin{xp}
    \definelegendplace{<tipo>}{<nome>}{<ref rótulo>}{<extensão>}
\end{xp}
Nesta versão, no entanto, o terceiro argumento foi excluído em favor do suporte mais extensivo do pacote `cleveref.sty` (não carregado pela classe). Veja o \ref{app:cleveref} para mais detalhes. A partir de agora, a nova sintaxe é
\begin{xp}
    \definelegendplace{<tipo>}{<nome>}{<extensão>}
\end{xp}
Esta classe continuará fornecendo suporte sintático, porém sem funcionalidade, para a versão anterior do referido comando, mas vai exibir uma mensagem de alerta pedindo para que o terceiro argumento seja removido.

Acerca do novo `\definelegendplace`: a instrução
\begin{xp}
    \definelegendplace{diagram}{Diagrama}{lod},
\end{xp}
por exemplo, criará o comando `\makelod` para escrever no documento as entradas do arquivo auxiliar `.lod`. Além disso, o tipo de legenda `diagram` estará disponível para ser usado em `\legend`: `\legend{diagram}{<título>}`.

\subsection{Figuras lado a lado}

O `abntexto` tem suporte nativo para inserção de figuras lado a lado com o ambiente \^`multiplace`. Elas se comportam como caracteres então podemos usar `\hfil` para centralizá-las. Veja o exemplo no \ref{cd:multiplace} e o resultado em \ref{out:result}.

\legend{code}{Figuras lado a lado}
\src\me
\label{cd:multiplace}
\begin{lst}
    \legend{figure}{Duas figuras}
    \src{Elaboração do autor.}

    \begin{multiplace}
        \sublegend{Um título} \label{fig:teste1}
        \subsrc{\texttt{www.site.com.br}. Acesso em 01 jan. 2020.}
        \hfil
        \begin{subplace}[.25\linewidth]
            \includegraphics[width=\linewidth]{example-image-a}
        \end{subplace}
        %
        \sublegend{Outro título} \label{fig:teste2}
        \subsrc{\texttt{www.site.com.br}. Acesso em 01 jan. 2020.}
        \hfil
        \begin{subplace}[.25\linewidth]
            \includegraphics[width=\linewidth]{example-image-b}
        \end{subplace}
    \end{multiplace}

    \parindent=1.5cm Ver \ref{fig:teste1} e \ref{fig:teste2}.
\end{lst}

\legend{output}{Resultado do \ref{cd:multiplace}}
\label{out:result}
\src\me
\begin{place}\vbox\bgroup
    \legend{figure}{Duas figuras}
    \src{Elaboração do autor.}
    %
    \begin{multiplace}
        \sublegend{Um título} \label{fig:teste1}
        \subsrc{\texttt{www.site.com.br}. Acesso em 01 jan. 2020.}
        \hfil
        \begin{subplace}[.25\linewidth]
            \includegraphics[width=\linewidth]{example-image-a}
        \end{subplace}
        %
        \sublegend{Outro título} \label{fig:teste2}
        \subsrc{\texttt{www.site.com.br}. Acesso em 01 jan. 2020.}
        \hfil
        \begin{subplace}[.25\linewidth]
            \includegraphics[width=\linewidth]{example-image-b}
        \end{subplace}
    \end{multiplace}

    \parindent=1.5cm Ver \ref{fig:teste1} e \ref{fig:teste2}.
\egroup\end{place}

Há uma legenda principal impressa com `\legend` e duas sublegendas para cada objeto inseridas com \^`\sublegend`\nobreak`{<título>}`, bem como um par de ambientes \^`subplace`\footnotemark. Além disso, a ordem das instruções é similar àquela aplicada no ambiente `place`: `subplace` deve ser posto após ambos `\sublegend` e \^`\subsrc``{<fonte>}`, e `\label` deve suceder `\legend`.

\footnotetext{A sintaxe \type{\\subplace{<conteúdo>}} está obsoleta. Use `subplace` na forma de ambiente: \type{\\begin{subplace}}.}

Linhas em branco produzem quebra de parágrafo em `multiplace` e, provavelmente, não são desejadas.

Normalmente, espaços em branco depois de `\end{<env>}` são significativos, mas, no caso do ambiente `subplace`, isso não é verdade, porque existe um `\ignorespacesafterend` em sua definição, um comando definido no LaTeX.

Existe, ainda, a possibilidade de um arranjo diferente no `multiplace` usando o parâmetro opcional de `\sublegend`, como se pode ver no \ref{cd:multialt} em conjunto com o resultado em \ref{out:alt}.

\legend{code}{Uso alternativo de \type{multiplace}}
\src{Elaborado pelo autor.}
\label{cd:multialt}
\begin{lst}
    \begin{multiplace}
        \sublegend[table]{Tabela de exemplo} \label{tab1}
        \subsrc{\me}
        \hfil
        \begin{subplace}[.25\linewidth]
            \includegraphics[width=\linewidth]{example-image-a}
        \end{subplace}
        %
        \sublegend[figure]{Título de figura} \label{fig1}
        \subsrc{\me}
        \hfil
        \begin{subplace}[.25\linewidth]
            \includegraphics[width=\linewidth]{example-image-b}
        \end{subplace}
        %
        \sublegend[figure]{Outra figura} \label{fig2}
        \subsrc{\me}
        \vadjust{\kern10pt}\\ \null \hfil
        \begin{subplace}[.25\linewidth]
            \includegraphics[width=\linewidth]{example-image-c}
        \end{subplace}
    \end{multiplace}
\end{lst}

\src{\me}
\begin{place}\vbox\bgroup
\legend{output}{Resultado do \ref{cd:multialt}} \label{out:alt}
\begin{multiplace}
    \sublegend[table]{Tabela de exemplo} \label{tab1}
    \subsrc{\me}
    \hfil
    \begin{subplace}[.25\linewidth]
        \includegraphics[width=\linewidth]{example-image-a}
    \end{subplace}
    \sublegend[figure]{Título de figura} \label{fig1}
    \subsrc{\me}
    \hfil
    \begin{subplace}[.25\linewidth]
        \includegraphics[width=\linewidth]{example-image-b}
    \end{subplace}
    \sublegend[figure]{Outra figura} \label{fig2}
    \subsrc{\me}
    \vadjust{\kern10pt}\\ \null \hfil
    \begin{subplace}[.25\linewidth]
        \includegraphics[width=\linewidth]{example-image-c}
    \end{subplace}
\end{multiplace}
\egroup\end{place}

A \ref{out:alt} possui a \ref{tab1}, \ref{fig1} e \ref{fig2}.

\subsection{Tabelas que se partem entre páginas}

No momento não existe suporte para legendar tabelas e outras estruturas que se partem entre páginas. Se o leitor tiver de usar `longtable.sty` \cite{longtable-package} ou `tabularray.sty` \cite{tabularray-package} (este último é mais flexível), será preciso legendar estruturas tabulares com as ferramentas fornecidas por esses pacotes.

\section{Formatação de \type{\\cite} e Referências}

O `abntexto` não possui macros para formatação de `\cite` e Referências. Utilize `\usepackage[style=abnt]{biblatex}`.

\section{Citações}

Visando seguir a norma \textcite{abnt-2023-10520-citacoes}, são fornecidos `\enquote` do pacote `csquotes`, carregado pela classe, e \^`\Enquote`. Para inserir aspas duplas em citações diretas, de até três linhas, use `\enquote{<texto>}\cite{<chave>}`. Porém, se tratando de citações diretas, com mais de três linhas, de parágrafo único, alterne para `\Enquote{<texto> \cite{<chave>}.}`.

\section{Apêndice e Anexo}

Com os comandos \^`\appendix` e \^`\annex` você pode inserir apêndices e anexos junto dos indicativos desejados: \enquote{APÊNDICE 1 --- \dots} ou \enquote{ANEXO 1 --- \dots} (ver \ref{cd:appendix}).

\legend{code}{Anexos e apêndices}
\src{\me}
\label{cd:appendix}

\begin{lst}
    \begin{document}
        \appendix{Lorem}
        \appendix{Lipsum}

        \annex{Lorem}
        \annex{Lipsum}
    \end{document}
\end{lst}

Até a versão anterior desta classe, era necessário ajustar a largura dos rótulos no Sumário manualmente redefinindo os comandos \^`\appendixlabelwidth` e \^`\annexlabelwidth`. Todavia, nesta versão, 4.0.6-beta, eles se ajustam automaticamente com o auxílio das novas macroinstruções \^`\eqbox` e \^`\eqboxsize`.

\section{Macros para identificação do documento}

\subsection{Ficha catalográfica}

Para produzir fichas catalográficas foi criado o ambiente \^`indexcard`. Ele sempre esteve disponível na forma `\indexcard{<conteúdo>}`, agora obsoleta, mas nunca foi devidamente documentado neste manual. Eis uma amostra de uso:
\begin{xp}
    \documentclass{abntexto}
    \usepackage{lipsum}
    \usepackage{xcolor}
    \begin{document}
        \def\icardWidth{12.5cm}
        \def\icardHeight{7.5cm}
        \def\icardPadding{1cm}
        \def\icardBorder{2pt}
        \appto\icardFont{\color{magenta}}
        \icardPrintGrid

        \noindent\hfil
        \begin{indexcard}
            \lipsum[1][1-4]
        \end{indexcard}
    \end{document}
\end{xp}
Que produz:

\bgroup
\def\icardWidth{12.5cm}
\def\icardHeight{7.5cm}
\def\icardPadding{1cm}
\def\icardBorder{2pt}
\appto\icardFont{\color{magenta}}
\icardPrintGrid
\noindent\hfil
\begin{indexcard}
    \lipsum[1][1-4]
\end{indexcard}
\egroup

Perceba que a instrução \^`\icardPrintGrid` desenha uma grade de quatro linhas para delimitar os contornos do bloco de texto, no intuito de verificar se ele não ultrapassou as margens internas. A definição original de `\icardFont` é
\begin{xp}
    \def\icardFont{\raggedright\parindent=2em\ttfamily}
\end{xp}

\subsection{Utilitário \type{\\Enter}}

O `abntexto` criou o comando \^`\Enter``[<número>]` para produzir um espaço vertical equivalente à tecla `Enter` do teclado. Ele é usado com frequência em `abntexto-exemplo.tex`. O argumento opcional `<número>` é o número de `Enters` que será impresso, por exemplo, `\Enter[2]` imprime um espaço vertical equivalente a dois `Enters` consecutivos. Como `\Enter` é, basicamente, `\vskip\baselineskip`, precedê-lo com `\par` não é necessário, porque `\vskip` possui, a grosso modo, um `\par` em seu funcionamento.

\subsection{Desativando a hifenização}

O comando \^`\nohyph` impede a hifenização dos parágrafos integralmente, sendo útil em títulos e parágrafos curtos. Seu uso pode ser feito assim:
\begin{xp}
    {\nohyph Parágrafo curto ou título\par} % Esse \par é necessário para
                                            % que a formatação do parágrafo
                                            % entre em ação.
\end{xp}

\subsection{Assinatura da banca}

Use \^`\judgeline``{<arg>}` para assinaturas na Folha de Aprovação. Exemplo:
\begin{xp}
    \judgeline{Prof. Dr. Nome Sobrenome \\ Instituição}\Enter
    \judgeline{Prof. Dr. Nome Sobrenome \\ Instituição}\Enter
    \judgeline{Prof. Dr. Nome Sobrenome \\ Instituição}
\end{xp}
Resultado:
\judgeline{Prof. Dr. Nome Sobrenome \\ Instituição}\Enter
\judgeline{Prof. Dr. Nome Sobrenome \\ Instituição}\Enter
\judgeline{Prof. Dr. Nome Sobrenome \\ Instituição}

Veja outro exemplo de `\judgeline` em `abntexto-exemplo.tex`.

\newpage
\begin{corrprint}
    \printbibliography
\end{corrprint}

\newrefsection

\appendix{Cleveref\label{app:cleveref}}

Até a versão passada, 3.2.1-beta, o `abntexto` tinha o poder de inferir o tipo de `\ref` e inserir o substantivo correspondente com auxílio do terceiro parâmetro de `\definelegendplace`. No entanto, esse parâmetro foi removido deixando o referido comando com três argumentos apenas. Como resgatar essa funcionalidade de inferência em `\ref`? Usando o pacote `cleveref.sty` e sua macro `\cref` \cite{cleveref-package}. Veja o exemplo em \ref{cd:cleveref}, cujo resultado é \enquote{Ler Seção 1}.

\legend{code}{Exemplo de uso do \type{cleveref.sty}}
\src\me
\label{cd:cleveref}
\begin{lst}
    \documentclass{abntexto}
    \usepackage[brazil]{babel}
    \usepackage[brazilian,capitalize,noabbrev]{cleveref}
    \begin{document}
        \section{Um título de exemplo\label{sec:ex}}
        Ler \cref{sec:ex}.
    \end{document}
\end{lst}

Percebe-se que a chave `brazil` não é reconhecida pelo `cleveref` compelindo-nos a usar `brazilian`. Além disso, ele deve ser carregado depois do `hyperref.sty`. Veja outras funcionalidades no manual do pacote.

\appendix{Lista de Abreviaturas}

Elemento opcional. O `glossaries-extra.sty` \cite{glossaries-extra-package} consegue criar uma Lista de Abreviaturas, conforme \ref{cd:abbrev}.

\legend{code}{Abreviaturas}
\src\me
\label{cd:abbrev}

\begin{lst}
    \documentclass{abntexto}

    \usepackage[brazil]{babel}
    \usepackage{hyperref}
    \usepackage[automake,shortcuts=all,abbreviations]{glossaries-extra}\makeglossaries
    \usepackage[brazilian,capitalize,noabbrev]{cleveref}

    \newabbreviation[category=abbreviations]{fil}{Fil.}{Filosofia}
    \newabbreviation[category=abbreviations]{el sf}{el.\,sf.}{elemento substantivo feminino}
    \newabbreviation[category=abbreviations]{eng el}{Eng.\,Elétr.}{Engenharia Elétrica}
    \newabbreviation[category=abbreviations]{med legal}{Med.\,Leg.}{Medicina Legal}

    \begin{document}
        Abreviaturas: \ab{fil},
                      \ab{el sf},
                      \ab{eng el},
                      \ab{med legal}.

        \begin{corrprint}\notoc % "Lista de Abreviaturas" não deve constar no Sumário.
            \printglossary[type=abbreviations]
        \end{corrprint}
    \end{document}
\end{lst}

Note que o `cleveref.sty` deve ser carregado depois de `glossaries-extra.sty`, e o `hyperref.sty`, antes.

\appendix{Lista de Siglas}

Elemento opcional. O `glossaries-extra.sty` \cite{glossaries-extra-package} também consegue criar uma Lista de Siglas, conforme \ref{cd:acro}.

\legend{code}{Siglas}
\src\me
\label{cd:acro}

\begin{lst}
    \documentclass{abntexto}

    \usepackage[brazil]{babel}
    \usepackage{hyperref}
    \usepackage[automake,shortcuts=all,acronym]{glossaries-extra}\makeglossaries
    \usepackage[brazilian,capitalize,noabbrev]{cleveref}

    \setabbreviationstyle[acronym]{long-short}

    \newacronym{ibge}{IBGE}{Instituto Brasileiro de Geografia e Estatística}
    \newacronym{html}{HTML}{HyperText Markup Language}
    \newacronym{ce}{CE}{Ceará}
    \newacronym{edo}{EDO}{Equação Diferencial Ordinária}

    \begin{document}
        1ª chamada: \ac{ibge}, \ac{html}, \ac{ce}, \ac{edo}.

        2ª chamada: \ac{ibge}, \ac{html}, \ac{ce}, \ac{edo}.

        \begin{corrprint}\notoc % "Lista de Siglas" não deve constar no Sumário.
            \printglossary[type=acronym]
        \end{corrprint}
    \end{document}
\end{lst}

Note que o `cleveref.sty` deve ser carregado depois de `glossaries-extra.sty`, e o `hyperref.sty`, antes.

\appendix{Lista de Símbolos}

Elemento opcional. O `glossaries-extra.sty` \cite{glossaries-extra-package} também consegue criar uma Lista de Símbolos, conforme \ref{cd:simbolos}.

\legend{code}{Símbolos}
\src\me
\label{cd:simbolos}

\begin{lst}
    \documentclass{abntexto}

    \usepackage[brazil]{babel}
    \usepackage{hyperref}
    \usepackage[automake,shortcuts=all,symbols]{glossaries-extra}\makeglossaries
    \usepackage[brazilian,capitalize,noabbrev]{cleveref}

    \def\newsymbol#1#2#3#4{\newglossaryentry{#1}{
        name={#2},sort={#3},description={#4},type=symbols,category=symbol}}

    \newsymbol{alpha}{\ensuremath{\alpha}}{a}{Letra grega alfa}
    \newsymbol{o(n)}{\ensuremath{O(n)}}{On}{Ordem de um algoritmo}
    \newsymbol{i}{\ensuremath{i}}{i}{Unidade imaginária}
    \newsymbol{A}{\ensuremath{\mathbf{A}}}{A}{Matriz $\mathbf{A}$}
    \newsymbol{N}{N}{N}{Newton}

    \def\symb#1{\gls{#1}}

    \begin{document}
        Símbolos: \symb{alpha}, \symb{o(n)}, \symb{i}, \symb{A}, \symb{N}.

        \begin{corrprint}\notoc % "Lista de Símbolos" não deve constar no Sumário.
            \printglossary[type=symbols]
        \end{corrprint}
    \end{document}
\end{lst}

Note que o `cleveref.sty` deve ser carregado depois de `glossaries-extra.sty`, e o `hyperref.sty`, antes.

\appendix{Glossário}

Elemento opcional. Você pode criar glossários com `glossaries-extra` \cite{glossaries-extra-package}, como em \ref{cd:gloss}.

\legend{code}{Glossário}
\src\me
\label{cd:gloss}

\begin{lst}
    \documentclass{abntexto}

    \usepackage[brazil]{babel}
    \usepackage{hyperref}
    \usepackage[automake,shortcuts=all]{glossaries-extra}\makeglossaries
    \usepackage[brazilian,capitalize,noabbrev]{cleveref}

    \newglossaryentry{metalismo}{
        name        = metalismo,
        description = {Tipo de sistema monetário cujo valor da moeda,
                       geralmente de uma nação ou de um país, é determinado
                       pela quantidade constante de certo metal (ouro ou
                       prata), desta forma se estabelece um valor fixo de
                       troca entre esses metais e o próprio dinheiro.}
    }
    \newglossaryentry{backend}{
        name        = \emph{backend},
        sort        = backend,
        description = {A parte de um site ou aplicação que não é visível
                       para o usuário e que gerencia a lógica de negócios
                       e a comunicação com o banco de dados.}
    }
    \newglossaryentry{cache}{
        name        = cache,
        description = {Um armazenamento temporário de dados para acesso
                       rápido. Ajuda a reduzir o tempo de carregamento.}
    }
    \newglossaryentry{autovetor}{
        name        = autovetor,
        description = {Um vetor não nulo que, quando multiplicado por uma
                       matriz, resulta em um vetor que é um múltiplo escalar do
                       vetor original.}
    }

    \begin{document}
        Glossário: \gls{metalismo}, \gls{backend}, \gls{cache}, \gls{autovetor}.

        \begin{corrprint}
            \printglossary[type=main]
        \end{corrprint}
    \end{document}
\end{lst}

Note que o `cleveref.sty` deve ser carregado depois de `glossaries-extra.sty`, e o `hyperref.sty`, antes.

\appendix{Ambiente \type{place} e o pacote \type{listings}}

O `listings` é um impressor de código-fonte para LaTeX \cite{listings-package}. Você pode usar `\legend` em conjunto com `listings` como mostrado em \ref{cd:listings}. Há algumas observações quanto a essa solução: primeiro, o `listings` não permite que `place` esteja na forma convencional, mas na forma `\place`-`\endplace`. Segundo, a região em \emph{source} terá largura fixa de `\linewidth`, que é a largura padrão de uma `\vbox`.

\legend{code}{\type{\\legend} em conjunto com \type{listings}}
\src\me
\label{cd:listings}

\lstnewenvironment{lstn}{%
    \let\medskipamount=\smallskipamount
    \place\vbox\bgroup
}{%
    \egroup\endplace
}

\begin{lstn}
    % No preâmbulo:
    \usepackage{listings}

    \lstnewenvironment{lst}[1][\placepos]{%
        \let\medskipamount=\smallskipamount
        \place[#1]\vbox\bgroup
    }{%
        \egroup\endplace
    }
    \definelegendplace{code}{Código}{loc}

    % No documento:
    \nonum\notoc\section{Lista de códigos}
    \makeloc

    \legend{code}{Título}
    \src{Fonte}
    \label{Chave}

    \begin{lst}
        int x = 1;
    \end{lst}
\end{lstn}

Na solução acima o ambiente `lst` é inquebrável, portanto quebras de páginas não são possíveis. Se você precisa de um ambiente que se parta entre páginas, veja a solução abaixo (\ref{cd:longlst}):

\legend{code}{O ambiente \type{longlst} quebrável como alternativa ao \type{lst}}
\src\me
\label{cd:longlst}

\begin{lst}
    \lstnewenvironment{longlst}[1][\placepos]{%
        \begingroup
        \processplacearg{#1}%
        \parskip=0pt \parindent=0pt
        \initplace
        \printlegendbox
        \endgraf\nointerlineskip \nobreak\vskip\medskipamount
        \resetORIlabel \ignorespaces
    }{%
        \unskip
        \savedplacewidth=\hsize
        \printsrcbox
        \finishplace
        \resetplace \gresetORIlabel
        \endgroup
    }
\end{lst}

\appendix{Ambiente \type{place} e o pacote \type{algorithm2e}}

O pacote `algorithm2e` foi criado para geração de algoritmos em LaTeX \cite{algorithm2e-package}. É possível usar `place` em conjunto com `algorithm2e` (\ref{cd:algorithm2e}):

\legend{code}{\type{algorithm2e} em \type{place}}
\src\me
\label{cd:algorithm2e}

\begin{lst}
    % No preâmbulo
    \usepackage{algorithm2e}

    \definelegendplace{alg}{Algoritmo}{loa}

    % No documento
    \nonum\notoc\section{Lista de algoritmos}
    \makeloa

    \legend{alg}{How to write algorithms}
    \src{Elaboração própria}
    \label{alg:chave}

    \begin{place}
    \begin{algorithm}[H]
    \SetAlgoLined
    \KwData{this text}
    \KwResult{how to write algorithm with \LaTeX2e }
    initialization\;
    \While{not at end of this document}{
        read current\;
        \eIf{understand}{
            go to next section\;
            current section becomes this one\;
        }{
            go back to the beginning of current section\;
        }
    }
    \end{algorithm}
    \end{place}
\end{lst}

A especificação `H` é obrigatória.

% MACROS PARA O CÓDIGO-FONTE
% ================================================

\pdfstringdefDisableCommands{\def\marks#1#2{}}
\appendix{\marks0{}\marks1{}Código-fonte\label{ap:codigo-fonte}}

\makeatletter
\newcounter{codesection}
\newcounter{codesubsection}[codesection]
\let\c@section=\c@codesection
\let\c@subsection=\c@codesubsection
\def\thecodesection{\arabic{codesection}}
\def\thecodesubsection{\thecodesection.\arabic{codesubsection}}
\edef\toclevel@codesection{\mainseclevel}
\edef\toclevel@codesubsection{\the\numexpr\mainseclevel-1\relax}

\appto\hooktocsection{\def\extleaders{}%
    \edef\savedskips{\rightskip=\the\rightskip \parfillskip=\the\parfillskip\relax}}
\appto\hooktocsubsection{\def\extleaders{}%
    \edef\savedskips{\rightskip=\the\rightskip \parfillskip=\the\parfillskip\relax}}
\appto\hookextline{\savedskips}

\def\toclabelbox{\lowercase{\eqbox{codetoc}}}
\def\tocsectionfont{\itshape\MakeUppercase}
\def\tocsubsectionfont{}
\def\sectionfont{\tocsectionfont}
\def\subsectionfont{\itshape}

% O \addcontentsline{#1}{#2}{#3} está programado para
% enviar marcadores apenas se #1 = toc. Podemos
% mudar isso com \hypersetup{bookmarkstype=<outro>}.
\hypersetup{bookmarkstype=codetoc}
\patchcmd{\targetsection}{\refstepcounter{section}}{\refstepcounter{codesection}}{}{}
\patchcmd{\targetsection}{\addcontentsline{toc}}{\addcontentsline{codetoc}}{}{}
\patchcmd{\targetsubsection}{\refstepcounter{subsection}}{\refstepcounter{codesubsection}}{}{}
\patchcmd{\targetsubsection}{\addcontentsline{toc}}{\addcontentsline{codetoc}}{}{}

\def\hooksection{}
\def\codesection{\section}
\def\codesubsection{\subsection}

\colorlet{C}{.}
\definecolor{tmp}{RGB}{60, 90, 110}  \colorlet{P}{tmp}
                                     \colorlet{B}{blue!7}

\expanded{\noexpand\lstdefinelanguage{mylang}{
    language    = [latex]tex,
    texcs       = {\directlua{for k,v in pairs(tex.primitives())
                   do tex.print(k, v.. ',')end}},
    deletetexcs = {end}
}}

\lstset{
    language    = mylang,
    frame       = single,
    numbers     = left,
    numbersep   = \dimexpr-5pt+\addnumbersep\relax,
    numberstyle = \fontsize{8bp}{9bp}\selectfont,
    literate    = {\^^M\%.}{}1,
}
\def\addnumbersep{9pt}
\mylstset{\linewidth}{0pt}{0pt}{3pt}{0pt}{3pt}{13pt+\addnumbersep}{0pt}{0pt}

\newcount\codecount \codecount=1
\newbox\mybox
\setbox\mybox=\hbox to\hsize{\hfil\abntsmall\ttfamily\bfseries\color{gray}abntexto.cls}
\lstnewenvironment{code}[1][]{%
    \lstset{
        title       = \copy\mybox\nobreak,
        firstnumber = \codecount,
        #1
    }%
}{%
    \global\codecount=\value{lstnumber}\relax
}
\appto\hookXP{
    frame     = none,
    numbers   = none,
    numbersep = 0pt,
    gobble    = 4,
}

\parskip=0pt plus.2ex minus.1ex
\@starttoc{codetoc}
{\def\excl{!}\catcode\1\!=14 \catcode\1\%=9 \input{abntexto.cls}}

% Não funciona:
% https://tex.stackexchange.com/questions/35279/biblatex-multiple-bibliographies-categorised-by-different-bib-files
% Complicado:
% https://tex.stackexchange.com/questions/299064/biblatex-two-bibliographies-with-different-styles-and-sortings?noredirect=1&lq=1
% Bom ponto de partida:
% https://tex.stackexchange.com/questions/6967/how-to-split-bibliography-into-works-cited-and-works-not-cited
% Bom ponto de partida: https://tex.stackexchange.com/questions/596128/separate-multiple-bib-files-reference-using-biblatex
% Solução definitiva:
% https://tex.stackexchange.com/questions/429970/is-there-a-way-i-can-manually-reset-or-clear-the-citations-in-biblatex
% Usaremos \newrefsection.
\newpage
\begin{corrprint}\def\refname{REFERÊNCIAS}
    \printbibliography
\end{corrprint}

\end{document}}

% Não funciona:
% https://tex.stackexchange.com/questions/35279/biblatex-multiple-bibliographies-categorised-by-different-bib-files
% Complicado:
% https://tex.stackexchange.com/questions/299064/biblatex-two-bibliographies-with-different-styles-and-sortings?noredirect=1&lq=1
% Bom ponto de partida:
% https://tex.stackexchange.com/questions/6967/how-to-split-bibliography-into-works-cited-and-works-not-cited
% Bom ponto de partida: https://tex.stackexchange.com/questions/596128/separate-multiple-bib-files-reference-using-biblatex
% Solução definitiva:
% https://tex.stackexchange.com/questions/429970/is-there-a-way-i-can-manually-reset-or-clear-the-citations-in-biblatex
% Usaremos \newrefsection.
\newpage
\begin{corrprint}\def\refname{REFERÊNCIAS}
    \printbibliography
\end{corrprint}

\end{document}}

% Não funciona:
% https://tex.stackexchange.com/questions/35279/biblatex-multiple-bibliographies-categorised-by-different-bib-files
% Complicado:
% https://tex.stackexchange.com/questions/299064/biblatex-two-bibliographies-with-different-styles-and-sortings?noredirect=1&lq=1
% Bom ponto de partida:
% https://tex.stackexchange.com/questions/6967/how-to-split-bibliography-into-works-cited-and-works-not-cited
% Bom ponto de partida: https://tex.stackexchange.com/questions/596128/separate-multiple-bib-files-reference-using-biblatex
% Solução definitiva:
% https://tex.stackexchange.com/questions/429970/is-there-a-way-i-can-manually-reset-or-clear-the-citations-in-biblatex
% Usaremos \newrefsection.
\newpage
\begin{corrprint}\def\refname{REFERÊNCIAS}
    \printbibliography
\end{corrprint}

\end{document}}

% Não funciona:
% https://tex.stackexchange.com/questions/35279/biblatex-multiple-bibliographies-categorised-by-different-bib-files
% Complicado:
% https://tex.stackexchange.com/questions/299064/biblatex-two-bibliographies-with-different-styles-and-sortings?noredirect=1&lq=1
% Bom ponto de partida:
% https://tex.stackexchange.com/questions/6967/how-to-split-bibliography-into-works-cited-and-works-not-cited
% Bom ponto de partida: https://tex.stackexchange.com/questions/596128/separate-multiple-bib-files-reference-using-biblatex
% Solução definitiva:
% https://tex.stackexchange.com/questions/429970/is-there-a-way-i-can-manually-reset-or-clear-the-citations-in-biblatex
% Usaremos \newrefsection.
\newpage
\begin{corrprint}\def\refname{REFERÊNCIAS}
    \printbibliography
\end{corrprint}

\end{document}